% Options for packages loaded elsewhere
\PassOptionsToPackage{unicode}{hyperref}
\PassOptionsToPackage{hyphens}{url}
\documentclass[
  ignorenonframetext,
  aspectratio=169,
]{beamer}
\newif\ifbibliography
\usepackage{pgfpages}
\setbeamertemplate{caption}[numbered]
\setbeamertemplate{caption label separator}{: }
\setbeamercolor{caption name}{fg=normal text.fg}
\beamertemplatenavigationsymbolsempty
% remove section numbering
\setbeamertemplate{part page}{
  \centering
  \begin{beamercolorbox}[sep=16pt,center]{part title}
    \usebeamerfont{part title}\insertpart\par
  \end{beamercolorbox}
}
\setbeamertemplate{section page}{
  \centering
  \begin{beamercolorbox}[sep=12pt,center]{section title}
    \usebeamerfont{section title}\insertsection\par
  \end{beamercolorbox}
}
\setbeamertemplate{subsection page}{
  \centering
  \begin{beamercolorbox}[sep=8pt,center]{subsection title}
    \usebeamerfont{subsection title}\insertsubsection\par
  \end{beamercolorbox}
}
% Prevent slide breaks in the middle of a paragraph
\widowpenalties 1 10000
\raggedbottom
\AtBeginPart{
  \frame{\partpage}
}
\AtBeginSection{
  \ifbibliography
  \else
    \frame{\sectionpage}
  \fi
}
\AtBeginSubsection{
  \frame{\subsectionpage}
}
\usepackage{iftex}
\ifPDFTeX
  \usepackage[T1]{fontenc}
  \usepackage[utf8]{inputenc}
  \usepackage{textcomp} % provide euro and other symbols
\else % if luatex or xetex
  \usepackage{unicode-math} % this also loads fontspec
  \defaultfontfeatures{Scale=MatchLowercase}
  \defaultfontfeatures[\rmfamily]{Ligatures=TeX,Scale=1}
\fi
\usepackage{lmodern}
\ifPDFTeX\else
  % xetex/luatex font selection
\fi
% Use upquote if available, for straight quotes in verbatim environments
\IfFileExists{upquote.sty}{\usepackage{upquote}}{}
\IfFileExists{microtype.sty}{% use microtype if available
  \usepackage[]{microtype}
  \UseMicrotypeSet[protrusion]{basicmath} % disable protrusion for tt fonts
}{}
\makeatletter
\@ifundefined{KOMAClassName}{% if non-KOMA class
  \IfFileExists{parskip.sty}{%
    \usepackage{parskip}
  }{% else
    \setlength{\parindent}{0pt}
    \setlength{\parskip}{6pt plus 2pt minus 1pt}}
}{% if KOMA class
  \KOMAoptions{parskip=half}}
\makeatother
\usepackage{color}
\usepackage{fancyvrb}
\newcommand{\VerbBar}{|}
\newcommand{\VERB}{\Verb[commandchars=\\\{\}]}
\DefineVerbatimEnvironment{Highlighting}{Verbatim}{commandchars=\\\{\}}
% Add ',fontsize=\small' for more characters per line
\newenvironment{Shaded}{}{}
\newcommand{\AlertTok}[1]{\textcolor[rgb]{1.00,0.00,0.00}{\textbf{#1}}}
\newcommand{\AnnotationTok}[1]{\textcolor[rgb]{0.38,0.63,0.69}{\textbf{\textit{#1}}}}
\newcommand{\AttributeTok}[1]{\textcolor[rgb]{0.49,0.56,0.16}{#1}}
\newcommand{\BaseNTok}[1]{\textcolor[rgb]{0.25,0.63,0.44}{#1}}
\newcommand{\BuiltInTok}[1]{\textcolor[rgb]{0.00,0.50,0.00}{#1}}
\newcommand{\CharTok}[1]{\textcolor[rgb]{0.25,0.44,0.63}{#1}}
\newcommand{\CommentTok}[1]{\textcolor[rgb]{0.38,0.63,0.69}{\textit{#1}}}
\newcommand{\CommentVarTok}[1]{\textcolor[rgb]{0.38,0.63,0.69}{\textbf{\textit{#1}}}}
\newcommand{\ConstantTok}[1]{\textcolor[rgb]{0.53,0.00,0.00}{#1}}
\newcommand{\ControlFlowTok}[1]{\textcolor[rgb]{0.00,0.44,0.13}{\textbf{#1}}}
\newcommand{\DataTypeTok}[1]{\textcolor[rgb]{0.56,0.13,0.00}{#1}}
\newcommand{\DecValTok}[1]{\textcolor[rgb]{0.25,0.63,0.44}{#1}}
\newcommand{\DocumentationTok}[1]{\textcolor[rgb]{0.73,0.13,0.13}{\textit{#1}}}
\newcommand{\ErrorTok}[1]{\textcolor[rgb]{1.00,0.00,0.00}{\textbf{#1}}}
\newcommand{\ExtensionTok}[1]{#1}
\newcommand{\FloatTok}[1]{\textcolor[rgb]{0.25,0.63,0.44}{#1}}
\newcommand{\FunctionTok}[1]{\textcolor[rgb]{0.02,0.16,0.49}{#1}}
\newcommand{\ImportTok}[1]{\textcolor[rgb]{0.00,0.50,0.00}{\textbf{#1}}}
\newcommand{\InformationTok}[1]{\textcolor[rgb]{0.38,0.63,0.69}{\textbf{\textit{#1}}}}
\newcommand{\KeywordTok}[1]{\textcolor[rgb]{0.00,0.44,0.13}{\textbf{#1}}}
\newcommand{\NormalTok}[1]{#1}
\newcommand{\OperatorTok}[1]{\textcolor[rgb]{0.40,0.40,0.40}{#1}}
\newcommand{\OtherTok}[1]{\textcolor[rgb]{0.00,0.44,0.13}{#1}}
\newcommand{\PreprocessorTok}[1]{\textcolor[rgb]{0.74,0.48,0.00}{#1}}
\newcommand{\RegionMarkerTok}[1]{#1}
\newcommand{\SpecialCharTok}[1]{\textcolor[rgb]{0.25,0.44,0.63}{#1}}
\newcommand{\SpecialStringTok}[1]{\textcolor[rgb]{0.73,0.40,0.53}{#1}}
\newcommand{\StringTok}[1]{\textcolor[rgb]{0.25,0.44,0.63}{#1}}
\newcommand{\VariableTok}[1]{\textcolor[rgb]{0.10,0.09,0.49}{#1}}
\newcommand{\VerbatimStringTok}[1]{\textcolor[rgb]{0.25,0.44,0.63}{#1}}
\newcommand{\WarningTok}[1]{\textcolor[rgb]{0.38,0.63,0.69}{\textbf{\textit{#1}}}}
\setlength{\emergencystretch}{3em} % prevent overfull lines
\providecommand{\tightlist}{%
  \setlength{\itemsep}{0pt}\setlength{\parskip}{0pt}}
\usetheme{Madrid}
\usecolortheme{dolphin}
\setbeamertemplate{navigation symbols}{}
\setbeamertemplate{footline}[frame number]
\setbeamersize{text margin left=5mm,text margin right=5mm}
\setbeamerfont{caption}{size=\scriptsize}
\setbeamerfont{footnote}{size=\tiny}
\usepackage{bookmark}
\IfFileExists{xurl.sty}{\usepackage{xurl}}{} % add URL line breaks if available
\urlstyle{same}
\hypersetup{
  pdftitle={What Four Research Attempts Taught Us},
  hidelinks,
  pdfcreator={LaTeX via pandoc}}

\title{\texorpdfstring{What Four Research Attempts Taught
Us}{What Four Research Attempts Taught Us}}
\subtitle{\texorpdfstring{Plain-language main talk with technical
appendix}{Plain-language main talk with technical appendix}}
\author{\texorpdfstring{}{}}
\date{2026-07-18}

\begin{document}
\frame{\titlepage}

\begin{frame}[allowframebreaks]
  \setcounter{tocdepth}{1}
  \tableofcontents
\end{frame}
\setcounter{tocdepth}{1}
\tableofcontents
}
\section{What We Learned From Four Research
Attempts}\label{what-we-learned-from-four-research-attempts}

\begin{frame}[fragile,allowframebreaks]{What We Learned From Four
Research Attempts}
Date: 2026-07-13

Presentation status: \textbf{CURRENT DRAFT -\/- not yet presented at a
meeting.}

Latest revision: 2026-07-18. This main section explains the decisions in
plain language. Exact commits, protocols, and evidence paths are
preserved in the technical appendix and
\texttt{docs/saq\_research\_direction\_registry.md}.
\end{frame}

\begin{frame}[allowframebreaks]{1. Today: What Do We Need To Decide?}
\protect\phantomsection\label{1-today-what-do-we-need-to-decide}
This is not a presentation of a finished method.

\begin{itemize}
\tightlist
\item
  Should any of the four ideas continue as the next paper direction?
\item
  Which evidence is strong enough to keep?
\item
  Which attractive interpretations are not supported?
\item
  What must a new direction prove before we invest again?
\end{itemize}
\end{frame}

\begin{frame}[allowframebreaks]{2. One-Minute Background}
\protect\phantomsection\label{2-one-minute-background}
SAQ compresses each database vector so search can compare many
candidates quickly. Compression saves work, but can distort which
candidates look close.

\begin{itemize}
\tightlist
\item
  \textbf{PCA means principal component analysis}: it rotates data to
  expose major directions of variation.
\item
  \textbf{Recall means the fraction of true nearest neighbors found}.
\item
  \textbf{QPS means queries per second}: higher is faster.
\item
  Our bar is not ``one metric improved.'' We need better quality at
  comparable speed, memory, and construction cost.
\end{itemize}
\end{frame}

\begin{frame}[allowframebreaks]{3. The Four-Attempt Map}
\protect\phantomsection\label{3-the-four-attempt-map}
\begin{itemize}
\tightlist
\item
  \textbf{Attempt}: 1

  \textbf{Plain-language question}: Can a better rotation or smaller
  representation fix SAQ\textquotesingle s error?
\item
  \textbf{Decision}: Stop
\item
  \textbf{Attempt}: 2

  \textbf{Plain-language question}: Can exact scalar training improve
  both compression and search?
\item
  \textbf{Decision}: Keep as a baseline; stop as a method
\item
  \textbf{Attempt}: 3

  \textbf{Plain-language question}: Did the chosen quality metric make
  an earlier result look worse?
\item
  \textbf{Decision}: Keep as measurement evidence; stop as a method
\item
  \textbf{Attempt}: 4

  \textbf{Plain-language question}: Can non-power-of-two choices use a
  fixed bit budget better?
\item
  \textbf{Decision}: Portfolio closed
\end{itemize}

The purpose is research selection, not presenting every experiment as a
win.
\end{frame}

\begin{frame}[allowframebreaks]{4. Attempt 1: The Intuition}
\protect\phantomsection\label{4-attempt-1-the-intuition}
Imagine viewing the same cloud of points from another angle. A better
angle might put the important variation into the parts SAQ represents
most accurately.

Attempt 1 tested two versions:

\begin{itemize}
\tightlist
\item
  rotate all dimensions differently; or
\item
  keep fewer dimensions and summarize what was removed.
\end{itemize}

A full-dimensional rotation preserves exact geometric distances.
Removing dimensions does not, so the second version faced a stricter
correctness test.
\end{frame}

\begin{frame}[allowframebreaks]{5. Attempt 1: Strongest Evidence And
Stop}
\protect\phantomsection\label{5-attempt-1-strongest-evidence-and-stop}
\begin{itemize}
\tightlist
\item
  A promising signal on one benchmark dataset (GIST) did not repeat on a
  second dataset (CIFAR).
\item
  The deliberately favorable reduced-dimension test still ranked worse
  than native full-dimensional SAQ before we built a search index.
\item
  Therefore there is no general ``better PCA'' or ``safe smaller SAQ''
  result.
\end{itemize}

What remains: a useful warning that one-dataset estimator gains may not
replicate. What stops: both proposed mechanisms.
\end{frame}

\begin{frame}[allowframebreaks]{6. Attempt 2: The Intuition}
\protect\phantomsection\label{6-attempt-2-the-intuition}
A scalar codebook replaces numbers with a small set of representatives.
Ordinary training can settle at a locally good answer; exact training
finds the best answer for the frozen one-dimensional histogram.

The question was whether that cleaner offline fit also improves actual
nearest-neighbor search, after counting its training cost.
\end{frame}

\begin{frame}[allowframebreaks]{7. Attempt 2: Strongest Evidence And
Stop}
\protect\phantomsection\label{7-attempt-2-strongest-evidence-and-stop}
\begin{itemize}
\tightlist
\item
  Exact training reduced reconstruction error---the difference between
  original and compressed values---in some regimes.
\item
  Search quality did not improve consistently; the direction of the
  effect changed across settings.
\item
  The inner exact one-dimensional optimizer is established prior work,
  so it is not our novelty.
\end{itemize}

What remains: a stronger offline baseline and evidence that lower
compression error does not guarantee better search. What stops:
presenting exact scalar training as a new database method.
\end{frame}

\begin{frame}[allowframebreaks]{8. Attempt 3: The Intuition}
\protect\phantomsection\label{8-attempt-3-the-intuition}
Recall asks whether the returned item identifiers exactly match the true
top items. A distance-ratio score instead asks how much farther the
returned items are than the true ones.

Two methods can have different Recall even when their returned distances
are similar. Attempt 3 checked whether this changed an earlier decision.
\end{frame}

\begin{frame}[allowframebreaks]{9. Attempt 3: Strongest Evidence And
Stop}
\protect\phantomsection\label{9-attempt-3-strongest-evidence-and-stop}
\begin{itemize}
\tightlist
\item
  At one GIST operating point, the distance-based interpretation
  differed from the exact-identifier Recall interpretation.
\item
  Two control settings on a second dataset remained negative.
\item
  No new encoding or search mechanism was introduced.
\end{itemize}

What remains: evidence that the conclusion can depend on the chosen
quality metric. What stops: claiming a replicated method advantage from
one changed interpretation.
\end{frame}

\begin{frame}[fragile,allowframebreaks]{10. Attempt 4: The (3,5) Versus
(4,4) Example}
\protect\phantomsection\label{10-attempt-4-the-35-versus-44-example}
Suppose two factors must share one fixed word with at most 16 joint
states.

\begin{itemize}
\tightlist
\item
  Power-of-two choices naturally allow \texttt{(4,4)}:
  \texttt{4\ ×\ 4\ =\ 16}.
\item
  Allowing arbitrary positive integers also allows \texttt{(3,5)}:
  \texttt{3\ ×\ 5\ =\ 15}.
\item
  The extra choice could fit uneven data better without increasing
  stored bits.
\end{itemize}

This proves only that the feasible choice set is larger. It does not
prove better search quality or lower system cost.
\end{frame}

\begin{frame}[fragile,allowframebreaks]{11. Attempt 4: Strongest
Evidence And Stop}
\protect\phantomsection\label{11-attempt-4-strongest-evidence-and-stop}
\begin{itemize}
\tightlist
\item
  The exact frozen construction pipeline was correct on its registered
  synthetic checks.
\item
  The original registered cost study projected about \texttt{24.17}
  CPU-hours, above its preregistered 24-hour ceiling, so it stopped.
\item
  The later recovery attempt produced no scientific or performance
  result. Its last run started but ended without a valid completion
  record, so its runtime outcome remains unknown.
\item
  The project decided to stop recovering the unresolved START-only
  experimental pipeline (\texttt{STOP\_NO\_FURTHER\_ARTIFACT\_RECOVERY}
  in the registry).
\end{itemize}

What remains: a mathematical possibility and a cost-stop record. What
stops: further recovery work and every method, search, or performance
claim.
\end{frame}

\begin{frame}[allowframebreaks]{12. Common Lesson Across The Four
Attempts}
\protect\phantomsection\label{12-common-lesson-across-the-four-attempts}
A better internal objective is not automatically a better search system.

The recurring failure modes were:

\begin{itemize}
\tightlist
\item
  a signal did not replicate;
\item
  discarded information mattered;
\item
  offline fit did not predict retrieval;
\item
  the metric changed interpretation but not mechanism; or
\item
  exact construction and evidence cost too much.
\end{itemize}
\end{frame}

\begin{frame}[allowframebreaks]{13. What We Should Keep}
\protect\phantomsection\label{13-what-we-should-keep}
\begin{itemize}
\tightlist
\item
  SAQ as a strong baseline.
\item
  Exact scalar training as a tougher offline comparison.
\item
  Both Recall and distance-quality views when they answer different
  questions.
\item
  Negative results with frozen thresholds and fair controls.
\item
  Early cost checks before building full experiment machinery.
\end{itemize}

These assets improve the next study even though none is a new method.
\end{frame}

\begin{frame}[fragile,allowframebreaks]{14. What We Cannot Claim}
\protect\phantomsection\label{14-what-we-cannot-claim}
Unsafe interpretation:

\begin{Shaded}
\begin{Highlighting}[]
\NormalTok{PCA is universally optimal, exact training never helps search, or arbitrary}
\NormalTok{cardinalities cannot help.}
\end{Highlighting}
\end{Shaded}

Safe interpretation: the tested mechanisms did not produce a general,
overhead-controlled improvement. The later Attempt 4 work has no valid
final execution record, scientific decision, performance result, or
production integration result.
\end{frame}

\begin{frame}[allowframebreaks]{15. Entry Gate For The Next Project}
\protect\phantomsection\label{15-entry-gate-for-the-next-project}
Before opening another implementation branch, require:

\begin{enumerate}
\tightlist
\item
  a clear problem that holds on at least two independent strong
  baselines;
\item
  the most relevant published work, and why our idea is more than a
  direct combination;
\item
  one small, cheap test that can quickly show the idea is not worth
  pursuing, with inputs and the stop threshold fixed in advance;
\item
  construction, memory, and query overhead counted from the start; and
\item
  a fair path to moving quality versus speed, not only an internal
  objective.
\end{enumerate}

No next Attempt 4 research or execution step is authorized.
\end{frame}

\begin{frame}[allowframebreaks]{16. Two Questions For Discussion}
\protect\phantomsection\label{16-two-questions-for-discussion}
\begin{enumerate}
\tightlist
\item
  Do we agree to close these four method lines while retaining their
  baselines and negative evidence?
\item
  For the next review, should we first select a new problem from primary
  work, with no implementation until a small, low-cost stop test is
  approved?
\end{enumerate}
\end{frame}

\section{Technical Appendix}\label{technical-appendix}

\begin{frame}[allowframebreaks]{Technical Appendix}
Normal presentation ends here. Use the following appendix only for
questions about exact definitions, protocols, commits, or experiment
details.
\end{frame}

\section{Post-SAQ Pivot: Attempts 1, 2, 3, And
4}\label{post-saq-pivot-attempts-1-2-3-and-4}

\begin{frame}[fragile,allowframebreaks]{Post-SAQ Pivot: Attempts 1, 2,
3, And 4}
PCA Replacement, Lossy Projection, Exact Scalar-Codebook DP, Recent ASQ
Prior-Art Audit, Distance-Quality Re-evaluation, And An
Arbitrary-Cardinality Fixed-Rate Quantization Gate

Date: 2026-07-13

Presentation status: \textbf{CURRENT DRAFT -\/- not yet presented at a
meeting.}

Latest revision: 2026-07-18. Incorporates the 2026-07-13 audit of
\href{https://arxiv.org/abs/2606.00289}{White and Singal,
arXiv:2606.00289v1} and its pinned official code, the terminal A4-1S
snapshot \texttt{saq-arbitrary-cardinality-analysis@f1b464b}, and the
corrected A4 V2 source plus terminal PAR snapshot
\texttt{saq-arbitrary-cardinality-feasibility-v2@30dfada}, independently
reviewed at \texttt{fd5367e}, and the

later exact executable-identity source repair \texttt{@e48df452},
independently reviewed at \texttt{c401dae}, the cache/staging
disposition target \texttt{@56210f8} and review record
\texttt{@249d5b8}, and the generic cache-policy source/static target
\texttt{@c33a2bf}, independently reviewed at \texttt{@5db3025}, and the
exact

PREP authorization target \texttt{@e7f940e}, independently reviewed at
\texttt{@16a8201}, followed by the host-rebind erratum target
\texttt{@5a47fed}, independently reviewed at \texttt{@b89dabe}, and the
exact HOST-I-AUTH target \texttt{@212a67b}, independently reviewed at
\texttt{@71e6bec}, followed by the

source-authority erratum target \texttt{@6fe8544}, independently
reviewed at \texttt{@9fa9528}, and the subsequent HOST-I source-static
target \texttt{@4602585}, whose direct-child independent review at
\texttt{@e8e9c79} recorded one HIGH and a terminal source-static target
review failure, followed by the bounded HOST-I-R1

correction-only repair protocol target \texttt{@ddfef99}, independently
reviewed at \texttt{@b1a7429} with
\texttt{HOST\_I\_R1\_CORRECTION\_ONLY\_REPAIR\_PROTOCOL\_REVIEW\_PASS},
and the exact correction target \texttt{@e17f887}, independently
reviewed at \texttt{@e10bde7} with
\texttt{HOST\_I\_R1\_CORRECTION\_ONLY\_REPAIR\_REVIEW\_PASS}, followed
by the bounded source-history epoch-correction protocol target
\texttt{@dcaed57}, independently reviewed

at \texttt{@e98a3e4} with
\texttt{SOURCE\_HISTORY\_EPOCH\_CORRECTION\_PROTOCOL\_REVIEW\_PASS}. The
first authorized source-history implementation then stopped before
target because its frozen ten-path closure was unsatisfiable.
Correction-only path-closure erratum target \texttt{@150e5b3},
independently reviewed at \texttt{@7749491}, reached
\texttt{SOURCE\_HISTORY\_I\_R1\_PATH\_CLOSURE\_ERRATUM\_REVIEW\_PASS}
and freezes only an eight-path/five-derived-object

correction boundary. The subsequently authorized exact correction target
\texttt{@e9b5c83}, independently reviewed at \texttt{@a1198c4}, reached
\texttt{SOURCE\_HISTORY\_EPOCH\_CORRECTION\_SOURCE\_STATIC\_REVIEW\_PASS}
with zero LOW-or-higher findings. Fresh dormant PREP authorization
target \texttt{@cd1757e}, independently reviewed at \texttt{@eb1b893},
then reached \texttt{PREP\_AUTHORIZATION\_EXACT\_TARGET\_REVIEW\_PASS}
with zero

LOW-or-higher findings. Its subsequently authorized unique probe
returned \texttt{eb1b893}, and the unique invocation was consumed but
left only durable START plus empty captures before the PREP body. Frozen
status is therefore
\texttt{CRASH\_OR\_UNKNOWN\ /\ NO\_PAR\_AUTHORITY\ /\ NO\_SCIENTIFIC\_DECISION}.
Additive crash-record protocol target

\texttt{@1982865}, independently reviewed at \texttt{@8bec546}, reached
\texttt{START\_ONLY\_CRASH\_RECORD\_PROTOCOL\_REVIEW\_PASS}; the actual
terminal record, retry, and source repair remain unauthorized. Audited
closure \texttt{@3577edd} now records portfolio state
\texttt{CLOSED\ /\ STOP\_NO\_FURTHER\_ARTIFACT\_RECOVERY} without
changing START-only runtime status or creating an artifact, scientific,

or performance result.

Audience assumption: familiar with vector search and vector quantization
at a high level, but not with SAQ\textquotesingle s transform,
segmentation, or the experiments in these repositories.

Deck status: Attempts 1, 2, and 3 are complete studies. Attempt 4
preserves the terminal A4-1S \texttt{NO\_GO\_EXACT\_SOLVER\_COST}, while
a new primary-source review supports a narrower A4 V2 protocol for a
differently defined synthetic comparative-instrument cost question. Its

additive PAR-report timing erratum is
\texttt{PROTOCOL\_ERRATUM\_INDEPENDENT\_REVIEW\_PASS}. Source-only
A4-V2-I reached \texttt{SOURCE\_IMPLEMENTED\_STATIC\_REVIEW\_PASS}, but
the later authorized PAR command stopped before build at frozen status
\texttt{ARTIFACT\_INVALID}. This is an artifact- validation failure with
no valid PAR authority and

no scientific decision. A4-V2-I-R1 subsequently repaired only the
CPython identity source and reached
\texttt{SOURCE\_REPAIR\_STATIC\_REVIEW\_PASS}; it did not run build or
parity. The later CACHE-P exact target reached only
\texttt{EXACT\_TARGET\_INDEPENDENT\_REVIEW\_PASS}: its protocol retains
the old quarantine, requires

staged sanitized remote isolation, and rejects direct PAR-R1. CACHE-I
then reached only
\texttt{GENERIC\_CACHE\_POLICY\_SOURCE\_STATIC\_REVIEW\_PASS} for 37
filesystem sources and 38 statically enumerated executable units. A
separate \texttt{compile()}-only syntax check accepted the exact locked
Python snapshots without importing or executing

them. The later three-path PREP authorization target and direct- child
review reached only
\texttt{PREP\_AUTHORIZATION\_EXACT\_TARGET\_REVIEW\_PASS}. Before PREP,
static checking found that the frozen \texttt{.el9\_8} Python leader had
been replaced by \texttt{.el9\_8.2}; PREP was not invoked and START

was not created. The additive host-rebind erratum target/review reached
only \texttt{HOST\_IDENTITY\_REBIND\_PROTOCOL\_REVIEW\_PASS}: host
identity is not rebound, the old invocation authority is
nontransferable, and the old probe is superseded and must never be
reused. The later HOST-I-AUTH target/review reached

only \texttt{AUTHORIZATION\_EXACT\_TARGET\_REVIEW\_PASS}; no five-source
or seven-derived-object rebind occurred, host identity is still not
rebound, and no Python or PREP ran. Its later HOST-I instruction stopped
before any source target because the preserved seven-component authority
closure could not admit

the required eighth component. The additive source-authority erratum
target/review reached only
\texttt{HOST\_I\_SOURCE\_AUTHORITY\_ERRATUM\_REVIEW\_PASS}. A
subsequently authorized 14-path source-static target changed the five
registered sources and seven derived authority objects without executing
Python, build, PREP, or data access. Its

exact direct-child review failed with one HIGH because the binding and
crosswalk cite a false SHA-256 for the governing erratum review. The
terminal status is
\texttt{SOURCE\_STATIC\_TARGET\_REVIEW\_FAIL\_AUTHORITY\_IDENTITY\_MISMATCH}
and the failed target itself remains immutable. The later
correction-only repair protocol

target/review \texttt{ddfef99/@b1a7429} passed with zero LOW-or-higher
findings. Exact five-path R1 target/review \texttt{e17f887/@e10bde7}
then passed with zero LOW-or-higher findings after reproducing the four
artifact substitutions, unchanged source/authority closure, and all 13
registered host objects. \texttt{HOST\_IDENTITY\_REBOUND} is established
only for

that registered 13-object bundle; whole-host identity is unestablished.
Source-history protocol target/review \texttt{dcaed57/@e98a3e4} then
passed with zero LOW-or-higher findings, freezing a strict
four-legacy/six-active epoch rule and rejecting adaptive or either-epoch
matching. The first authorized \texttt{A4-V2-SOURCE-HISTORY-I} stopped
before target because

required runtime-schema and maximal-witness changes had no legitimate
semantic, output, or ledger delta. Path-closure erratum target/review
\texttt{150e5b3/@7749491} then passed with zero LOW-or-higher findings,
freezing a necessary-and-sufficient eight-path/five-derived-object
cascade while preserving the two unchanged runtime artifacts exactly.
Exact R1 target/review

\texttt{e9b5c83/@a1198c4} then passed static review with zero
LOW-or-higher findings after changing only the one verifier source plus
the five derived authority objects and two instruction files. The
reviewer reproduced the five legacy overrides, strict
four-legacy/six-active role partition, sole

ninth protocol component, and canonical 37-source tree while the runtime
schema and maximal witness remained byte-identical. This corrects the
known source-history rule only in reviewed static source; no Python,
syntax/import, build, verifier, PREP, data, or scientific execution
occurred, so

cache-verifier/PAR readiness remains unestablished. Fresh R1 PREP
authorization/review remained documentation governance. Its one later
authorized probe/invocation was consumed and stopped at START-only
\texttt{CRASH\_OR\_UNKNOWN}, before the PREP body and without scientific
evidence. Reviewed crash-record protocol \texttt{1982865/@8bec546}
closes only

the missing future record topology. Actual CRASH-I, source repair,
retry, CACHE-BIND, PAR-R1, SRUN, real-base reads, and SAQ integration
remain unauthorized.
\end{frame}

\begin{frame}[fragile,allowframebreaks]{1. Meeting Goal}
\protect\phantomsection\label{1-meeting-goal}
This is not a presentation of a finished method.

The goal is to explain four research attempts in enough detail to decide
what they teach us, what they rule out, and what evidence is still
missing:

\begin{Shaded}
\begin{Highlighting}[]
\NormalTok{Attempt 1: Is PCA a practical limitation of SAQ?}
\NormalTok{           1A. Replace full{-}D raw{-}data PCA with another isometric transform.}
\NormalTok{           1B. Physically project D {-}\textgreater{} d and summarize the omitted tail.}

\NormalTok{Attempt 2: Does histogram{-}exact scalar{-}codebook optimization improve the}
\NormalTok{           quantization and retrieval result beyond Lloyd training, and what}
\NormalTok{           remains after recent ASQ and exact{-}1D{-}DP prior work?}
\end{Highlighting}
\end{Shaded}

\begin{Shaded}
\begin{Highlighting}[]

\NormalTok{Attempt 3: Does paper{-}exact distance quality change conclusions based on}
\NormalTok{           exact{-}identifier Recall@k?}
\NormalTok{Attempt 4: Can arbitrary integer scalar cardinalities use a fixed B\_g{-}bit}
\NormalTok{           group word more effectively than power{-}of{-}two cardinalities while}
\NormalTok{           retaining one random{-}access word and one lookup per group?}
\end{Highlighting}
\end{Shaded}

Main message:

\begin{Shaded}
\begin{Highlighting}[]
\NormalTok{Attempts 1{-}{-}3 found useful local evidence, but none established a new method.}
\NormalTok{For Attempt 2, the local ANN audit remains informative, while the inner exact{-}}
\NormalTok{DP primitive is established prior work and cannot be claimed as the}
\NormalTok{contribution. Attempt 4\textquotesingle{}s A4{-}1S formulation is a completed synthetic pipeline{-}}
\NormalTok{cost falsification study, not a natural{-}data, ANN, or systems result. A4 V2}
\NormalTok{preserves that stop and preregisters a new cost/evidence boundary. Its}
\NormalTok{reporting{-}closure erratum is reviewed protocol evidence, not a feasibility}
\end{Highlighting}
\end{Shaded}

\begin{Shaded}
\begin{Highlighting}[]
\NormalTok{result. Its corrected source and manifest passed independent static review,}
\NormalTok{but the one authorized PAR invocation rejected the \textasciigrave{}/bin/python\textasciigrave{} symlink}
\NormalTok{before any B/P worker. No build or parity fixture ran, so this adds no}
\NormalTok{quantization evidence. A later exact source repair passed static review only;}
\NormalTok{it has not been executed and does not change the terminal PAR result. CACHE{-}P}
\NormalTok{then froze a governance{-}only isolation protocol after reviewing cache semantics}
\NormalTok{and frozen source without opening the quarantine. It rejects direct PAR{-}R1.}
\end{Highlighting}
\end{Shaded}

\begin{Shaded}
\begin{Highlighting}[]
\NormalTok{CACHE{-}I later implemented and statically closed only that generic governance}
\NormalTok{layer; its review and syntax{-}only check add no parity, cost, or scientific}
\NormalTok{evidence. A later exact PREP authorization record passed independent review,}
\NormalTok{but that is authorization{-}document governance only and authorizes no clone,}
\NormalTok{preparation, or execution. The subsequent host{-}rebind protocol and HOST{-}I{-}AUTH}
\NormalTok{reviews likewise remain documentation governance: they rebind neither source}
\NormalTok{nor host and authorize no Python or PREP. The attempted HOST{-}I source edge was}
\end{Highlighting}
\end{Shaded}

\begin{Shaded}
\begin{Highlighting}[]
\NormalTok{stopped before edits on a seven{-}versus{-}eight contract contradiction; its}
\NormalTok{reviewed additive erratum repairs only that source{-}authority contract and}
\NormalTok{enabled a later source{-}static target. That target\textquotesingle{}s direct{-}child review found}
\NormalTok{one HIGH authority{-}identity mismatch, so that target remains failed evidence.}
\NormalTok{A later correction{-}only target passed exact review and establishes rebound}
\NormalTok{only for the registered 13{-}object bundle. It ran no execution and does not}
\NormalTok{establish whole{-}host, PREP, cache{-}verifier/PAR, or scientific readiness. A}
\end{Highlighting}
\end{Shaded}

\begin{Shaded}
\begin{Highlighting}[]
\NormalTok{later reviewed source{-}history protocol specifies a fail{-}closed two{-}epoch}
\NormalTok{repair. Its first implementation authorization stopped before target on an}
\NormalTok{unsatisfiable ten{-}path projection; the reviewed path{-}closure erratum freezes}
\NormalTok{the corrected eight{-}path boundary. The separately authorized exact R1 target}
\NormalTok{then implemented only that static correction and passed direct{-}child review;}
\NormalTok{it did not run the verifier or establish downstream readiness.}
\NormalTok{Fresh dormant PREP authorization target/review \textasciigrave{}cd1757e/@eb1b893\textasciigrave{} then passed}
\end{Highlighting}
\end{Shaded}

\begin{Shaded}
\begin{Highlighting}[]
\NormalTok{exact direct{-}child documentation review with zero LOW{-}or{-}higher findings. The}
\NormalTok{later unique immediate probe passed and the unique invocation was consumed,}
\NormalTok{but it stopped at durable START{-}only \textasciigrave{}CRASH\_OR\_UNKNOWN\textasciigrave{} before the PREP body.}
\NormalTok{Crash{-}record protocol target/review \textasciigrave{}1982865/@8bec546\textasciigrave{} subsequently passed;}
\NormalTok{it forms governance for a future terminal record, not execution evidence.}
\end{Highlighting}
\end{Shaded}

Speaker notes:

\begin{itemize}
\tightlist
\item
  The purpose is research selection, not presenting every experiment as
  a win.
\item
  Attempt 1 belongs to the SAQ repository.
\item
  Attempt 2 comes from the sibling \texttt{vectordb} repository and
  tests a more classical scalar-quantization premise. The new
  related-work audit separates the local two-DP pipeline from the
  already-established inner optimizer.
\item
  Attempt 3 returns to SAQ and asks whether Recall understated the
  geometric quality of any previously rejected result set.
\item
  Attempt 4 is problem-first and fixed-rate. A4-1S passed implementation
  parity but failed its pipeline-cost gate. A4 V2 has a reviewed parent
  protocol and additive timing-closure erratum. Its PAR attempt failed
  only at executable-identity admission. I-R1 repairs that source

  path but remains unexecuted. CACHE-P selects staged sanitized remote
  isolation and rejects direct PAR-R1. CACHE-I closes its generic
  source/schema layer under static review only; none of these facts is
  quantization-quality, natural-data, ANN, or systems evidence. The
  later PREP
\item
  authorization record also passed only documentation review. A
  subsequent static host mismatch prevented PREP and START; the reviewed
  host-rebind erratum is protocol governance only, and the old delayed
  probe is superseded and may never run or be reused.

  The later HOST-I-AUTH review freezes only the contract and projection
  for a possible future source target; it grants no HOST-I authority and
  performed no source/authority rebind, Python, or PREP. That attempted
  source target then stopped before edits on
\item
  an exact seven-versus-eight component contradiction. The reviewed
  source-authority erratum admits the eighth component. The later
  14-path HOST-I source-static target was created, but its exact review
  failed with one HIGH because two authority documents bind the wrong
  governing-review SHA-256.

  A later reviewed five-path R1 correction repaired only that provenance
  defect and passed independent review across all 13 registered host
  objects. This is bounded bundle identity, not whole- host or execution
  readiness. The later source-history protocol passed independent
\item
  review, but its first implementation authorization stopped before
  target on an unsatisfiable ten-path closure. The reviewed path-closure
  erratum froze a strict eight-path/five-derived-object repair. The
  later exact R1 target passed direct-child static review and corrects
  that known source rule

  only; no verifier, PREP, cache/PAR readiness, or scientific claim
  follows. Fresh dormant PREP authorization target/review
  \texttt{cd1757e/@eb1b893} subsequently passed exact documentation
  review. Its unique probe/invocation was later consumed and left only
  START plus empty captures, so the status is
\item
  \texttt{CRASH\_OR\_UNKNOWN\ /\ NO\_SCIENTIFIC\_DECISION}. Crash-record
  protocol \texttt{1982865/@8bec546} passed review, but CRASH-I, repair,
  retry, CACHE-BIND, and PAR-R1 remain unauthorized.
\end{itemize}
\end{frame}

\begin{frame}[fragile,allowframebreaks]{2. Current Research Position}
\protect\phantomsection\label{2-current-research-position}
Project-level decision before this deck:

\begin{Shaded}
\begin{Highlighting}[]
\NormalTok{SAQ{-}centric incremental method search: STOP}
\NormalTok{SAQ as a strong baseline:              RETAIN}
\NormalTok{Attempt 4 A4{-}1S formulation:           TERMINAL COST STOP}
\NormalTok{Attempt 4 A4 V2 PAR:                   ARTIFACT\_INVALID; NO SCIENTIFIC DECISION}
\NormalTok{Attempt 4 A4 V2 I{-}R1:                  SOURCE\_REPAIR\_STATIC\_REVIEW\_PASS ONLY}
\NormalTok{Attempt 4 A4 V2 CACHE{-}P target:        EXACT\_TARGET\_INDEPENDENT\_REVIEW\_PASS}
\NormalTok{Attempt 4 A4 V2 CACHE{-}I:               GENERIC\_CACHE\_POLICY\_SOURCE\_STATIC\_REVIEW\_PASS}
\end{Highlighting}
\end{Shaded}

\begin{Shaded}
\begin{Highlighting}[]
\NormalTok{Attempt 4 A4 V2 PREP authorization:    PREP\_AUTHORIZATION\_EXACT\_TARGET\_REVIEW\_PASS}
\NormalTok{Attempt 4 A4 V2 HOST{-}P erratum:        HOST\_IDENTITY\_REBIND\_PROTOCOL\_REVIEW\_PASS}
\NormalTok{Attempt 4 A4 V2 HOST{-}I{-}AUTH:           AUTHORIZATION\_EXACT\_TARGET\_REVIEW\_PASS}
\NormalTok{Attempt 4 A4 V2 HOST{-}I{-}ERRATUM:        HOST\_I\_SOURCE\_AUTHORITY\_ERRATUM\_REVIEW\_PASS}
\NormalTok{Attempt 4 A4 V2 host rebound:          ESTABLISHED FOR REGISTERED 13{-}OBJECT BUNDLE ONLY}
\NormalTok{Attempt 4 A4 V2 HOST{-}I:                SOURCE\_STATIC\_TARGET\_REVIEW\_FAIL\_AUTHORITY\_IDENTITY\_MISMATCH}
\NormalTok{Attempt 4 A4 V2 HOST{-}I{-}R1{-}P:           HOST\_I\_R1\_CORRECTION\_ONLY\_REPAIR\_PROTOCOL\_REVIEW\_PASS}
\end{Highlighting}
\end{Shaded}

\begin{Shaded}
\begin{Highlighting}[]
\NormalTok{Attempt 4 A4 V2 HOST{-}I{-}R1:             HOST\_I\_R1\_CORRECTION\_ONLY\_REPAIR\_REVIEW\_PASS}
\NormalTok{Attempt 4 A4 V2 SOURCE{-}HISTORY{-}P:      SOURCE\_HISTORY\_EPOCH\_CORRECTION\_PROTOCOL\_REVIEW\_PASS}
\NormalTok{Attempt 4 A4 V2 SOURCE{-}HISTORY{-}I:      STOPPED PRE{-}TARGET; TEN{-}PATH CLOSURE UNSATISFIABLE}
\NormalTok{Attempt 4 A4 V2 SOURCE{-}HISTORY{-}I{-}R1{-}P: SOURCE\_HISTORY\_I\_R1\_PATH\_CLOSURE\_ERRATUM\_REVIEW\_PASS}
\NormalTok{Attempt 4 A4 V2 SOURCE{-}HISTORY{-}I{-}R1:   SOURCE\_HISTORY\_EPOCH\_CORRECTION\_SOURCE\_STATIC\_REVIEW\_PASS}
\NormalTok{Attempt 4 A4 V2 CACHE{-}PREP{-}R1{-}AUTH:    PREP\_AUTHORIZATION\_EXACT\_TARGET\_REVIEW\_PASS}
\NormalTok{Attempt 4 A4 V2 PREP{-}ISO{-}R1:           CRASH\_OR\_UNKNOWN / START\_ONLY / NO SCIENTIFIC DECISION}
\end{Highlighting}
\end{Shaded}

\begin{Shaded}
\begin{Highlighting}[]
\NormalTok{Attempt 4 A4 V2 PREP{-}ISO{-}R1{-}CRASH{-}P:   START\_ONLY\_CRASH\_RECORD\_PROTOCOL\_REVIEW\_PASS}
\NormalTok{Attempt 4 A4 V2 CRASH{-}I/BIND:          NOT AUTHORIZED / NOT RUN}
\NormalTok{Attempt 4 A4 V2 cache/PAR readiness:   NOT ESTABLISHED; CORRECTED SOURCE NOT EXECUTED}
\NormalTok{Attempt 4 A4 V2 direct PAR{-}R1:         NOT AUTHORIZED / NOT RUN}
\NormalTok{Attempt 4 A4 V2 SRUN/data:             NOT AUTHORIZED / NOT RUN}
\NormalTok{Attempt 4 A4 V2 portfolio:             CLOSED / STOP\_NO\_FURTHER\_ARTIFACT\_RECOVERY}
\NormalTok{SAQ/index/search integration:          NOT AUTHORIZED}
\end{Highlighting}
\end{Shaded}

Why still present these attempts:

\begin{enumerate}
\tightlist
\item
  Attempts 1-\/-3 test three intuitive claims that may arise in a
  meeting.
\item
  They distinguish optimization guarantees from retrieval guarantees.
\item
  They provide quantitative stop evidence instead of relying on
  intuition.
\item
  Attempt 3 tests whether an evaluation choice, rather than a quantizer
  mechanism, caused one earlier negative conclusion.
\item
  Attempt 4 shows both a preregistered synthetic cost stop and how a
  later protocol must preserve that result while redefining its FOM and
  evidence boundary before any new authorization.
\end{enumerate}

Unsafe interpretation:

\begin{Shaded}
\begin{Highlighting}[]
\NormalTok{PCA is universally optimal, or exact DP never helps ANN.}
\end{Highlighting}
\end{Shaded}

Safe interpretation:

\begin{Shaded}
\begin{Highlighting}[]
\NormalTok{The specific PCA{-}replacement, lossy{-}projection, and histogram{-}DP mechanisms}
\NormalTok{tested here do not supply a sufficiently general, overhead{-}controlled method.}
\NormalTok{The 2026 ASQ paper and pinned public code do not expose an integrated Attempt 2}
\NormalTok{pipeline, but the paper and earlier exact{-}1D{-}clustering work remove the inner}
\NormalTok{DP as a novelty claim.}
\NormalTok{Paper{-}exact distance quality changes one GIST operating{-}point interpretation,}
\NormalTok{but does not establish a new mechanism or a replicated method advantage.}
\end{Highlighting}
\end{Shaded}

\begin{Shaded}
\begin{Highlighting}[]
\NormalTok{For Attempt 4, allowing positive{-}integer cardinalities only enlarges the}
\NormalTok{dyadic feasible set for a factorized reconstruction objective. Mixed{-}radix}
\NormalTok{addressing and non{-}power{-}of{-}two scalar alphabets are prior art, and imply}
\NormalTok{neither ANN improvement nor a systems advantage. A4 V2 only asks whether a}
\NormalTok{complete four{-}arm synthetic comparative instrument fits an internal cost cap;}
\NormalTok{the old \textasciigrave{}5/2\textasciigrave{} real{-}data projection does not transfer.}
\end{Highlighting}
\end{Shaded}
\end{frame}

\begin{frame}[fragile,allowframebreaks]{3. Attempt Map And Evidence
Status}
\protect\phantomsection\label{3-attempt-map-and-evidence-status}
\begin{itemize}
\tightlist
\item
  \textbf{Attempt}: 1A. Full-D transform replacement

  \textbf{Core question}: Does residual PCA or another isometric basis
  improve SAQ estimator behavior over raw-data PCA?
\item
  \textbf{Main regimes}: GIST sample50k, then preregistered CIFAR60K
  replication

  \textbf{Decision / current status}: Closed after replication failure
\item
  \textbf{Attempt}: 1B. Lossy \texttt{D\ -\textgreater{}\ d} projection

  \textbf{Core question}: Can a PCA head plus a compact tail surrogate
  beat native full-D SAQ?
\item
  \textbf{Main regimes}: GIST sample50k,
  \texttt{960\ -\textgreater{}\ 576}, favorable exact surrogate

  \textbf{Decision / current status}: Gate A failed; projected SAQ not
  built
\item
  \textbf{Attempt}: 2. Exact scalar-codebook DP

  \textbf{Core question}: Does histogram-exact 1D DP improve shared
  dimensionwise scalar quantization over Lloyd, and is any method
  novelty left after prior work?
\item
  \textbf{Main regimes}: audio, PCA CIFAR60K, PCA DEEP1M; arXiv/code
  audit

  \textbf{Decision / current status}: Stronger offline baseline; inner
  DP is prior art; no stable recall dominance
\item
  \textbf{Attempt}: 3. Distance-quality re-evaluation

  \textbf{Core question}: Does \texttt{1/Ratio@k} change a frozen
  Recall-based Pareto conclusion?
\item
  \textbf{Main regimes}: GIST sample100k B=4; DEEP sample100k B=4/B=5
  controls

  \textbf{Decision / current status}: Closed as metric-sensitivity
  evidence
\item
  \textbf{Attempt}: 4. Arbitrary-cardinality fixed-rate words

  \textbf{Core question}: Does the power-of-two restriction waste
  material capacity at unchanged fixed payload and lookup granularity?
\item
  \textbf{Main regimes}: A4-0 synthetic witness; terminal A4-1S
  correctness/cost gate; reviewed A4 V2 protocol/source, terminal
  pre-build PAR identity failure, static-only corrections, CACHE
  governance, and the consumed START-only PREP plus reviewed
  crash-record protocol; registered real-base gates remain unauthorized

  \textbf{Decision / current status}: preserve A4-1S
  \texttt{NO\_GO\_EXACT\_SOLVER\_COST} and A4-V2-PAR
  \texttt{ARTIFACT\_INVALID\ /\ NO\_SCIENTIFIC\_DECISION};
  source-history R1 passed static review only; the later unique PREP
  invocation stopped at
  \texttt{CRASH\_OR\_UNKNOWN\ /\ START\_ONLY\ /\ NO\_SCIENTIFIC\_DECISION};
  crash-record protocol \texttt{1982865/@8bec546} passed, but CRASH-I,
  repair/retry, CACHE-BIND, and PAR-R1 remain unauthorized
\end{itemize}

Speaker notes:

\begin{itemize}
\tightlist
\item
  Attempt 1A and 1B are related but not equivalent.
\item
  A full-D orthogonal transform preserves exact L2 geometry.
\item
  A physical \texttt{D\ -\textgreater{}\ d} projection changes the exact
  distance unless the omitted interaction is somehow reconstructed.
\item
  Attempt 4 changes the feasible scalar-product cardinalities inside a
  fixed word; it does not yet change SAQ, CAQ, the index, or the search
  path.
\end{itemize}
\end{frame}

\begin{frame}[fragile,allowframebreaks]{4. Minimal SAQ Background}
\protect\phantomsection\label{4-minimal-saq-background}
SAQ compresses database vectors for approximate L2 search.

High-level pipeline:

\begin{Shaded}
\begin{Highlighting}[]
\NormalTok{raw vectors}
\NormalTok{  {-}\textgreater{} full{-}dimensional PCA}
\NormalTok{  {-}\textgreater{} IVF residuals}
\NormalTok{  {-}\textgreater{} contiguous 64{-}dimensional segments}
\NormalTok{  {-}\textgreater{} nonuniform bits across segments}
\NormalTok{  {-}\textgreater{} CAQ encoding inside each positive{-}bit segment}
\NormalTok{  {-}\textgreater{} progressive distance estimation during search}
\end{Highlighting}
\end{Shaded}

Example GIST \texttt{B=4} plan:

\begin{Shaded}
\begin{Highlighting}[]
\NormalTok{64@11 | 192@6 | 320@4 | 256@2 | 128@0}
\end{Highlighting}
\end{Shaded}

Meaning:

\begin{itemize}
\tightlist
\item
  \textbf{Segment}: 1

  \textbf{PCA coordinates}: 0-63
\item
  \textbf{Dimensions}: 64

  \textbf{Bits/dimension}: 11
\item
  \textbf{Segment}: 2

  \textbf{PCA coordinates}: 64-255
\item
  \textbf{Dimensions}: 192

  \textbf{Bits/dimension}: 6
\item
  \textbf{Segment}: 3

  \textbf{PCA coordinates}: 256-575
\item
  \textbf{Dimensions}: 320

  \textbf{Bits/dimension}: 4
\item
  \textbf{Segment}: 4

  \textbf{PCA coordinates}: 576-831
\item
  \textbf{Dimensions}: 256

  \textbf{Bits/dimension}: 2
\item
  \textbf{Segment}: 5

  \textbf{PCA coordinates}: 832-959
\item
  \textbf{Dimensions}: 128

  \textbf{Bits/dimension}: 0
\end{itemize}

PCA is therefore not a cosmetic preprocessing step. It determines which
subspaces receive high bits, low bits, or no inner-product code.
\end{frame}

\begin{frame}[fragile,allowframebreaks]{5. Common Variables}
\protect\phantomsection\label{5-common-variables}
\begin{itemize}
\tightlist
\item
  \textbf{Symbol}: \texttt{N}

  \textbf{Meaning}: number of base vectors
\item
  \textbf{Values used here}: GIST 50,000/100,000; CIFAR 60,000; DEEP
  100,000/1,000,000
\item
  \textbf{Symbol}: \texttt{D}

  \textbf{Meaning}: original dimension
\item
  \textbf{Values used here}: GIST 960; CIFAR 512; DEEP 256; audio 192
\item
  \textbf{Symbol}: \texttt{d}

  \textbf{Meaning}: physically retained dimension after lossy projection
\item
  \textbf{Values used here}: GIST 576
\item
  \textbf{Symbol}: \texttt{K\_IVF}

  \textbf{Meaning}: number of IVF clusters
\item
  \textbf{Values used here}: 512 in Attempts 1 and 3
\item
  \textbf{Symbol}: \texttt{B}

  \textbf{Meaning}: nominal average bits/dimension
\item
  \textbf{Values used here}: 4 in Attempt 1; 2/4/8 in Attempt 2; 4/5 in
  Attempt 3
\item
  \textbf{Symbol}: \texttt{Q}

  \textbf{Meaning}: number of held-out evaluation queries
\item
  \textbf{Values used here}: 128 or 1,000 in Attempts 1 and 3
\item
  \textbf{Symbol}: \texttt{H}

  \textbf{Meaning}: equal-count histogram bins for scalar DP
\item
  \textbf{Values used here}: default 256
\item
  \textbf{Symbol}: \texttt{K\_b}

  \textbf{Meaning}: scalar centroids at bitwidth \texttt{b}
\item
  \textbf{Values used here}: up to \texttt{2\^{}b}
\item
  \textbf{Symbol}: \texttt{B\_g}

  \textbf{Meaning}: fixed bits in one two-factor word for Attempt 4
\item
  \textbf{Values used here}: 4 or 8
\item
  \textbf{Symbol}: \texttt{S}

  \textbf{Meaning}: joint-state capacity of one fixed word,
  \texttt{2\^{}B\_g}
\item
  \textbf{Values used here}: 16 or 256
\item
  \textbf{Symbol}: \texttt{K\_j}

  \textbf{Meaning}: scalar levels assigned to factor \texttt{j}
\item
  \textbf{Values used here}: dyadic or any positive integer
\item
  \textbf{Symbol}: \texttt{r}

  \textbf{Meaning}: scalar factors grouped into one fixed word
\item
  \textbf{Values used here}: 2 in A4-1
\item
  \textbf{Symbol}: \texttt{w\_j}

  \textbf{Meaning}: allocation weight for dimension \texttt{j}
\item
  \textbf{Values used here}: 1 or historical rank-boundary weight
\item
  \textbf{Symbol}: \texttt{R}

  \textbf{Meaning}: full-D orthogonal transform matrix
\item
  \textbf{Values used here}: transform-dependent
\item
  \textbf{Symbol}: \texttt{k}

  \textbf{Meaning}: number of returned neighbors evaluated
\item
  \textbf{Values used here}: 100 in Attempt 3
\item
  \textbf{Symbol}: \texttt{nprobe}

  \textbf{Meaning}: IVF cells scanned per query
\item
  \textbf{Values used here}: GIST 20-\/-400; DEEP 50-\/-400
\item
  \textbf{Symbol}: \texttt{d\_i(q)}

  \textbf{Meaning}: true Euclidean distance to exact neighbor at
  distance rank \texttt{i}
\item
  \textbf{Values used here}: recomputed in float64
\item
  \textbf{Symbol}: \texttt{d\_tilde\_i(q)}

  \textbf{Meaning}: true Euclidean distance at rank \texttt{i} within
  the returned set
\item
  \textbf{Values used here}: recomputed in float64
\end{itemize}

Important notation distinction:

\begin{Shaded}
\begin{Highlighting}[]
\NormalTok{K\_IVF = IVF clusters}
\NormalTok{K\_b   = scalar quantizer centroids}
\end{Highlighting}
\end{Shaded}
\end{frame}

\begin{frame}[fragile,allowframebreaks]{6. Attempt 1: Research Question}
\protect\phantomsection\label{6-attempt-1-research-question}
SAQ uses raw-data PCA because it:

\begin{Shaded}
\begin{Highlighting}[]
\NormalTok{decorrelates coordinates}
\NormalTok{orders coordinates by decreasing variance}
\NormalTok{creates a low{-}variance tail}
\NormalTok{supports contiguous mixed{-}bit segmentation}
\end{Highlighting}
\end{Shaded}

But PCA optimizes variance/reconstruction objectives, not directly:

\begin{Shaded}
\begin{Highlighting}[]
\NormalTok{CAQ angular error}
\NormalTok{progressive{-}prefix distance error}
\NormalTok{top{-}k ranking stability}
\NormalTok{Recall{-}QPS{-}bytes}
\end{Highlighting}
\end{Shaded}

Attempt 1 asks two increasingly aggressive questions:

\begin{Shaded}
\begin{Highlighting}[]
\NormalTok{1A. Can another full{-}D orthogonal transform improve SAQ while preserving L2?}

\NormalTok{1B. If low{-}variance tail coordinates are physically removed, can a compact}
\NormalTok{    tail summary retain enough ranking information to improve work/space?}
\end{Highlighting}
\end{Shaded}

Speaker notes:

\begin{itemize}
\tightlist
\item
  We first test whether the premise is real before training a new
  transform.
\item
  A learned transform is not authorized merely because PCA optimizes a
  different textbook objective.
\end{itemize}
\end{frame}

\begin{frame}[fragile,allowframebreaks]{7. Attempt 1A: Transform
Definitions}
\protect\phantomsection\label{7-attempt-1a-transform-definitions}
Current SAQ transform:

\[T_{raw}(x)=(x-\mu_{raw})R_{raw},\]

where \texttt{R\_raw} contains the eigenvectors of raw base-vector
covariance.

Residual-PCA alternative:

\[r_x=x-c_{cid(x)},\]

\[T_{res}(x)=(x-\mu_{res})R_{res},\]

where \texttt{R\_res} is trained from IVF residual statistics rather
than raw-vector statistics.

For either full-D orthogonal matrix:

\[R^TR=I,\]

so applying the same affine transform to base vectors, queries, and
centroids preserves squared L2 up to float roundoff:

\[\lVert T(q)-T(x)\rVert_2^2=\lVert q-x\rVert_2^2.\]

Controls:

\begin{Shaded}
\begin{Highlighting}[]
\NormalTok{current raw{-}data PCA}
\NormalTok{global residual PCA}
\NormalTok{identity coordinates}
\NormalTok{seeded random orthogonal transform}
\end{Highlighting}
\end{Shaded}
\end{frame}

\begin{frame}[fragile,allowframebreaks]{8. Attempt 1A Running Example}
\protect\phantomsection\label{8-attempt-1a-running-example}
Illustrative two-dimensional dataset:

\begin{Shaded}
\begin{Highlighting}[]
\NormalTok{raw covariance variances:       axis 1 = 100, axis 2 = 1}
\NormalTok{within{-}IVF residual variances:   axis 1 =   4, axis 2 = 9}
\end{Highlighting}
\end{Shaded}

Raw PCA ordering:

\begin{Shaded}
\begin{Highlighting}[]
\NormalTok{axis 1 first, axis 2 second}
\end{Highlighting}
\end{Shaded}

Residual PCA ordering:

\begin{Shaded}
\begin{Highlighting}[]
\NormalTok{axis 2 first, axis 1 second}
\end{Highlighting}
\end{Shaded}

If SAQ assigns more bits to earlier coordinates, the two transforms can
produce different segment importance even though both preserve exact L2.

This motivates the test, but it is not evidence. The empirical question
is:

\begin{Shaded}
\begin{Highlighting}[]
\NormalTok{Does residual ordering improve the actual packed SAQ estimator and ranking,}
\NormalTok{after plan, bytes, and internal rotations are controlled?}
\end{Highlighting}
\end{Shaded}
\end{frame}

\begin{frame}[fragile,allowframebreaks]{9. Attempt 1A Phase 1 Protocol:
GIST}
\protect\phantomsection\label{9-attempt-1a-phase-1-protocol-gist}
\begin{itemize}
\tightlist
\item
  \textbf{Item}: Dataset

  \textbf{Frozen value}: GIST sample50k, \texttt{N=50,000},
  \texttt{D=960}
\item
  \textbf{Item}: IVF

  \textbf{Frozen value}: \texttt{K\_IVF=512}, one recovered canonical
  codebook and fixed assignments
\item
  \textbf{Item}: Budget

  \textbf{Frozen value}: \texttt{B=4}
\item
  \textbf{Item}: Queries

  \textbf{Frozen value}: first 128 held-out queries; evaluation only
\item
  \textbf{Item}: Fixed probes

  \textbf{Frozen value}: top 16 canonical IVF clusters/query
\item
  \textbf{Item}: Fixed candidates

  \textbf{Frozen value}: 442,823 candidate rows/configuration
\item
  \textbf{Item}: Ranking cutoff

  \textbf{Frozen value}: top 100
\item
  \textbf{Item}: Transforms

  \textbf{Frozen value}: current PCA, residual PCA, identity, random
  orthogonal
\item
  \textbf{Item}: Plan controls

  \textbf{Frozen value}: native, frozen-PCA, uniform \texttt{960@4}
\item
  \textbf{Item}: Segment rotations

  \textbf{Frozen value}: logical seeds \texttt{0..9}, plus \texttt{off}
\item
  \textbf{Item}: Configurations

  \textbf{Frozen value}:
  \texttt{4\ transforms\ x\ 3\ plans\ x\ 11\ controls\ =\ 132}
\item
  \textbf{Item}: Inference

  \textbf{Frozen value}: query-paired percentile bootstrap, 10,000
  replicates
\end{itemize}

The runner replays production packed estimators for:

\begin{Shaded}
\begin{Highlighting}[]
\NormalTok{variance estimate}
\NormalTok{every fast prefix}
\NormalTok{every accurate prefix}
\NormalTok{full code}
\end{Highlighting}
\end{Shaded}

It is a fixed-candidate estimator experiment, not an end-to-end QPS run.
\end{frame}

\begin{frame}[fragile,allowframebreaks]{10. Attempt 1A Code And Artifact
Map}
\protect\phantomsection\label{10-attempt-1a-code-and-artifact-map}
Branch:

\begin{Shaded}
\begin{Highlighting}[]
\NormalTok{saq{-}transform{-}analysis}
\NormalTok{evidence checkpoint: 3d94840}
\end{Highlighting}
\end{Shaded}

Key code:

\begin{Shaded}
\begin{Highlighting}[]
\NormalTok{script/prepare\_phase1\_transform\_views.py}
\NormalTok{src/phase1\_transform\_diagnostic.cpp}
\NormalTok{script/summarize\_phase1\_transform.py}
\NormalTok{script/evaluate\_phase1b\_gate.py}
\end{Highlighting}
\end{Shaded}

Primary evidence:

\begin{Shaded}
\begin{Highlighting}[]
\NormalTok{docs/saq\_transform\_phase1\_limitation\_evidence\_2026\_07\_10.md}
\NormalTok{docs/saq\_transform\_phase1b\_external\_replication\_evidence\_2026\_07\_10.md}
\end{Highlighting}
\end{Shaded}

Validity controls include:

\begin{Shaded}
\begin{Highlighting}[]
\NormalTok{held{-}out affine{-}transform reconstruction}
\NormalTok{orthogonality/isometry checks}
\NormalTok{identical IVF assignments and probe lists}
\NormalTok{canonical raw{-}space ranking labels}
\NormalTok{identical serialized bytes for matched transform pairs}
\NormalTok{distinct internal{-}rotation streams}
\end{Highlighting}
\end{Shaded}
\end{frame}

\begin{frame}[fragile,allowframebreaks]{11. GIST Native Plans And Stage
Results}
\protect\phantomsection\label{11-gist-native-plans-and-stage-results}
Current PCA and residual PCA select the same plan and serialized bytes:

\begin{Shaded}
\begin{Highlighting}[]
\NormalTok{64@11 | 192@6 | 320@4 | 256@2 | 128@0}
\NormalTok{serialized index = 34,201,289 bytes}
\end{Highlighting}
\end{Shaded}

Seed-averaged mean-query results:

\begin{itemize}
\tightlist
\item
  \textbf{Transform}: current PCA

  \textbf{Stage}: fast all
\item
  \textbf{Logical bytes/candidate}: 124

  \textbf{RMSE}: 0.386347
\item
  \textbf{Top-100 agreement}: 0.593133

  \textbf{Boundary inversion}: \texttt{1.9294e-2}
\item
  \textbf{Transform}: residual PCA

  \textbf{Stage}: fast all
\item
  \textbf{Logical bytes/candidate}: 124

  \textbf{RMSE}: 0.386987
\item
  \textbf{Top-100 agreement}: 0.589813

  \textbf{Boundary inversion}: \texttt{1.9522e-2}
\item
  \textbf{Transform}: current PCA

  \textbf{Stage}: accurate prefix 1
\item
  \textbf{Logical bytes/candidate}: 212

  \textbf{RMSE}: 0.0818108
\item
  \textbf{Top-100 agreement}: 0.904930

  \textbf{Boundary inversion}: \texttt{9.5544e-4}
\item
  \textbf{Transform}: residual PCA

  \textbf{Stage}: accurate prefix 1
\item
  \textbf{Logical bytes/candidate}: 212

  \textbf{RMSE}: 0.0813185
\item
  \textbf{Top-100 agreement}: 0.906188

  \textbf{Boundary inversion}: \texttt{9.3807e-4}
\item
  \textbf{Transform}: current PCA

  \textbf{Stage}: full
\item
  \textbf{Logical bytes/candidate}: 508

  \textbf{RMSE}: 0.00168088
\item
  \textbf{Top-100 agreement}: 0.994617

  \textbf{Boundary inversion}: \texttt{4.0869e-6}
\item
  \textbf{Transform}: residual PCA

  \textbf{Stage}: full
\item
  \textbf{Logical bytes/candidate}: 508

  \textbf{RMSE}: 0.00167530
\item
  \textbf{Top-100 agreement}: 0.994758

  \textbf{Boundary inversion}: \texttt{4.1323e-6}
\end{itemize}

Residual PCA improves distance RMSE at accurate/full stages, but worsens
the fast stage.
\end{frame}

\begin{frame}[fragile,allowframebreaks]{12. GIST Paired Effects And
Gate}
\protect\phantomsection\label{12-gist-paired-effects-and-gate}
All deltas are \texttt{residual\ PCA\ -\ current\ PCA}.

\begin{itemize}
\tightlist
\item
  \textbf{Stage}: fast all

  \textbf{Mean-query RMSE delta}: \texttt{+6.398e-4}
\item
  \textbf{Relative effect}: worse

  \textbf{Ranking interpretation}: top-100 and inversion also worse
\item
  \textbf{Stage}: accurate prefix 1

  \textbf{Mean-query RMSE delta}: \texttt{-4.922e-4}
\item
  \textbf{Relative effect}: \texttt{-0.602\%}

  \textbf{Ranking interpretation}: top-100/inversion CIs include zero
\item
  \textbf{Stage}: full

  \textbf{Mean-query RMSE delta}: \texttt{-5.572e-6}
\item
  \textbf{Relative effect}: \texttt{-0.331\%}

  \textbf{Ranking interpretation}: no reliable top-100 or inversion gain
\end{itemize}

Additional controls:

\begin{Shaded}
\begin{Highlighting}[]
\NormalTok{accurate{-}prefix RMSE improvement: 10/10 rotation seeds}
\NormalTok{rotation{-}off accurate RMSE:       favorable}
\NormalTok{fast prefixes:                    consistently worse}
\NormalTok{identity/random with native plan: faster{-}looking fast stage, much worse full code}
\NormalTok{identity/random with frozen plan: substantially worse ranking}
\end{Highlighting}
\end{Shaded}

Phase 1 interpretation:

\begin{Shaded}
\begin{Highlighting}[]
\NormalTok{A small estimator{-}objective mismatch exists on GIST, but no progressive or}
\NormalTok{ranking Pareto improvement is established. Standard residual PCA already}
\NormalTok{explains the signal, so a learned transform is not yet justified.}
\end{Highlighting}
\end{Shaded}

The only authorized continuation was one preregistered external
replication.
\end{frame}

\begin{frame}[fragile,allowframebreaks]{13. Attempt 1A Phase 1b
Protocol: CIFAR60K}
\protect\phantomsection\label{13-attempt-1a-phase-1b-protocol-cifar60k}
The replication freezes the narrow current-PCA versus residual-PCA
comparison.

\begin{itemize}
\tightlist
\item
  \textbf{Item}: Dataset

  \textbf{Frozen value}: CIFAR60K, \texttt{N=60,000}, \texttt{D=512}
\item
  \textbf{Item}: IVF

  \textbf{Frozen value}: \texttt{K\_IVF=512}, same recovered
  codebook/assignments
\item
  \textbf{Item}: Budget

  \textbf{Frozen value}: \texttt{B=4}
\item
  \textbf{Item}: Queries

  \textbf{Frozen value}: all 1,000 held-out queries
\item
  \textbf{Item}: Fixed probes

  \textbf{Frozen value}: 16/query
\item
  \textbf{Item}: Candidate pairs

  \textbf{Frozen value}: 2,130,639/configuration
\item
  \textbf{Item}: Ranking cutoff

  \textbf{Frozen value}: top 100
\item
  \textbf{Item}: Plan

  \textbf{Frozen value}: `64@9
\item
  \textbf{Item}: Transforms

  \textbf{Frozen value}: current PCA, residual PCA
\item
  \textbf{Item}: Rotation controls

  \textbf{Frozen value}: 10 paired seeds plus \texttt{off}
\item
  \textbf{Item}: Configurations

  \textbf{Frozen value}: \texttt{2\ x\ 11\ =\ 22}
\item
  \textbf{Item}: Bootstrap

  \textbf{Frozen value}: 10,000 query-paired replicates
\end{itemize}

Registered success required:

\begin{Shaded}
\begin{Highlighting}[]
\NormalTok{confidence interval in the favorable direction}
\NormalTok{at least 8/10 seeds favor residual PCA}
\NormalTok{rotation{-}off direction also favorable}
\NormalTok{no harm at the fast stage}
\end{Highlighting}
\end{Shaded}
\end{frame}

\begin{frame}[fragile,allowframebreaks]{14. CIFAR60K Replication Result}
\protect\phantomsection\label{14-cifar60k-replication-result}
All deltas are \texttt{residual\ PCA\ -\ current\ PCA}; lower RMSE is
better.

\begin{itemize}
\tightlist
\item
  \textbf{Stage}: fast all

  \textbf{Current RMSE}: 0.13355750
\item
  \textbf{Residual RMSE}: 0.13359980

  \textbf{Relative delta}: \texttt{+0.0317\%} worse
\item
  \textbf{Seeds favoring residual}: 3/10
\item
  \textbf{Stage}: accurate prefix 1

  \textbf{Current RMSE}: 0.03010816
\item
  \textbf{Residual RMSE}: 0.03010955

  \textbf{Relative delta}: \texttt{+0.00463\%} worse
\item
  \textbf{Seeds favoring residual}: 4/10
\item
  \textbf{Stage}: full

  \textbf{Current RMSE}: 0.001108523
\item
  \textbf{Residual RMSE}: 0.001108262

  \textbf{Relative delta}: \texttt{-0.0235\%}
\item
  \textbf{Seeds favoring residual}: 5/10
\end{itemize}

Confidence and controls:

\begin{Shaded}
\begin{Highlighting}[]
\NormalTok{accurate{-}prefix CI spans zero}
\NormalTok{full{-}code CI spans zero}
\NormalTok{rotation{-}off accurate/full directions are unfavorable}
\NormalTok{fast{-}stage degradation is statistically supported}
\NormalTok{top{-}100 and boundary{-}inversion intervals span zero or point worse}
\end{Highlighting}
\end{Shaded}

Registered decision:

\begin{Shaded}
\begin{Highlighting}[]
\NormalTok{close\_one\_dataset\_estimator\_effect}
\end{Highlighting}
\end{Shaded}

The GIST effect does not replicate across the second frozen spectral
regime.
\end{frame}

\begin{frame}[fragile,allowframebreaks]{15. Full-D Transform Complexity
And Overhead}
\protect\phantomsection\label{15-full-d-transform-complexity-and-overhead}
For a dense full-D PCA implementation:

\begin{Shaded}
\begin{Highlighting}[]
\NormalTok{covariance accumulation: O(N D\^{}2)}
\NormalTok{eigendecomposition:      O(D\^{}3)}
\NormalTok{stored operator:         O(D\^{}2)}
\NormalTok{base transformation:     O(N D\^{}2)}
\NormalTok{raw{-}query transformation: O(D\^{}2) per query}
\end{Highlighting}
\end{Shaded}

Residual PCA additionally constructs residuals in \texttt{O(ND)}, then
has the same dominant covariance/eigendecomposition scale.

Concrete CIFAR operator state:

\begin{Shaded}
\begin{Highlighting}[]
\NormalTok{(D\^{}2 + D) float32 values}
\NormalTok{= (512\^{}2 + 512) * 4}
\NormalTok{= 1,050,624 bytes}
\end{Highlighting}
\end{Shaded}

Concrete GIST dense state:

\begin{Shaded}
\begin{Highlighting}[]
\NormalTok{(960\^{}2 + 960) * 4 = 3,690,240 bytes}
\end{Highlighting}
\end{Shaded}

Current and residual PCA have the same operator shape, so operator size
does not explain their paired result. However, the fixed-candidate study
does not support an end-to-end QPS claim because raw-query
transformation was not timed inside a complete

search path.
\end{frame}

\begin{frame}[fragile,allowframebreaks]{16. Attempt 1A Decision}
\protect\phantomsection\label{16-attempt-1a-decision}
\begin{Shaded}
\begin{Highlighting}[]
\NormalTok{Decision: CLOSE FULL{-}D PCA REPLACEMENT AS A MAIN DIRECTION}
\end{Highlighting}
\end{Shaded}

Why:

\begin{enumerate}
\tightlist
\item
  GIST shows only small estimator RMSE changes.
\item
  Fast prefixes worsen.
\item
  Ranking gains are not statistically stable.
\item
  Standard residual PCA explains the only stable GIST signal.
\item
  The preregistered CIFAR replication fails confidence, seed,
  rotation-off, and no-harm requirements.
\end{enumerate}

What was deliberately not implemented:

\begin{Shaded}
\begin{Highlighting}[]
\NormalTok{SAQ{-}aware learned rotation}
\NormalTok{joint transform{-}plan optimization}
\NormalTok{OPQ/ITQ{-}style learner integration}
\NormalTok{broad transform/dataset sweep}
\NormalTok{new persisted transform format}
\end{Highlighting}
\end{Shaded}

Claim boundary:

\begin{Shaded}
\begin{Highlighting}[]
\NormalTok{This does not prove PCA universally optimal. It says the measured premise is}
\NormalTok{not strong enough to justify learned replacement development.}
\end{Highlighting}
\end{Shaded}
\end{frame}

\begin{frame}[fragile,allowframebreaks]{17. Attempt 1B: Why Test Lossy
\texttt{D\ -\textgreater{}\ d} Projection?}
\protect\phantomsection\label{17-attempt-1b-why-test-lossy-d---d-projection}
Full-D transforms preserve dimension and therefore cannot directly
reduce code work or dense query-state size.

More aggressive idea:

\begin{Shaded}
\begin{Highlighting}[]
\NormalTok{retain only the high{-}variance PCA head}
\NormalTok{summarize the omitted tail with a small statistic}
\NormalTok{run quantization/search in d \textless{} D dimensions}
\end{Highlighting}
\end{Shaded}

Potential benefit:

\begin{Shaded}
\begin{Highlighting}[]
\NormalTok{fewer dimensions transformed}
\NormalTok{fewer dimensions encoded}
\NormalTok{fewer code/factor bytes}
\NormalTok{less estimator work}
\end{Highlighting}
\end{Shaded}

Risk:

\begin{Shaded}
\begin{Highlighting}[]
\NormalTok{physical projection changes original{-}space distances and nearest neighbors}
\end{Highlighting}
\end{Shaded}

Related-work boundary:

\begin{Shaded}
\begin{Highlighting}[]
\NormalTok{truncated PCA, LeanVec{-}style projection, and MRQ{-}style head/tail refinement}
\NormalTok{already occupy the broad design space}
\end{Highlighting}
\end{Shaded}

Therefore the first test is an intentionally favorable quality upper
bound, not a projected index implementation.
\end{frame}

\begin{frame}[allowframebreaks]{18. Lossy Projection Distance
Decomposition}
\protect\phantomsection\label{18-lossy-projection-distance-decomposition}
Split PCA coordinates into retained head and omitted tail:

\[q=(q_h,q_t),\qquad x=(x_h,x_t).\]

Original exact squared L2:

\[D_0=
\lVert q_h-x_h\rVert^2+
\lVert q_t\rVert^2+
\lVert x_t\rVert^2-
2\langle q_t,x_t\rangle.\]

Pure-head projection:

\[D_{head}=\lVert q_h-x_h\rVert^2.\]

Exact tail-norm surrogate:

\[D_{head+norm}=
\lVert q_h-x_h\rVert^2+
\lVert q_t\rVert^2+
\lVert x_t\rVert^2.\]

The missing information is exactly:

\[-2\langle q_t,x_t\rangle.\]

The oracle uses exact float64 head distances and exact tail norms. It
adds no head quantization error, making it strictly more favorable than
a deployed projected SAQ implementation.
\end{frame}

\begin{frame}[fragile,allowframebreaks]{19. Lossy Projection Running
Example}
\protect\phantomsection\label{19-lossy-projection-running-example}
Suppose two candidates have identical head distance and equal tail norm:

\begin{Shaded}
\begin{Highlighting}[]
\NormalTok{head distance:       1.00 for both}
\NormalTok{query tail norm\^{}2:   0.50}
\NormalTok{base tail norm\^{}2:    0.50 for both}
\end{Highlighting}
\end{Shaded}

But their tail correlations differ:

\begin{Shaded}
\begin{Highlighting}[]
\NormalTok{candidate p: \textless{}q\_t, p\_t\textgreater{} = 0.45}
\NormalTok{candidate n: \textless{}q\_t, n\_t\textgreater{} = 0.35}
\end{Highlighting}
\end{Shaded}

Exact tail contributions:

\begin{Shaded}
\begin{Highlighting}[]
\NormalTok{p: 0.50 + 0.50 {-} 2*0.45 = 0.10}
\NormalTok{n: 0.50 + 0.50 {-} 2*0.35 = 0.30}
\end{Highlighting}
\end{Shaded}

Exact distances:

\begin{Shaded}
\begin{Highlighting}[]
\NormalTok{p = 1.10}
\NormalTok{n = 1.30}
\end{Highlighting}
\end{Shaded}

Norm-only surrogate:

\begin{Shaded}
\begin{Highlighting}[]
\NormalTok{p = 2.00}
\NormalTok{n = 2.00}
\end{Highlighting}
\end{Shaded}

Tail norms remove average bias but cannot distinguish candidates whose
omitted tail inner products carry ranking information.
\end{frame}

\begin{frame}[fragile,allowframebreaks]{20. LP-0 Gate A Protocol}
\protect\phantomsection\label{20-lp-0-gate-a-protocol}
\begin{itemize}
\tightlist
\item
  \textbf{Item}: Dataset

  \textbf{Frozen value}: GIST sample50k
\item
  \textbf{Item}: Original dimension

  \textbf{Frozen value}: \texttt{D=960}
\item
  \textbf{Item}: Retained head

  \textbf{Frozen value}: \texttt{d=576}, exactly at an SAQ segment
  boundary
\item
  \textbf{Item}: Physical reduction

  \textbf{Frozen value}: 40\% of dimensions removed
\item
  \textbf{Item}: IVF

  \textbf{Frozen value}: \texttt{K\_IVF=512}
\item
  \textbf{Item}: Budget of native comparator

  \textbf{Frozen value}: \texttt{B=4}
\item
  \textbf{Item}: Queries

  \textbf{Frozen value}: first 128 held-out queries
\item
  \textbf{Item}: Fixed probes

  \textbf{Frozen value}: 16/query
\item
  \textbf{Item}: Fixed candidates

  \textbf{Frozen value}: 442,823
\item
  \textbf{Item}: Ranking cutoff

  \textbf{Frozen value}: top 100
\item
  \textbf{Item}: Oracle head

  \textbf{Frozen value}: exact float64 distance over first 576 PCA
  coordinates
\item
  \textbf{Item}: Tail

  \textbf{Frozen value}: exact float64 norms; float32 deployed summary
  also measured
\item
  \textbf{Item}: Comparator

  \textbf{Frozen value}: seed-averaged native full-D SAQ full-code
  estimator
\end{itemize}

Native GIST plan:

\begin{Shaded}
\begin{Highlighting}[]
\NormalTok{64@11 | 192@6 | 320@4 | 256@2 | 128@0}
\NormalTok{                       \^{}}
\NormalTok{                    d = 576}
\end{Highlighting}
\end{Shaded}

Projection removes both the \texttt{256@2} segment and the
\texttt{128@0} tail physically.
\end{frame}

\begin{frame}[fragile,allowframebreaks]{21. LP-0 Ranking Result}
\protect\phantomsection\label{21-lp-0-ranking-result}
\begin{itemize}
\tightlist
\item
  \textbf{Estimator}: exact projected head + exact tail norms

  \textbf{Top-100 agreement}: 0.992890625
\item
  \textbf{Boundary inversion rate}: \texttt{6.85248e-6}
\item
  \textbf{Estimator}: native full-D SAQ full code

  \textbf{Top-100 agreement}: 0.994617188
\item
  \textbf{Boundary inversion rate}: \texttt{4.08687e-6}
\end{itemize}

Registered deltas, projected oracle minus native SAQ:

\begin{Shaded}
\begin{Highlighting}[]
\NormalTok{top{-}100 agreement delta       = {-}0.0017265625}
\NormalTok{one{-}sided 95\% lower bound     = {-}0.002578125}

\NormalTok{boundary inversion delta      = +2.76561e{-}6}
\NormalTok{one{-}sided 95\% upper bound      = +4.21688e{-}6}
\end{Highlighting}
\end{Shaded}

Both ranking gates fail. Rotation-off has the same unfavorable
directions:

\begin{Shaded}
\begin{Highlighting}[]
\NormalTok{agreement delta  = {-}0.00140625}
\NormalTok{inversion delta  = +2.70090e{-}6}
\end{Highlighting}
\end{Shaded}

This is not a quantizer implementation failure: the favorable exact
surrogate itself is below native SAQ ranking quality.
\end{frame}

\begin{frame}[fragile,allowframebreaks]{22. LP-0 Error Attribution}
\protect\phantomsection\label{22-lp-0-error-attribution}
Pooled absolute-error summary:

\begin{itemize}
\tightlist
\item
  \textbf{Component}: pure head, \texttt{D\_head-D0}

  \textbf{Bias}: -0.0339155
\item
  \textbf{MAE}: 0.0339155

  \textbf{Pooled RMSE}: 0.0444650
\item
  \textbf{Component}: exact norm surrogate, \texttt{D\_head+norm-D0}

  \textbf{Bias}: \texttt{-3.25e-6}
\item
  \textbf{MAE}: 0.00150736

  \textbf{Pooled RMSE}: 0.00248180
\item
  \textbf{Component}: float32 tail-summary precision

  \textbf{Bias}: \texttt{5.37e-11}
\item
  \textbf{MAE}: \texttt{1.05e-7}

  \textbf{Pooled RMSE}: \texttt{1.51e-7}
\item
  \textbf{Component}: deployed norm surrogate total

  \textbf{Bias}: \texttt{-3.25e-6}
\item
  \textbf{MAE}: 0.00150736

  \textbf{Pooled RMSE}: 0.00248180
\end{itemize}

Interpretation:

\begin{Shaded}
\begin{Highlighting}[]
\NormalTok{tail norms remove nearly all systematic bias}
\NormalTok{float32 summary precision is negligible}
\NormalTok{remaining error comes from omitted tail inner{-}product variation}
\end{Highlighting}
\end{Shaded}

The problem is not storing the tail norm inaccurately. The problem is
that one norm per vector cannot recover query-candidate tail correlation
near the top-k boundary.
\end{frame}

\begin{frame}[fragile,allowframebreaks]{23. Lossy Projection Complexity
And Gate Logic}
\protect\phantomsection\label{23-lossy-projection-complexity-and-gate-logic}
An actual dense \texttt{D\ -\textgreater{}\ d} projection would require:

\begin{Shaded}
\begin{Highlighting}[]
\NormalTok{stored operator: O(Dd)}
\NormalTok{base transformation: O(NDd)}
\NormalTok{raw{-}query transformation: O(Dd)}
\NormalTok{head quantization/search: method{-}dependent O(d)}
\NormalTok{tail metadata: at least O(1) per vector for norm{-}only design}
\end{Highlighting}
\end{Shaded}

But Gate A deliberately removes projected-head quantization and runtime
from the quality question:

\begin{Shaded}
\begin{Highlighting}[]
\NormalTok{exact head + exact tail norms already loses to native SAQ}
\end{Highlighting}
\end{Shaded}

Therefore adding:

\begin{Shaded}
\begin{Highlighting}[]
\NormalTok{projected{-}head CAQ error}
\NormalTok{progressive{-}estimator error}
\NormalTok{projection/index metadata}
\NormalTok{query transform and preparation work}
\end{Highlighting}
\end{Shaded}

cannot rescue the registered quality premise without introducing a
different tail mechanism. Gate B, projected index construction and
Recall-QPS evaluation, was not authorized.
\end{frame}

\begin{frame}[fragile,allowframebreaks]{24. Attempt 1B Decision}
\protect\phantomsection\label{24-attempt-1b-decision}
\begin{Shaded}
\begin{Highlighting}[]
\NormalTok{Decision: FAIL GATE A AND STOP THE REGISTERED D=576 LINE}
\end{Highlighting}
\end{Shaded}

Why:

\begin{enumerate}
\tightlist
\item
  The exact head plus exact tail-norm oracle ranks candidates worse than
  native full-D SAQ.
\item
  The missing tail inner product, not summary precision, dominates the
  residual error.
\item
  A real projected quantizer would add more error and overhead.
\item
  Sweeping \texttt{d}, \texttt{B}, probes, or tail rules after this
  registered failure would be post-hoc rescue.
\end{enumerate}

What the result does not prove:

\begin{Shaded}
\begin{Highlighting}[]
\NormalTok{all lossy projection is ineffective}
\NormalTok{all dimensions d fail}
\NormalTok{all richer tail representations fail}
\NormalTok{PCA truncation is universally inferior}
\end{Highlighting}
\end{Shaded}

A future projection study would require a distinct related-work-grounded
mechanism and a new protocol, not continuation of this line.
\end{frame}

\begin{frame}[fragile,allowframebreaks]{25. Attempt 1 Unified
Conclusion}
\protect\phantomsection\label{25-attempt-1-unified-conclusion}
\begin{itemize}
\tightlist
\item
  \textbf{Sub-attempt}: 1A full-D residual PCA

  \textbf{Positive signal}: GIST accurate/full RMSE improves
  0.60\%/0.33\%
\item
  \textbf{Failed requirement}: no stable ranking gain; fast harm; CIFAR
  replication fails

  \textbf{Final status}: Closed
\item
  \textbf{Sub-attempt}: 1B lossy \texttt{960\ -\textgreater{}\ 576}

  \textbf{Positive signal}: exact tail norms reduce projection RMSE from
  0.0445 to 0.00248
\item
  \textbf{Failed requirement}: exact favorable surrogate still ranks
  worse than native SAQ

  \textbf{Final status}: Closed
\end{itemize}

Unified research conclusion:

\begin{Shaded}
\begin{Highlighting}[]
\NormalTok{PCA is not proven universally optimal, but the measured evidence does not}
\NormalTok{support spending additional degrees of freedom on either a learned full{-}D}
\NormalTok{replacement or the registered norm{-}only lossy projection.}
\end{Highlighting}
\end{Shaded}

Strict-reviewer view:

\begin{Shaded}
\begin{Highlighting}[]
\NormalTok{The full{-}D signal is small and non{-}replicated.}
\NormalTok{The lossy signal fails before quantization is introduced.}
\NormalTok{Neither supports a new transform method or a system{-}level Pareto claim.}
\end{Highlighting}
\end{Shaded}

Speaker notes:

\begin{itemize}
\tightlist
\item
  Attempt 1 is a useful example of staged falsification.
\item
  The more aggressive variant was tested with a more favorable oracle,
  not a weaker implementation.
\end{itemize}
\end{frame}

\begin{frame}[fragile,allowframebreaks]{26. Attempt 2: Research
Question}
\protect\phantomsection\label{26-attempt-2-research-question}
Attempt 2 comes from the sibling repository:

\begin{Shaded}
\begin{Highlighting}[]
\NormalTok{/rwproject/kdd{-}db/kluaq/vectordb}
\NormalTok{main checkpoint: f51b487}
\NormalTok{exact{-}hist implementation checkpoint: 9a7026d}
\NormalTok{cross{-}dataset comparison checkpoint:    870d829}
\end{Highlighting}
\end{Shaded}

Original empirical question:

\begin{Shaded}
\begin{Highlighting}[]
\NormalTok{If each coordinate uses a scalar codebook, can exact 1D dynamic programming}
\NormalTok{produce a better codebook than Lloyd refinement, and does lower raw}
\NormalTok{reconstruction SSE improve ANN recall?}
\end{Highlighting}
\end{Shaded}

The 2026-07-13 related-work audit adds a distinct novelty question:

\begin{Shaded}
\begin{Highlighting}[]
\NormalTok{Does White and Singal\textquotesingle{}s 2026 inner{-}product{-}aware quantization paper, together}
\NormalTok{with the exact 1D k{-}means literature it uses, already subsume Attempt 2?}
\end{Highlighting}
\end{Shaded}

Short answer:

\begin{Shaded}
\begin{Highlighting}[]
\NormalTok{public artifacts: no integrated two{-}DP pipeline is documented or exposed}
\NormalTok{combination claim: no novelty is established from that absence}
\NormalTok{inner exact scalar{-}codebook DP: already established; no novelty claim remains}
\NormalTok{local value that remains: controlled baseline and objective{-}mismatch evidence}
\end{Highlighting}
\end{Shaded}

The paper was submitted on 2026-05-29, before the \texttt{vectordb}
exact-hist implementation on 2026-06-26. More importantly, exact 1D
quantization DP predates both by decades.

Speaker notes:

\begin{itemize}
\tightlist
\item
  Attempt 2 is not SAQ\textquotesingle s segment-planner DP.
\item
  It is also not the CAQ exact direction-code oracle from the later
  branch.
\item
  "No integrated public pipeline" must not be misread as "novel as a
  component combination."
\end{itemize}
\end{frame}

\begin{frame}[fragile,allowframebreaks]{27. Two Different Dynamic
Programs}
\protect\phantomsection\label{27-two-different-dynamic-programs}
The \texttt{vectordb} pipeline contains two DPs.

\textbf{Inner DP: scalar codebook training}

For one dimension \texttt{j} and bitwidth \texttt{b}:

\begin{Shaded}
\begin{Highlighting}[]
\NormalTok{input:  one sorted raw scalar column; the trainer internally constructs its}
\NormalTok{        equal{-}count weighted histogram}
\NormalTok{output: K\_b centroids minimizing whole{-}bin partition SSE, with cuts restricted}
\NormalTok{        to histogram{-}bin boundaries}
\end{Highlighting}
\end{Shaded}

Only when each raw scalar is its own bin does this equal globally
minimal raw-sample 1D reconstruction SSE. After DP backtracking, raw
nearest-centroid SSE is separately recomputed under the deployed
encoding semantics.

\textbf{Outer DP: bit allocation}

Given one error value \texttt{E{[}j,b{]}} per dimension and bitwidth:

\[A[j,c]=\min_b\left(A[j-1,c-b]+w_jE[j,b]\right),\]

subject to:

\[\sum_j b_j = D B.\]

The outer DP chooses the bitwidth of every dimension under an exact
total bit budget. Attempt 2 changes only the inner trainer from Lloyd to
exact-hist; the outer algorithm, objective form, and weights are held
fixed as part

of the evaluation pipeline. The selected bit vector is not fixed: it
changes when the new trainer changes \texttt{E{[}j,b{]}}.

Novelty accounting must therefore be layer-specific:

\begin{Shaded}
\begin{Highlighting}[]
\NormalTok{inner DP: exact 1D scalar clustering / quantization is prior work}
\NormalTok{outer DP: classical discrete rate allocation / multiple{-}choice knapsack}
\NormalTok{local experiment: asks whether the replacement changes ANN recall}
\end{Highlighting}
\end{Shaded}
\end{frame}

\begin{frame}[fragile,allowframebreaks]{28. Dimensionwise
Scalar-Quantization Setting}
\protect\phantomsection\label{28-dimensionwise-scalar-quantization-setting}
Write the base matrix as:

\[X\in\mathbb{R}^{N\times D}.\]

Attempt 2 trains on one column at a time:

\[X[:,j]=(x_{1j},\ldots,x_{Nj}),\]

and shares the resulting codebook across all \texttt{N} base vectors.
For database vector:

\[x=(x_1,\ldots,x_D).\]

Each dimension has an independent codebook:

\[C_{j,b}=\{c_{j,b,1},\ldots,c_{j,b,K_b}\}.\]

Base encoding:

\[\hat x_j=\operatorname*{argmin}_{c\in C_{j,b_j}}(x_j-c)^2.\]

The query remains float. Compressed squared-L2 scan uses:

\[\hat d(q,x)=\sum_j(q_j-\hat x_j)^2.\]

Candidate bitwidths in variable mode:

\begin{Shaded}
\begin{Highlighting}[]
\NormalTok{b in \{0,1,2,3,4,5,6,7,8\}}
\NormalTok{total budget = D * B bits/vector}
\end{Highlighting}
\end{Shaded}

Historical implementation boundary:

\begin{Shaded}
\begin{Highlighting}[]
\NormalTok{unpacked uint8 codes}
\NormalTok{brute{-}force compressed scan}
\NormalTok{L2 only}
\end{Highlighting}
\end{Shaded}

This is an algorithmic diagnostic, not a QPS-comparable SAQ index.

The column-wise training direction matters for the paper comparison. The
paper\textquotesingle s formal ASQ problem instead chooses a
quantization set from the coordinates of one input vector; its official
code also contains a separate shared-block helper, audited later.
\end{frame}

\begin{frame}[fragile,allowframebreaks]{29. Exact-Hist Construction}
\protect\phantomsection\label{29-exact-hist-construction}
For one sorted dimension with \texttt{N} values:

\begin{Shaded}
\begin{Highlighting}[]
\NormalTok{H = min(max\_histogram\_bins, N)}
\NormalTok{default H = 256}
\end{Highlighting}
\end{Shaded}

Step 1: split sorted values into \texttt{H} equal-count bins.

For each bin, record:

\begin{Shaded}
\begin{Highlighting}[]
\NormalTok{weight}
\NormalTok{sum}
\NormalTok{sum\_sq}
\end{Highlighting}
\end{Shaded}

and their prefix sums.

Step 2: any contiguous bin interval \texttt{{[}a,b)} has one-centroid
optimum:

\[\mu(a,b)=\frac{sum(a,b)}{weight(a,b)},\]

\[SSE(a,b)=sum\_sq(a,b)-\frac{sum(a,b)^2}{weight(a,b)}.\]

Step 3: partition the \texttt{H} ordered bins into \texttt{K\_b}
contiguous clusters.

Step 4: backtrack split points and use each interval mean as its
centroid.

Implementation detail that affects the guarantee:

\begin{Shaded}
\begin{Highlighting}[]
\NormalTok{The histogram DP value is not inserted directly into E[j,b].}
\NormalTok{After recovering centroids, midpoint Voronoi boundaries are rebuilt, every raw}
\NormalTok{value is encoded by nearest centroid, and the raw{-}sample SSE of those encodings}
\NormalTok{is recomputed and passed to the outer DP.}
\end{Highlighting}
\end{Shaded}

Thus the reported training SSE measures the realized raw encoding, while
the global optimality claim applies only to split points between
indivisible bins. The recomputation evaluates the selected centroid set;
it does not reoptimize that set over raw-sample

split points.
\end{frame}

\begin{frame}[fragile,allowframebreaks]{30. Scalar DP Recurrence And
Guarantee}
\protect\phantomsection\label{30-scalar-dp-recurrence-and-guarantee}
Define:

\begin{Shaded}
\begin{Highlighting}[]
\NormalTok{F[k,r] = minimum weighted SSE for the first r bins using k centroids}
\end{Highlighting}
\end{Shaded}

Recurrence:

\[F[k,r]=\min_{m<r}\left(F[k-1,m]+SSE(m,r)\right).\]

The implementation uses monotone split points and divide-and-conquer DP
optimization to compute each layer.

Exact guarantee:

\begin{Shaded}
\begin{Highlighting}[]
\NormalTok{For fixed H equal{-}count bins and K\_b centroids, the DP returns the minimum}
\NormalTok{raw{-}value SSE among contiguous partitions whose split points lie only at bin}
\NormalTok{boundaries. Prefix weight/sum/sum\_sq statistics evaluate each allowed}
\NormalTok{interval\textquotesingle{}s raw SSE exactly; approximation enters through forbidden in{-}bin}
\NormalTok{split points.}
\end{Highlighting}
\end{Shaded}

It does not guarantee:

\begin{Shaded}
\begin{Highlighting}[]
\NormalTok{minimum SSE over all N raw values when H \textless{} N}
\NormalTok{minimum query{-}distance error}
\NormalTok{minimum top{-}k boundary inversion}
\NormalTok{maximum recall}
\end{Highlighting}
\end{Shaded}

If \texttt{H=N}, every bin contains one sample and the formulation
becomes full-sample exact 1D k-means. That expensive configuration was
not run in the recorded experiments.

Strict novelty implication:

\begin{Shaded}
\begin{Highlighting}[]
\NormalTok{The missing full{-}scale H=N run is a limitation of this local baseline,}
\NormalTok{not an open algorithmic problem in the literature.}
\end{Highlighting}
\end{Shaded}

The recent paper\textquotesingle s official repository directly vendors
and exposes a raw exact 1D k-means solver primitive from earlier work;
no public experiment caller integrates or evaluates that primitive.
Scaling the local code by removing histogram coarsening would therefore

close an implementation gap, not create a method contribution.
\end{frame}

\begin{frame}[fragile,allowframebreaks]{31. White And Singal 2026:
Formal Problem}
\protect\phantomsection\label{31-white-and-singal-2026-formal-problem}
\href{https://arxiv.org/html/2606.00289v1}{Inner Product Aware
Quantization: Provably Fast, Accurate, and Adaptive Algorithms} studies
a different formal quantization problem.

This is a May 2026 arXiv v1 preprint by Nathan White and Krish Singal;
this deck makes no venue-acceptance claim. Its contributions are the
MDV/ADV inner-product objectives, rounding-distribution results,
exact/approximate algorithms, and practical ASQ acceleration, not the
invention

of exact 1D k-means.

For one vector \texttt{w} and a quantization-set budget
\texttt{\textbar{}Q(w)\textbar{}\ \textless{}=\ s}, standard stochastic
quantization independently rounds each coordinate to its adjacent lower
or upper point in \texttt{Q(w)}, using the unique probabilities that
make the rounding unbiased. Several algorithms

then return a set of size \texttt{s}.

For an unseen input/query distribution, written as
\texttt{\textbackslash{}mathcal\ X} below, its Average Directional
Variance is:

\[\operatorname{ADV}_{\mathcal X}(w,Q)
=\sum_i \lambda_i
(w_i^\uparrow-w_i)(w_i-w_i^\downarrow),
\qquad
\lambda_i=\mathbb E_{x\sim\mathcal X}[x_i^2].\]

The paper also studies Maximum Directional Variance, which controls the
worst coordinate variance rather than their weighted sum.

Key semantic boundary:

\begin{Shaded}
\begin{Highlighting}[]
\NormalTok{paper:     adaptive, unbiased, adjacent stochastic rounding}
\NormalTok{Attempt 2: shared, biased, deterministic nearest{-}centroid encoding}
\end{Highlighting}
\end{Shaded}
\end{frame}

\begin{frame}[fragile,allowframebreaks]{32. Related DPs, Different
Interval Costs}
\protect\phantomsection\label{32-related-dps-different-interval-costs}
After jointly sorting coordinate-weight pairs so that
\texttt{w\_1\ \textless{}=\ ...\ \textless{}=\ w\_d} while each
\texttt{lambda\_i} remains attached to its original coordinate, define
the paper\textquotesingle s ADV interval cost:

\[C[j,k]=\sum_{i=j}^{k}\lambda_i
(w_k-w_i)(w_i-w_j).\]

A conventional equivalent form of its quantization-set DP is:

\[G[t,k]=\min_{j<k}\left(G[t-1,j]+C[j,k]\right).\]

Here \texttt{G{[}t,k{]}} is the best cost for \texttt{t} selected
quantization points ending at \texttt{w\_k}, with \texttt{G{[}1,1{]}=0};
the answer is \texttt{G{[}s,d{]}}.

Attempt 2 instead uses a best-mean interval cost:

\[F[t,r]=\min_{m<r}
\left(F[t-1,m]+SSE_{\mathrm{mean}}(m,r)\right).\]

The relationship is structural, not semantic:

\begin{itemize}
\tightlist
\item
  \textbf{Property}: Ordered objects

  \textbf{Paper ADV/ASQ}: coordinates of one vector
\item
  \textbf{Attempt 2 inner DP}: weighted bins from one database column
\item
  \textbf{Property}: Interval representative

  \textbf{Paper ADV/ASQ}: two stochastic endpoints
\item
  \textbf{Attempt 2 inner DP}: one deterministic mean
\item
  \textbf{Property}: Interval cost

  \textbf{Paper ADV/ASQ}: weighted rounding variance
\item
  \textbf{Attempt 2 inner DP}: centroid reconstruction SSE
\item
  \textbf{Property}: Shared skeleton

  \textbf{Paper ADV/ASQ}: sorted contiguous partition with Monge
  structure
\item
  \textbf{Attempt 2 inner DP}: sorted contiguous partition with monotone
  splits
\end{itemize}

The paper explicitly notes that if biased quantization is allowed, the
optimal scheme is 1D k-means plus nearest-cluster assignment. That
observation places Attempt 2\textquotesingle s mathematical core inside
established related work.
\end{frame}

\begin{frame}[fragile,allowframebreaks]{33. The Training Axis Is
Transposed}
\protect\phantomsection\label{33-the-training-axis-is-transposed}
Let \texttt{X} be an \texttt{N\ x\ D} database matrix.

\begin{itemize}
\tightlist
\item
  \textbf{Question}: Optimization input

  \textbf{Attempt 2}: column \texttt{X{[}:,j{]}}, containing \texttt{N}
  database scalars
\item
  \textbf{Paper formal model}: row/vector \texttt{X{[}i,:{]}},
  containing \texttt{D} coordinates
\item
  \textbf{Question}: Learned set

  \textbf{Attempt 2}: \texttt{C{[}j,b{]}} for one dimension and bitwidth
\item
  \textbf{Paper formal model}: \texttt{Q(X{[}i,:{]})} for one vector
\item
  \textbf{Question}: Sharing

  \textbf{Attempt 2}: all \texttt{N} database vectors share
  \texttt{C{[}j,b{]}}
\item
  \textbf{Paper formal model}: different vectors may have different
  \texttt{Q}
\item
  \textbf{Question}: Rate choice

  \textbf{Attempt 2}: heterogeneous \texttt{b\_j}, exact total budget
\item
  \textbf{Paper formal model}: one cardinality budget \texttt{s}, with `
\item
  \textbf{Question}: Encoding/objective semantics

  \textbf{Attempt 2}: biased deterministic nearest-centroid SSE
\item
  \textbf{Paper formal model}: unbiased stochastic rounding optimized
  for IP variance
\end{itemize}

The paper\textquotesingle s formal vector-search limitation follows from
this orientation: it requires per-vector \texttt{Q} storage. Its
discussion therefore leaves adaptation to potentially dynamic
multi-vector quantization as
\href{https://arxiv.org/html/2606.00289v1\#S5}{future work}.

This formal difference is real, but it is not enough to establish
Attempt 2 novelty: dimensionwise shared scalar codebooks and discrete
rate allocation are themselves classical constructions.
\end{frame}

\begin{frame}[fragile,allowframebreaks]{34. Official-Code Audit:
Stronger Adjacency}
\protect\phantomsection\label{34-official-code-audit-stronger-adjacency}
The pinned
\href{https://github.com/nathanllww/Inner-Product-Aware-Quantization/tree/e92ed904b85ad12469a0a350c81071ff9ef6aa37}{official
repository} contains more than the paper\textquotesingle s formal
per-vector interface:

\begin{Shaded}
\begin{Highlighting}[]
\NormalTok{batch\_row\_quant(X, ...)}
\NormalTok{  one Q per row/vector; matches the formal adaptive direction}

\NormalTok{block\_quant(X, m, s, ...)}
\NormalTok{  flatten all N*m values in each consecutive m{-}column block}
\NormalTok{  and train one shared size{-}s Q for that block}
\end{Highlighting}
\end{Shaded}

\begin{Shaded}
\begin{Highlighting}[]
\NormalTok{kmeans\_wilber\_1d(sorted\_points, k)}
\NormalTok{  raw{-}point globally optimal 1D k{-}means implementation}
\NormalTok{  requires ascending input; the wrapper does not sort}
\end{Highlighting}
\end{Shaded}

Pinned source:
\href{https://github.com/nathanllww/Inner-Product-Aware-Quantization/blob/e92ed904b85ad12469a0a350c81071ff9ef6aa37/cython/pq.pyx\#L16-L202}{shared-block
and per-row interfaces},
\href{https://github.com/nathanllww/Inner-Product-Aware-Quantization/blob/e92ed904b85ad12469a0a350c81071ff9ef6aa37/cython/kmeans1d.pyx\#L38-L85}{raw
exact 1D k-means wrapper}.

Consequently, \texttt{block\_quant(...,\ m=1,\ ...)} has exactly the
same training axis as a fixed-rate shared per-dimension codebook.
However:

\begin{Shaded}
\begin{Highlighting}[]
\NormalTok{it calls vmix\_approx for ADV rather than centroid{-}SSE k{-}means}
\NormalTok{all blocks receive the same s}
\NormalTok{there is no per{-}dimension error table or outer bit{-}allocation DP}
\end{Highlighting}
\end{Shaded}

The \texttt{block\_quant} docstring says \texttt{exact\_adv}, while the
pinned implementation calls \texttt{vmix\_approx}; the implementation is
the evidence used here. The public repository contains no caller or
experiment script that proves which path generated the
paper\textquotesingle s GloVe table,

so this deck makes no such attribution.

The paper calls its GloVe result preliminary evidence and reports
Recall@100 for L2 and maximum-inner-product search at an average four
bits per coordinate. It does not report a controlled Lloyd-versus-exact
trainer ablation, QPS/build cost, or complete quantization-set storage

accounting. That table therefore does not answer the specific Attempt 2
end-to-end question.
\end{frame}

\begin{frame}[fragile,allowframebreaks]{35. Prior-Art And Layer-By-Layer
Verdict}
\protect\phantomsection\label{35-prior-art-and-layer-by-layer-verdict}
\begin{itemize}
\tightlist
\item
  \textbf{Attempt 2 layer}: Raw exact 1D k-means optimizer

  \textbf{Formal paper}: Recognizes it as older related work
\item
  \textbf{Pinned public code}: Exposes \texttt{kmeans\_wilber\_1d} from
  the reused exact-1D library

  \textbf{Strict-reviewer verdict}: Novelty foreclosed by Wu 1991 / GLM
  2017; not a White-Singal contribution
\item
  \textbf{Attempt 2 layer}: Shared per-dimension fixed-rate set

  \textbf{Formal paper}: Formal unit is one vector
\item
  \textbf{Pinned public code}: \texttt{block\_quant(m=1)} is directly
  adjacent but uses ASQ semantics

  \textbf{Strict-reviewer verdict}: Not a formal guarantee; shared
  scalar codebooks are not enough for novelty
\item
  \textbf{Attempt 2 layer}: Potential diagonal query-second-moment
  weighting

  \textbf{Formal paper}: ADV supplies diagonal second moments under
  per-vector independent unbiased rounding
\item
  \textbf{Pinned public code}: \texttt{block\_quant} uses all-one
  weights; no such experiment is exposed

  \textbf{Strict-reviewer verdict}: High objective-level adjacency, not
  direct pipeline coverage or a validity proof for biased SQ
\item
  \textbf{Attempt 2 layer}: Multi-bit error table \texttt{E{[}j,b{]}}

  \textbf{Formal paper}: Fixes one cardinality budget \texttt{s}; no
  table
\item
  \textbf{Pinned public code}: Repeated fixed-\texttt{s} calls are an
  API inference, not integrated or evaluated

  \textbf{Strict-reviewer verdict}: Not claimed novel here; separate
  rate-distortion prior support is required
\item
  \textbf{Attempt 2 layer}: Outer heterogeneous-bit DP

  \textbf{Formal paper}: Not presented
\item
  \textbf{Pinned public code}: Not exposed

  \textbf{Strict-reviewer verdict}: Not covered by this paper, but
  \href{https://arxiv.org/abs/2509.12086}{SAQ} already provides DP
  segmentation/bit allocation
\item
  \textbf{Attempt 2 layer}: Controlled Lloyd/exact-hist L2 recall audit

  \textbf{Formal paper}: Only preliminary GloVe L2/MIPS comparison
  against PQ
\item
  \textbf{Pinned public code}: No public experiment caller/configuration

  \textbf{Strict-reviewer verdict}: Useful local negative/baseline
  evidence, not a quantizer contribution
\item
  \textbf{Attempt 2 layer}: Packed variable-rate ANN system

  \textbf{Formal paper}: Per-vector storage is listed as a
  limitation/future-work issue
\item
  \textbf{Pinned public code}: Fixed-rate block helper and 256-code
  search helper only

  \textbf{Strict-reviewer verdict}: A possible problem layer, not an
  achieved result in either artifact
\end{itemize}

Timeline:

\begin{Shaded}
\begin{Highlighting}[]
\NormalTok{paper and official code:       2026{-}05{-}29}
\NormalTok{vectordb exact{-}hist checkpoint: 2026{-}06{-}26}
\end{Highlighting}
\end{Shaded}

The stronger point is not temporal priority. The paper itself traces
exact 1D k-means matrix-search DP to Wu 1991. Separately, the pinned
code exposes a solver primitive from the reused
\href{https://arxiv.org/abs/1701.07204}{2017 exact-1D-clustering
implementation}, without an integrated public experiment path.

Bottom line:

\begin{Shaded}
\begin{Highlighting}[]
\NormalTok{public{-}artifact level: no integrated shared{-}codebook + heterogeneous{-}bit +}
\NormalTok{                       outer{-}allocation pipeline is documented or exposed}
\NormalTok{novelty level:         absence is not combination novelty; the inner optimizer}
\NormalTok{                       is foreclosed by older prior art, and remaining system}
\NormalTok{                       layers were not built or validated by Attempt 2}
\end{Highlighting}
\end{Shaded}
\end{frame}

\begin{frame}[fragile,allowframebreaks]{36. Attempt 2 Code Map}
\protect\phantomsection\label{36-attempt-2-code-map}
Sibling repository code:

\begin{Shaded}
\begin{Highlighting}[]
\NormalTok{src/scalar\_quantizer.cpp}
\NormalTok{  BuildEqualCountHistogram(...)}
\NormalTok{  HistogramIntervalSse(...)}
\NormalTok{  ComputeExactKMeansLayer(...)}
\NormalTok{  TrainExactScalarQuantizerFromSorted(...)}
\NormalTok{  midpoint boundaries, raw nearest{-}centroid reassignment, raw SSE}
\end{Highlighting}
\end{Shaded}

\begin{Shaded}
\begin{Highlighting}[]
\NormalTok{src/dimensionwise\_quantizer.cpp}
\NormalTok{  per{-}dimension/per{-}bitwidth training}

\NormalTok{src/bit\_allocator.cpp}
\NormalTok{  total{-}budget bit{-}allocation DP}

\NormalTok{src/eval\_dimensionwise\_quantizer.cpp}
\end{Highlighting}
\end{Shaded}

\begin{Shaded}
\begin{Highlighting}[]
\NormalTok{  training, encoding, brute{-}force scan, recall output}
\end{Highlighting}
\end{Shaded}

Main evidence:

\begin{Shaded}
\begin{Highlighting}[]
\NormalTok{vectordb/docs/dimensionwise\_quantizer\_smoke\_run\_2026\_06\_26.md}
\NormalTok{vectordb/reports/scalar\_training\_exact\_hist\_audit\_2026\_06\_30/README.md}
\NormalTok{vectordb/reports/scalar\_training\_exact\_hist\_audit\_2026\_06\_30/comparison.csv}
\end{Highlighting}
\end{Shaded}

CLI:

\begin{Shaded}
\begin{Highlighting}[]
\NormalTok{eval\_dimensionwise\_quantizer ... [lloyd|exact{-}hist] [histogram\_bins]}
\end{Highlighting}
\end{Shaded}
\end{frame}

\begin{frame}[fragile,allowframebreaks]{37. Attempt 2 Complexity}
\protect\phantomsection\label{37-attempt-2-complexity}
For one dimension and one bitwidth after sorting:

\begin{itemize}
\tightlist
\item
  \textbf{Stage}: Equal-count histogram build

  \textbf{Complexity}: \texttt{O(N)}
\item
  \textbf{Stage}: Interval SSE query

  \textbf{Complexity}: \texttt{O(1)}
\item
  \textbf{Stage}: Divide-and-conquer histogram DP

  \textbf{Complexity}: approximately \texttt{O(K\_b\ H\ log\ H)}
\item
  \textbf{Stage}: Backtracking

  \textbf{Complexity}: \texttt{O(K\_b)}
\item
  \textbf{Stage}: Final raw-value SSE

  \textbf{Complexity}: \texttt{O(N\ log\ K\_b)}
\end{itemize}

Including the one-time dimension sort:

\begin{Shaded}
\begin{Highlighting}[]
\NormalTok{O(N log N + N log K\_b + K\_b H log H)}
\end{Highlighting}
\end{Shaded}

Variable mode trains candidate bitwidths \texttt{0..8}:

\begin{Shaded}
\begin{Highlighting}[]
\NormalTok{O(D * [N log N + sum\_b(N log K\_b + K\_b H log H)])}
\end{Highlighting}
\end{Shaded}

The outer allocation DP has approximately:

\begin{Shaded}
\begin{Highlighting}[]
\NormalTok{time   O(D * (D B) * |bitwidth choices|)}
\NormalTok{cost rows with rolling layers: O(D B)}
\NormalTok{backpointers for reconstruction: O(D * D B)}
\NormalTok{current total memory:           O(D * D B)}
\end{Highlighting}
\end{Shaded}

Its optimality is conditional:

\begin{Shaded}
\begin{Highlighting}[]
\NormalTok{Given one trainer{-}specific raw SSE table E[j,b], fixed weights w\_j, the}
\NormalTok{candidate set b in \{0,...,8\}, and a reachable budget, the outer DP finds the}
\NormalTok{minimum table objective. It does not make the inner codebooks or Recall}
\NormalTok{globally optimal.}
\end{Highlighting}
\end{Shaded}

Full-sample \texttt{H=N} exact training would replace the small
histogram term with approximately \texttt{O(K\_b\ N\ log\ N)} per
dimension/bitwidth. It was judged too expensive for full DEEP1M and was
not evaluated.
\end{frame}

\begin{frame}[fragile,allowframebreaks]{38. Exact-Hist Running Example}
\protect\phantomsection\label{38-exact-hist-running-example}
One-dimensional sorted values:

\begin{Shaded}
\begin{Highlighting}[]
\NormalTok{0, 1, 2, 10, 11, 12}
\end{Highlighting}
\end{Shaded}

Let:

\begin{Shaded}
\begin{Highlighting}[]
\NormalTok{H = 6 bins, one value/bin}
\NormalTok{K\_b = 2 centroids}
\end{Highlighting}
\end{Shaded}

The DP evaluates every legal split through the recurrence. The optimum
is:

\begin{Shaded}
\begin{Highlighting}[]
\NormalTok{cluster 1 = \{0,1,2\},  centroid = 1}
\NormalTok{cluster 2 = \{10,11,12\}, centroid = 11}
\end{Highlighting}
\end{Shaded}

SSE:

\begin{Shaded}
\begin{Highlighting}[]
\NormalTok{(0{-}1)\^{}2 + (1{-}1)\^{}2 + (2{-}1)\^{}2}
\NormalTok{+ (10{-}11)\^{}2 + (11{-}11)\^{}2 + (12{-}11)\^{}2}
\NormalTok{= 4}
\end{Highlighting}
\end{Shaded}

With \texttt{H=2}, each bin would already contain three values. The same
partition is available, but an optimal raw split inside a bin would not
be. This is the difference between histogram exactness and raw-sample
exactness.
\end{frame}

\begin{frame}[fragile,allowframebreaks]{39. Audio Result: Lloyd Versus
Exact-Hist}
\protect\phantomsection\label{39-audio-result-lloyd-versus-exact-hist}
Dataset:

\begin{Shaded}
\begin{Highlighting}[]
\NormalTok{audio base = 53,387 x 192}
\NormalTok{audio query = 200 x 192}
\NormalTok{metric = L2}
\NormalTok{mode = variable}
\NormalTok{B = 4}
\NormalTok{H = 256}
\NormalTok{source = initial H=256 comparison run}
\end{Highlighting}
\end{Shaded}

\begin{itemize}
\tightlist
\item
  \textbf{Trainer}: Lloyd

  \textbf{MSE/value}: 280,023
\item
  \textbf{R@10}: 0.9260

  \textbf{R@100}: 0.95275
\item
  \textbf{Training time}: 2,004.68 ms
\item
  \textbf{Trainer}: exact-hist

  \textbf{MSE/value}: 236,478
\item
  \textbf{R@10}: 0.9270

  \textbf{R@100}: 0.95675
\item
  \textbf{Training time}: 1,864.28 ms
\end{itemize}

At this operating point:

\begin{Shaded}
\begin{Highlighting}[]
\NormalTok{MSE improvement \textasciitilde{}= 15.5\%}
\NormalTok{R@10 delta       = +0.0010}
\NormalTok{R@100 delta      = +0.0040}
\end{Highlighting}
\end{Shaded}

This is a positive example, but it does not establish stable recall
dominance or a new scalar-quantization contribution.
\end{frame}

\begin{frame}[fragile,allowframebreaks]{40. Audio Histogram-Bin Sweep}
\protect\phantomsection\label{40-audio-histogram-bin-sweep}
\texttt{B=4}, variable allocation, later sweep rerun:

\begin{itemize}
\tightlist
\item
  \textbf{Histogram bins H}: 256

  \textbf{MSE/value}: 236,478
\item
  \textbf{R@10}: 0.9270

  \textbf{R@100}: \textbf{0.95675}
\item
  \textbf{Training time}: 1,776 ms
\item
  \textbf{Histogram bins H}: 512

  \textbf{MSE/value}: 235,292
\item
  \textbf{R@10}: 0.9340

  \textbf{R@100}: 0.95615
\item
  \textbf{Training time}: 2,418 ms
\item
  \textbf{Histogram bins H}: 1,024

  \textbf{MSE/value}: 234,918
\item
  \textbf{R@10}: \textbf{0.9345}

  \textbf{R@100}: 0.95610
\item
  \textbf{Training time}: 3,858 ms
\item
  \textbf{Histogram bins H}: 2,048

  \textbf{MSE/value}: \textbf{234,844}
\item
  \textbf{R@10}: 0.9335

  \textbf{R@100}: 0.95625
\item
  \textbf{Training time}: 7,018 ms
\end{itemize}

Observation:

\begin{Shaded}
\begin{Highlighting}[]
\NormalTok{Increasing H improves the reconstruction objective almost monotonically,}
\NormalTok{but R@100 is not monotonic and training time grows substantially.}
\end{Highlighting}
\end{Shaded}

At \texttt{B=8}, \texttt{H=256} is especially coarse because
\texttt{K\_b=256}: one centroid per histogram bin leaves little
partition freedom. Increasing \texttt{H} improves the result, but the
best MSE and best recall still occur at different bin counts.

The \texttt{H=256} training time here and on the previous slide comes
from separate runs; it is not an internal consistency check on timing
noise.
\end{frame}

\begin{frame}[fragile,allowframebreaks]{41. PCA CIFAR/DEEP Comparison}
\protect\phantomsection\label{41-pca-cifardeep-comparison}
All rows use \texttt{H=256} and variable bit allocation. Raw SSE ratio
is the reported unweighted nearest-centroid \texttt{training\_sse},
\texttt{exact-hist\ /\ Lloyd}; recall delta is
\texttt{exact-hist\ -\ Lloyd}.

\begin{itemize}
\tightlist
\item
  \textbf{Dataset}: CIFAR60K

  \textbf{B}: 2
\item
  \textbf{Allocation objective}: reconstruction

  \textbf{Raw SSE ratio}: 0.9737
\item
  \textbf{R@100 delta}: +0.00313
\item
  \textbf{Dataset}: CIFAR60K

  \textbf{B}: 2
\item
  \textbf{Allocation objective}: rank-boundary

  \textbf{Raw SSE ratio}: 0.9336
\item
  \textbf{R@100 delta}: +0.00378
\item
  \textbf{Dataset}: CIFAR60K

  \textbf{B}: 4
\item
  \textbf{Allocation objective}: reconstruction

  \textbf{Raw SSE ratio}: 1.0425
\item
  \textbf{R@100 delta}: -0.00024
\item
  \textbf{Dataset}: CIFAR60K

  \textbf{B}: 4
\item
  \textbf{Allocation objective}: rank-boundary

  \textbf{Raw SSE ratio}: 0.9371
\item
  \textbf{R@100 delta}: -0.00046
\item
  \textbf{Dataset}: DEEP1M

  \textbf{B}: 2
\item
  \textbf{Allocation objective}: reconstruction

  \textbf{Raw SSE ratio}: 0.9774
\item
  \textbf{R@100 delta}: +0.00751
\item
  \textbf{Dataset}: DEEP1M

  \textbf{B}: 2
\item
  \textbf{Allocation objective}: rank-boundary

  \textbf{Raw SSE ratio}: 0.9477
\item
  \textbf{R@100 delta}: +0.00681
\item
  \textbf{Dataset}: DEEP1M

  \textbf{B}: 4
\item
  \textbf{Allocation objective}: reconstruction

  \textbf{Raw SSE ratio}: 0.8847
\item
  \textbf{R@100 delta}: +0.00304
\item
  \textbf{Dataset}: DEEP1M

  \textbf{B}: 4
\item
  \textbf{Allocation objective}: rank-boundary

  \textbf{Raw SSE ratio}: \textbf{0.8016}
\item
  \textbf{R@100 delta}: \textbf{-0.00215}
\end{itemize}

Two different effects appear:

\begin{enumerate}
\tightlist
\item
  CIFAR \texttt{B=4} reconstruction shows histogram exactness does not
  guarantee lower raw-sample SSE.
\item
  DEEP \texttt{B=4} rank-boundary shows that much lower unweighted raw
  reconstruction SSE still does not guarantee higher recall.
\end{enumerate}

Important correction to the earlier deck interpretation:

\begin{Shaded}
\begin{Highlighting}[]
\NormalTok{DEEP B=4 rank{-}boundary allocation objective}
\NormalTok{Lloyd:      1070.81}
\NormalTok{exact{-}hist: 1551.93}
\end{Highlighting}
\end{Shaded}

The weighted outer objective worsened in that cell. Therefore it is
evidence against using raw SSE as a recall surrogate, not evidence that
a lower rank-boundary allocation objective failed to improve recall.
\end{frame}

\begin{frame}[fragile,allowframebreaks]{42. Why Lower Reconstruction SSE
Need Not Improve Recall}
\protect\phantomsection\label{42-why-lower-reconstruction-sse-need-not-improve-recall}
Define base-vector quantization error:

\[e_x=\hat x-x.\]

Only the base is quantized; the query remains float. Distance error is:

\[\Delta_x=
\hat d(q,x)-d(q,x)
=-2(q-x)^Te_x+\lVert e_x\rVert^2.\]

Scalar k-means minimizes the second-order reconstruction term aggregated
over base vectors:

\[\sum_x\lVert e_x\rVert^2.\]

It does not directly control:

\[-2(q-x)^Te_x,\]

which depends on query direction and signed correlation with the
quantization error.

For a positive \texttt{p} and negative \texttt{n}, exact margin is:

\[m=d(q,n)-d(q,p)>0.\]

Estimated margin is:

\[\hat m=m+\Delta_n-\Delta_p.\]

Recall depends on the sign of this margin, not average reconstruction
SSE.
\end{frame}

\begin{frame}[fragile,allowframebreaks]{43. Ranking Running Example}
\protect\phantomsection\label{43-ranking-running-example}
Exact distances near the retrieval boundary:

\begin{Shaded}
\begin{Highlighting}[]
\NormalTok{positive p = 1.000}
\NormalTok{negative n = 1.010}
\NormalTok{margin     = 0.010}
\end{Highlighting}
\end{Shaded}

Method A has smaller but oppositely signed errors:

\begin{Shaded}
\begin{Highlighting}[]
\NormalTok{Delta\_p = +0.006}
\NormalTok{Delta\_n = {-}0.006}

\NormalTok{estimated p = 1.006}
\NormalTok{estimated n = 1.004}
\NormalTok{result: inversion}
\end{Highlighting}
\end{Shaded}

Method B has larger common-mode error:

\begin{Shaded}
\begin{Highlighting}[]
\NormalTok{Delta\_p = +0.020}
\NormalTok{Delta\_n = +0.020}

\NormalTok{estimated p = 1.020}
\NormalTok{estimated n = 1.030}
\NormalTok{result: ordering preserved}
\end{Highlighting}
\end{Shaded}

This is why a codebook or coupled variable-bit pipeline with lower raw
reconstruction SSE can nevertheless have lower recall: ranking depends
on differential error across close candidates.
\end{frame}

\begin{frame}[fragile,allowframebreaks]{44. Bit-Allocation Coupling}
\protect\phantomsection\label{44-bit-allocation-coupling}
Changing the scalar trainer changes the full error table:

\[E[j,b]=\text{raw nearest-centroid SSE for dimension j at bitwidth b}.\]

The outer DP then selects a different bit vector:

\[(b_1,\ldots,b_D).\]

DEEP1M \texttt{B=4}, historical rank-boundary objective:

\begin{Shaded}
\begin{Highlighting}[]
\NormalTok{Lloyd histogram:}
\NormalTok{1b:61, 2b:36, 3b:33, 4b:27, 5b:19, 6b:19, 7b:13, 8b:48}

\NormalTok{Exact{-}hist histogram:}
\NormalTok{1b:40, 2b:48, 3b:33, 4b:30, 5b:26, 6b:28, 7b:37, 8b:14}
\end{Highlighting}
\end{Shaded}

Consequently recall differences combine:

\begin{Shaded}
\begin{Highlighting}[]
\NormalTok{centroid movement}
\NormalTok{encoding{-}boundary movement}
\NormalTok{dimension bit reallocation}
\NormalTok{surrogate{-}objective mismatch}
\end{Highlighting}
\end{Shaded}

Given this trainer-specific table and fixed weights, the outer DP is
exact over the enumerated bitwidths and reachable total budget. It does
not isolate a pure inner-trainer effect: changing the trainer also
changes the selected bit vector.

Historical disclosure:

\begin{Shaded}
\begin{Highlighting}[]
\NormalTok{rank{-}boundary weights use query/ground{-}truth{-}derived pairs.}
\NormalTok{They are evidence about objective behavior, not an admissible method under the}
\NormalTok{current query{-}unaware project constraint.}
\end{Highlighting}
\end{Shaded}
\end{frame}

\begin{frame}[fragile,allowframebreaks]{45. Attempt 2 Decision}
\protect\phantomsection\label{45-attempt-2-decision}
\begin{Shaded}
\begin{Highlighting}[]
\NormalTok{Decision: KEEP EXACT{-}HIST AS A STRONG OFFLINE BASELINE, NOT A MAIN METHOD}
\end{Highlighting}
\end{Shaded}

What the experiment establishes:

\begin{enumerate}
\tightlist
\item
  The bin-restricted DP is implemented and can materially lower reported
  raw reconstruction SSE after nearest-centroid reassignment.
\item
  It sometimes improves recall, especially at lower bit budgets.
\item
  Historical rank-boundary gains over reconstruction survive both Lloyd
  and exact-hist, so they are not only a Lloyd artifact.
\item
  Lower raw SSE is not a stable surrogate for ANN ranking quality.
\end{enumerate}

What the related-work audit establishes:

\begin{enumerate}
\tightlist
\item
  The White-Singal paper and pinned public code do not document or
  expose an integrated shared-codebook, heterogeneous-bit,
  outer-allocation pipeline.
\item
  The pinned code separately exposes a directly adjacent shared-block
  helper and a full-raw-data exact 1D k-means solver primitive; neither
  is shown as an integrated or evaluated Attempt 2 pipeline.
\item
  Exact 1D scalar clustering is much older than either project; the
  formal paper explicitly recognizes that prior.
\end{enumerate}

Why it is not a method contribution:

\begin{enumerate}
\tightlist
\item
  Split points are exact only under a chosen histogram boundary
  restriction.
\item
  Full-sample exact DP was not evaluated at full scale.
\item
  Lower SSE does not provide a recall guarantee.
\item
  Recall gains are non-monotonic across \texttt{H}, datasets, budgets,
  and allocation objectives.
\item
  The inner optimizer\textquotesingle s novelty is foreclosed by Wu 1991
  / GLM 2017; outer DP novelty is not claimed here, and SAQ already
  supplies direct ANN prior.
\item
  The historical rank-boundary outer objective is
  query/ground-truth-aware.
\item
  Unpacked byte codes and brute-force L2 scan do not establish a
  competitive packed variable-rate ANN system.
\end{enumerate}

Strict-reviewer summary:

\begin{Shaded}
\begin{Highlighting}[]
\NormalTok{No integrated end{-}to{-}end pipeline is visible in the paper/public code, but}
\NormalTok{that absence does not establish combination novelty. The inner optimizer is}
\NormalTok{prior art and the unbuilt system layer remains unvalidated. This is a useful}
\NormalTok{baseline and objective{-}mismatch result, not a new ANN quantizer or}
\NormalTok{database{-}systems method.}
\end{Highlighting}
\end{Shaded}
\end{frame}

\begin{frame}[fragile,allowframebreaks]{46. Attempt 3: Research
Question}
\protect\phantomsection\label{46-attempt-3-research-question}
Recall@\texttt{k} measures exact identifier overlap. For exact
top-\texttt{k} set \texttt{E\_k(q)} and returned set \texttt{A\_k(q)}:

\begin{Shaded}
\begin{Highlighting}[]
\NormalTok{Recall@k(q) = |E\_k(q) intersect A\_k(q)| / k.}
\end{Highlighting}
\end{Shaded}

This can penalize a geometrically near-equivalent replacement at a dense
top-\texttt{k} boundary. Attempt 3 therefore asks:

\begin{Shaded}
\begin{Highlighting}[]
\NormalTok{Does paper{-}exact distance quality change the Pareto interpretation of any}
\NormalTok{previously frozen SAQ comparison, when both plans are evaluated on the same}
\NormalTok{search{-}effort grid?}
\end{Highlighting}
\end{Shaded}

The premise is about evaluation validity, not a new quantizer:

\begin{Shaded}
\begin{Highlighting}[]
\NormalTok{Recall:       Did we return the same identifiers?}
\NormalTok{1/Ratio:      How far are the returned vectors compared with the exact ranks?}
\end{Highlighting}
\end{Shaded}
\end{frame}

\begin{frame}[fragile,allowframebreaks]{47. Attempt 3 Metric:
\texttt{1/Ratio@k}}
\protect\phantomsection\label{47-attempt-3-metric-1ratiok}
For query \texttt{q}, define:

\begin{itemize}
\tightlist
\item
  \texttt{d\_i(q)}: true Euclidean distance to the exact neighbor at
  distance rank \texttt{i}, for \texttt{i=1,...,k};
\item
  \texttt{d\_tilde\_i(q)}: true Euclidean distance at rank \texttt{i}
  after independently sorting the \texttt{k} returned identifiers by
  true distance.
\end{itemize}

The metric from \href{https://arxiv.org/abs/2606.04522v1}{ANN Search:
Recall What Matters} is:

\begin{Shaded}
\begin{Highlighting}[]
\NormalTok{Ratio@k(q) = (1/k) * sum\_\{i=1\}\^{}k d\_tilde\_i(q) / d\_i(q)}

\NormalTok{1/Ratio@k(q) = k / sum\_\{i=1\}\^{}k d\_tilde\_i(q) / d\_i(q).}
\end{Highlighting}
\end{Shaded}

Properties under valid exact top-\texttt{k} input:

\begin{Shaded}
\begin{Highlighting}[]
\NormalTok{0 \textless{} 1/Ratio@k \textless{}= 1}
\NormalTok{higher is better}
\NormalTok{1 means equal distance quality, even if tied identifiers differ}
\end{Highlighting}
\end{Shaded}

Critical semantic detail: SAQ internally reports squared L2 distances,
but this metric uses Euclidean distances. The evaluator must take square
roots before forming the ratios.
\end{frame}

\begin{frame}[fragile,allowframebreaks]{48. Attempt 3 Running Example}
\protect\phantomsection\label{48-attempt-3-running-example}
Let \texttt{k=3}. Suppose the exact result distances are:

\begin{Shaded}
\begin{Highlighting}[]
\NormalTok{d = [1.00, 1.10, 1.20].}
\end{Highlighting}
\end{Shaded}

The ANN result shares only one exact identifier, but after recomputing
and sorting true distances it has:

\begin{Shaded}
\begin{Highlighting}[]
\NormalTok{d\_tilde = [1.00, 1.11, 1.21].}
\end{Highlighting}
\end{Shaded}

Then:

\begin{Shaded}
\begin{Highlighting}[]
\NormalTok{Recall@3 = 1/3 = 0.3333}

\NormalTok{1/Ratio@3}
\NormalTok{  = 3 / (1.00/1.00 + 1.11/1.10 + 1.21/1.20)}
\NormalTok{  = 0.9942.}
\end{Highlighting}
\end{Shaded}

Interpretation:

\begin{Shaded}
\begin{Highlighting}[]
\NormalTok{identifier recovery is poor}
\NormalTok{geometric result quality is nearly exact}
\end{Highlighting}
\end{Shaded}

This example does not imply that Recall is wrong. It shows that the two
metrics answer different scientific questions.
\end{frame}

\begin{frame}[fragile,allowframebreaks]{49. Related Work And Novelty
Gate}
\protect\phantomsection\label{49-related-work-and-novelty-gate}
The metric paper already contributes:

\begin{enumerate}
\tightlist
\item
  the judge-free, hyperparameter-free \texttt{1/Ratio@k} definition;
\item
  evaluation of Annoy, SuCo, HNSW, RaBitQ, and SymphonyQG on six
  datasets;
\item
  build, memory, distance-work, classification, and RAG analyses; and
\item
  evidence that less search effort can satisfy a distance-quality target
  than an equal-valued Recall target.
\end{enumerate}

At \texttt{k=100} and quality \texttt{0.95}, it reports the following
average ratios between the effort required by Recall and by
\texttt{1/Ratio}:

\begin{itemize}
\tightlist
\item
  \textbf{Method}: HNSW

  \textbf{Relative distance-computation effort}: \texttt{9.36x}
\item
  \textbf{Method}: RaBitQ

  \textbf{Relative distance-computation effort}: \texttt{2.48x}
\item
  \textbf{Method}: SymphonyQG

  \textbf{Relative distance-computation effort}: \texttt{2.38x}
\item
  \textbf{Method}: SuCo

  \textbf{Relative distance-computation effort}: \texttt{1.86x}
\item
  \textbf{Method}: Annoy

  \textbf{Relative distance-computation effort}: \texttt{3.22x}
\end{itemize}

Strict-reviewer constraint:

\begin{Shaded}
\begin{Highlighting}[]
\NormalTok{The metric is prior work, and the paper reports that relative algorithm}
\NormalTok{rankings are usually stable. Re{-}evaluating SAQ is not a method contribution.}
\end{Highlighting}
\end{Shaded}

Attempt 3 was therefore authorized only as a frozen retrospective
falsification study. It could clarify prior evidence, but could not fit
a plan to benchmark queries or claim the metric as novelty.
\end{frame}

\begin{frame}[fragile,allowframebreaks]{50. Frozen Attempt 3 Protocol}
\protect\phantomsection\label{50-frozen-attempt-3-protocol}
Stages:

\begin{Shaded}
\begin{Highlighting}[]
\NormalTok{A3{-}0  validate the paper{-}exact evaluator on deterministic fixtures}
\NormalTok{A3{-}1  replay the frozen GIST default{-}versus{-}fac{-}error comparison}
\NormalTok{A3{-}2  test frozen DEEP B=4/B=5 rejected plans as negative controls}
\NormalTok{A3{-}3  apply the preregistered cross{-}dataset decision gate}
\end{Highlighting}
\end{Shaded}

\begin{itemize}
\tightlist
\item
  \textbf{Variable}: base vectors \texttt{N}

  \textbf{GIST A3-1}: 100,000
\item
  \textbf{DEEP A3-2}: 100,000
\item
  \textbf{Variable}: queries \texttt{Q}

  \textbf{GIST A3-1}: 1,000
\item
  \textbf{DEEP A3-2}: 1,000
\item
  \textbf{Variable}: dimension \texttt{D}

  \textbf{GIST A3-1}: 960
\item
  \textbf{DEEP A3-2}: 256
\item
  \textbf{Variable}: IVF clusters \texttt{K\_IVF}

  \textbf{GIST A3-1}: 512
\item
  \textbf{DEEP A3-2}: 512
\item
  \textbf{Variable}: nominal budget \texttt{B}

  \textbf{GIST A3-1}: 4
\item
  \textbf{DEEP A3-2}: 4 and 5
\item
  \textbf{Variable}: result size \texttt{k}

  \textbf{GIST A3-1}: 100
\item
  \textbf{DEEP A3-2}: 100
\item
  \textbf{Variable}: \texttt{nprobe} grid

  \textbf{GIST A3-1}: 20, 50, 100, 160, 200, 220, 240, 280, 300, 320,
  400
\item
  \textbf{DEEP A3-2}: 50, 100, 200, 400
\end{itemize}

Frozen common settings:

\begin{Shaded}
\begin{Highlighting}[]
\NormalTok{threads = 24}
\NormalTok{QPS repetitions = 10 per row}
\NormalTok{searcher\_safe\_block\_min\_mode = 2}
\NormalTok{searcher\_vars\_bound\_m = 4}
\NormalTok{PCA = enabled, retaining all D dimensions}
\end{Highlighting}
\end{Shaded}

\texttt{nprobe} is the number of IVF cells scanned. Mode \texttt{2} is
the previously fixed correctness-preserving block-min path, and
\texttt{m=4} is SAQ\textquotesingle s unchanged variance bound
multiplier. Neither is a parameter of \texttt{1/Ratio}; no value was
selected from Attempt

3 outcomes.
\end{frame}

\begin{frame}[fragile,allowframebreaks]{51. Evaluator And Search-Code
Semantics}
\protect\phantomsection\label{51-evaluator-and-search-code-semantics}
Core evaluator logic on branch \texttt{saq-ratio-metric-analysis}:

\begin{Shaded}
\begin{Highlighting}[]
\NormalTok{exact }\OperatorTok{=} \BuiltInTok{sorted}\NormalTok{(true\_l2(base[groundtruth\_ids], query))}
\NormalTok{returned }\OperatorTok{=} \BuiltInTok{sorted}\NormalTok{(true\_l2(base[result\_ids], query))}
\NormalTok{inverse\_ratio }\OperatorTok{=}\NormalTok{ k }\OperatorTok{/} \BuiltInTok{sum}\NormalTok{(returned[i] }\OperatorTok{/}\NormalTok{ exact[i] }\ControlFlowTok{for}\NormalTok{ i }\KeywordTok{in} \BuiltInTok{range}\NormalTok{(k))}
\end{Highlighting}
\end{Shaded}

The implementation in \texttt{script/evaluate\_inverse\_ratio.py}:

\begin{enumerate}
\tightlist
\item
  recomputes both sets of distances in float64;
\item
  independently sorts by true Euclidean distance;
\item
  rejects duplicate, invalid, truncated, or non-finite result rows;
\item
  stops explicitly if an exact distance is zero; and
\item
  preserves every query-level value instead of only the mean.
\end{enumerate}

The existing upstream \texttt{utils::get\_ratio} was not used as the
scientific artifact because it reports forward \texttt{Ratio}, skips
tiny squared distances with an epsilon, and discards query-level values.

\texttt{src/test\_qps.cpp} exports result IDs only after the timed
search region. Thus the evaluator can replay exact returned sets without
inflating measured QPS. Eighteen deterministic Python tests cover exact,
tied-distance, reordered, squared-L2, malformed-input, provenance, and
frontier cases.
\end{frame}

\begin{frame}[fragile,allowframebreaks]{52. Attempt 3 Complexity And
Overhead}
\protect\phantomsection\label{52-attempt-3-complexity-and-overhead}
Given \texttt{Q} queries, returned top-\texttt{k}, and dimension
\texttt{D}:

\begin{Shaded}
\begin{Highlighting}[]
\NormalTok{recompute exact and returned distances: O(Q k D) time}
\NormalTok{sort each query\textquotesingle{}s distance lists:       O(Q k log k) time}
\NormalTok{Recall and ratio aggregation:           O(Q k) time}
\NormalTok{stored result identifiers:              O(Q k)}
\NormalTok{stored per{-}query metric values:         O(Q)}
\end{Highlighting}
\end{Shaded}

No \texttt{O(QND)} exhaustive search is required when exact
top-\texttt{k} identifiers already exist. Only the \texttt{2Qk}
selected-vector distances are recomputed.

Separation of costs:

\begin{Shaded}
\begin{Highlighting}[]
\NormalTok{ANN query latency / QPS: unchanged; result export is outside the timer}
\NormalTok{benchmark evaluation:    O(Q k D), reported separately}
\NormalTok{index build and storage: unchanged}
\end{Highlighting}
\end{Shaded}

The evaluator has no learned parameters. \texttt{k} is the benchmark
result size, and the \texttt{nprobe} values describe the measured
search-effort curve rather than a fitted decision rule.
\end{frame}

\begin{frame}[fragile,allowframebreaks]{53. GIST Comparison: Two Frozen
Global Plans}
\protect\phantomsection\label{53-gist-comparison-two-frozen-global-plans}
Setting: \texttt{gist\_sample100k}, \texttt{K\_IVF=512}, \texttt{B=4},
\texttt{D=960}, \texttt{k=100}.

\begin{Shaded}
\begin{Highlighting}[]
\NormalTok{SAQ default: 64x11 | 192x6 | 320x4 | 256x2 | 128x0}
\NormalTok{fac{-}error:   192x9 | 512x4 | 256x0}
\end{Highlighting}
\end{Shaded}

Each \texttt{dimensions\ x\ bits} term means that every coordinate in
the contiguous PCA segment receives that many code bits. Both plan
strings sum to 960 dimensions and were materialized as frozen SAQ
\texttt{B=4} indexes before Attempt 3. The plan string

alone is not total index-byte accounting because SAQ also stores
per-segment factors.

The comparison changes only the global plan:

\begin{Shaded}
\begin{Highlighting}[]
\NormalTok{same PCA vectors and IVF assignments}
\NormalTok{same CAQ encoder and serialized index format}
\NormalTok{same distance estimator, pruning, heap, and safe{-}search mode}
\NormalTok{same nprobe grid, threads, and repetitions}
\end{Highlighting}
\end{Shaded}

The fac-error plan has two positive-bit segments instead of four, so it
is a plausible lower-search-work shape. Its name refers to the
historical offline plan artifact; Attempt 3 neither refits nor endorses
that empirical objective.
\end{frame}

\begin{frame}[fragile,allowframebreaks]{54. GIST Complete Common-Grid
Result}
\protect\phantomsection\label{54-gist-complete-common-grid-result}
\begin{itemize}
\tightlist
\item
  \textbf{nprobe}: 20

  \textbf{default R@100}: 0.75957
\item
  \textbf{fac R@100}: 0.75952

  \textbf{default 1/Ratio}: 0.993860620
\item
  \textbf{fac 1/Ratio}: 0.993858359

  \textbf{default QPS}: 31234.2
\item
  \textbf{fac QPS}: 36687.1
\item
  \textbf{nprobe}: 50

  \textbf{default R@100}: 0.92809
\item
  \textbf{fac R@100}: 0.92769

  \textbf{default 1/Ratio}: 0.998671018
\item
  \textbf{fac 1/Ratio}: 0.998668346

  \textbf{default QPS}: 20104.2
\item
  \textbf{fac QPS}: 25581.3
\item
  \textbf{nprobe}: 100

  \textbf{default R@100}: 0.98140
\item
  \textbf{fac R@100}: 0.98069

  \textbf{default 1/Ratio}: 0.999773561
\item
  \textbf{fac 1/Ratio}: 0.999769939

  \textbf{default QPS}: 13825.0
\item
  \textbf{fac QPS}: 18572.6
\item
  \textbf{nprobe}: 160

  \textbf{default R@100}: 0.99036
\item
  \textbf{fac R@100}: 0.98958

  \textbf{default 1/Ratio}: 0.999964379
\item
  \textbf{fac 1/Ratio}: 0.999960461

  \textbf{default QPS}: 10524.5
\item
  \textbf{fac QPS}: 14189.3
\item
  \textbf{nprobe}: 200

  \textbf{default R@100}: 0.99132
\item
  \textbf{fac R@100}: 0.99059

  \textbf{default 1/Ratio}: 0.999985301
\item
  \textbf{fac 1/Ratio}: 0.999981452

  \textbf{default QPS}: 9264.9
\item
  \textbf{fac QPS}: 12333.9
\item
  \textbf{nprobe}: 220

  \textbf{default R@100}: 0.99149
\item
  \textbf{fac R@100}: 0.99079

  \textbf{default 1/Ratio}: 0.999987377
\item
  \textbf{fac 1/Ratio}: 0.999983555

  \textbf{default QPS}: 8777.6
\item
  \textbf{fac QPS}: 11739.0
\item
  \textbf{nprobe}: 240

  \textbf{default R@100}: 0.99153
\item
  \textbf{fac R@100}: 0.99084

  \textbf{default 1/Ratio}: 0.999988670
\item
  \textbf{fac 1/Ratio}: 0.999984881

  \textbf{default QPS}: 8313.8
\item
  \textbf{fac QPS}: 11073.8
\item
  \textbf{nprobe}: 280

  \textbf{default R@100}: 0.99163
\item
  \textbf{fac R@100}: 0.99091

  \textbf{default 1/Ratio}: 0.999991688
\item
  \textbf{fac 1/Ratio}: 0.999987869

  \textbf{default QPS}: 7579.7
\item
  \textbf{fac QPS}: 9986.9
\item
  \textbf{nprobe}: 300

  \textbf{default R@100}: 0.99164
\item
  \textbf{fac R@100}: 0.99091

  \textbf{default 1/Ratio}: 0.999992743
\item
  \textbf{fac 1/Ratio}: 0.999988921

  \textbf{default QPS}: 7329.5
\item
  \textbf{fac QPS}: 9522.1
\item
  \textbf{nprobe}: 320

  \textbf{default R@100}: 0.99164
\item
  \textbf{fac R@100}: 0.99091

  \textbf{default 1/Ratio}: 0.999993271
\item
  \textbf{fac 1/Ratio}: 0.999989447

  \textbf{default QPS}: 7035.6
\item
  \textbf{fac QPS}: 9086.9
\item
  \textbf{nprobe}: 400

  \textbf{default R@100}: 0.99164
\item
  \textbf{fac R@100}: 0.99091

  \textbf{default 1/Ratio}: 0.999993271
\item
  \textbf{fac 1/Ratio}: 0.999989447

  \textbf{default QPS}: 6111.4
\item
  \textbf{fac QPS}: 7784.3
\end{itemize}

At every equal \texttt{nprobe}, fac-error is faster and slightly worse
under both quality metrics. Equal search effort is therefore not the
scientific comparison; the relevant question is QPS at matched quality.
\end{frame}

\begin{frame}[fragile,allowframebreaks]{55. GIST Metric-Sensitive
Operating Point}
\protect\phantomsection\label{55-gist-metric-sensitive-operating-point}
Use the freshly replayed default \texttt{nprobe=200} row as the frozen
target:

\begin{Shaded}
\begin{Highlighting}[]
\NormalTok{default Recall@100:    0.991320000000}
\NormalTok{default 1/Ratio@100:   0.999985300968}
\NormalTok{default QPS:           9264.923}
\end{Highlighting}
\end{Shaded}

Recall conclusion:

\begin{Shaded}
\begin{Highlighting}[]
\NormalTok{maximum measured fac{-}error Recall = 0.99091}
\NormalTok{target 0.99132 is unreachable}
\NormalTok{decision: reject at this reference quality}
\end{Highlighting}
\end{Shaded}

Distance-quality conclusion:

\begin{Shaded}
\begin{Highlighting}[]
\NormalTok{fac{-}error nprobe=280  1/Ratio@100 = 0.999987869279}
\NormalTok{fac{-}error nprobe=280  QPS         = 9986.910}
\NormalTok{quality               \textgreater{} default target}
\NormalTok{QPS ratio             = 1.07793x}
\end{Highlighting}
\end{Shaded}

This is a directly measured point, not interpolation. Linear
interpolation between fac-error \texttt{nprobe=240} and \texttt{280}
estimates \texttt{1.17876x} QPS exactly at the target, but this is
reported only as secondary evidence.
\end{frame}

\begin{frame}[fragile,allowframebreaks]{56. GIST Query-Level
Attribution}
\protect\phantomsection\label{56-gist-query-level-attribution}
For fac-error \texttt{nprobe=280} minus default \texttt{nprobe=200}:

\begin{itemize}
\tightlist
\item
  \textbf{Statistic of paired 1/Ratio@100 delta}: mean

  \textbf{Value}: \texttt{+2.5683e-6}
\item
  \textbf{Statistic of paired 1/Ratio@100 delta}: median

  \textbf{Value}: \texttt{0}
\item
  \textbf{Statistic of paired 1/Ratio@100 delta}: minimum

  \textbf{Value}: \texttt{-3.1166e-4}
\item
  \textbf{Statistic of paired 1/Ratio@100 delta}: p01

  \textbf{Value}: \texttt{-1.1754e-4}
\item
  \textbf{Statistic of paired 1/Ratio@100 delta}: p05

  \textbf{Value}: \texttt{-4.1993e-5}
\item
  \textbf{Statistic of paired 1/Ratio@100 delta}: queries worse

  \textbf{Value}: \texttt{39.6\%}
\item
  \textbf{Statistic of paired 1/Ratio@100 delta}: queries equal

  \textbf{Value}: \texttt{23.0\%}
\item
  \textbf{Statistic of paired 1/Ratio@100 delta}: queries better

  \textbf{Value}: \texttt{37.4\%}
\end{itemize}

The alternative improves the aggregate mean and its separately measured
marginal lower-tail statistics, but it does not dominate query by query.

Strict interpretation:

\begin{Shaded}
\begin{Highlighting}[]
\NormalTok{supported:   the historical GIST operating{-}point rejection is metric{-}sensitive}
\NormalTok{unsupported: fac{-}error is uniformly better or provides a per{-}query guarantee}
\end{Highlighting}
\end{Shaded}

This paired distribution is why the mean metric alone is insufficient
for a new method claim.
\end{frame}

\begin{frame}[fragile,allowframebreaks]{57. DEEP Negative Controls}
\protect\phantomsection\label{57-deep-negative-controls}
Frozen plans:

\begin{Shaded}
\begin{Highlighting}[]
\NormalTok{B=4 default:   64x6 | 192x3     candidate: 128x4 | 128x3}
\NormalTok{B=5 default:   64x7 | 192x4     candidate: 128x5 | 128x4}
\end{Highlighting}
\end{Shaded}

\begin{itemize}
\tightlist
\item
  \textbf{B}: 4

  \textbf{nprobe}: 50
\item
  \textbf{default R}: 0.94550

  \textbf{candidate R}: 0.92381
\item
  \textbf{default 1/Ratio}: 0.998803057

  \textbf{candidate 1/Ratio}: 0.998540024
\item
  \textbf{QPS ratio}: 1.075x
\item
  \textbf{B}: 4

  \textbf{nprobe}: 100
\item
  \textbf{default R}: 0.97061

  \textbf{candidate R}: 0.94401
\item
  \textbf{default 1/Ratio}: 0.999686070

  \textbf{candidate 1/Ratio}: 0.999416524
\item
  \textbf{QPS ratio}: 1.095x
\item
  \textbf{B}: 4

  \textbf{nprobe}: 200
\item
  \textbf{default R}: 0.97641

  \textbf{candidate R}: 0.94884
\item
  \textbf{default 1/Ratio}: 0.999904123

  \textbf{candidate 1/Ratio}: 0.999633936
\item
  \textbf{QPS ratio}: 1.071x
\item
  \textbf{B}: 4

  \textbf{nprobe}: 400
\item
  \textbf{default R}: 0.97713

  \textbf{candidate R}: 0.94940
\item
  \textbf{default 1/Ratio}: 0.999930322

  \textbf{candidate 1/Ratio}: 0.999658856
\item
  \textbf{QPS ratio}: 1.046x
\item
  \textbf{B}: 5

  \textbf{nprobe}: 50
\item
  \textbf{default R}: 0.95180

  \textbf{candidate R}: 0.94251
\item
  \textbf{default 1/Ratio}: 0.998842845

  \textbf{candidate 1/Ratio}: 0.998776284
\item
  \textbf{QPS ratio}: 1.108x
\item
  \textbf{B}: 5

  \textbf{nprobe}: 100
\item
  \textbf{default R}: 0.97973

  \textbf{candidate R}: 0.96687
\item
  \textbf{default 1/Ratio}: 0.999728710

  \textbf{candidate 1/Ratio}: 0.999659199
\item
  \textbf{QPS ratio}: 1.092x
\item
  \textbf{B}: 5

  \textbf{nprobe}: 200
\item
  \textbf{default R}: 0.98670

  \textbf{candidate R}: 0.97241
\item
  \textbf{default 1/Ratio}: 0.999948614

  \textbf{candidate 1/Ratio}: 0.999877534
\item
  \textbf{QPS ratio}: 1.077x
\item
  \textbf{B}: 5

  \textbf{nprobe}: 400
\item
  \textbf{default R}: 0.98752

  \textbf{candidate R}: 0.97304
\item
  \textbf{default 1/Ratio}: 0.999974940

  \textbf{candidate 1/Ratio}: 0.999903505
\item
  \textbf{QPS ratio}: 1.043x
\end{itemize}

Neither candidate reaches its default \texttt{nprobe=200}
\texttt{1/Ratio} target anywhere on the measured curve. At
\texttt{nprobe=200}, 97.1\% of B=4 queries and 87.3\% of B=5 queries are
worse. Therefore \texttt{1/Ratio} is not merely accepting every faster,
lower-Recall plan: both

DEEP negative decisions remain negative.
\end{frame}

\begin{frame}[fragile,allowframebreaks]{58. Attempt 3 Decision}
\protect\phantomsection\label{58-attempt-3-decision}
\begin{itemize}
\tightlist
\item
  \textbf{Predeclared condition}: at least one previous conclusion
  changes

  \textbf{Outcome}: pass: GIST historical reference
\item
  \textbf{Predeclared condition}: change survives a complete measured
  curve

  \textbf{Outcome}: pass: 11-point GIST union grid
\item
  \textbf{Predeclared condition}: positive result appears on two
  datasets or independent baselines

  \textbf{Outcome}: fail: DEEP provides controls, not a second positive
\item
  \textbf{Predeclared condition}: query-level tails show no hidden
  material loss

  \textbf{Outcome}: mixed: 39.6\% of paired GIST queries worsen
\item
  \textbf{Predeclared condition}: a mechanism-level gap remains after
  related work

  \textbf{Outcome}: fail: no mechanism beyond the old fac-error plan
\end{itemize}

\begin{Shaded}
\begin{Highlighting}[]
\NormalTok{Decision: CLOSE\_AS\_METRIC\_SENSITIVITY\_EVIDENCE}
\end{Highlighting}
\end{Shaded}

What survives:

\begin{enumerate}
\tightlist
\item
  report Recall and \texttt{1/Ratio} together in future ANN studies;
\item
  distinguish identifier-boundary loss from geometric degradation; and
\item
  retain GIST as a concrete metric-sensitive SAQ case and DEEP as
  controls.
\end{enumerate}

What does not reopen:

\begin{Shaded}
\begin{Highlighting}[]
\NormalTok{fac{-}error plan sweeps}
\NormalTok{query{-}fitted metric{-}aware planning}
\NormalTok{lossy projection or exact{-}hist as methods}
\NormalTok{high{-}overhead CAQ, mixed{-}plan, graph{-}prefix, or search{-}bound directions}
\end{Highlighting}
\end{Shaded}

The metric paper owns the general contribution. Our evidence is a
bounded SAQ case study, not an independent database-systems method.
\end{frame}

\begin{frame}[fragile,allowframebreaks]{59. Attempt 4: Research
Question, Old Stop, And V2 Boundary}
\protect\phantomsection\label{59-attempt-4-research-question-old-stop-and-v2-boundary}
Current R12 reviewed snapshot:
\texttt{saq-arbitrary-cardinality-feasibility-v2@3577edd}

Authoritative project-stop decision:
\texttt{saq-arbitrary-cardinality-feasibility-v2@3577edd}:
\texttt{docs/saq\_a4\_v2\_project\_stop\_no\_further\_artifact\_recovery\_2026\_07\_18.md}

START-only crash-record protocol target/direct-child review:
\texttt{saq-arbitrary-cardinality-feasibility-v2@1982865/@8bec546}

Fresh R1 PREP authorization target/direct-child review:
\texttt{saq-arbitrary-cardinality-feasibility-v2@cd1757e/@eb1b893}

Source-history R1 correction target/direct-child review:
\texttt{saq-arbitrary-cardinality-feasibility-v2@e9b5c83/@a1198c4}

Source-history R1 path-closure erratum target/direct-child review:
\texttt{saq-arbitrary-cardinality-feasibility-v2@150e5b3/@7749491}

Source-history epoch-correction protocol target/direct-child review:
\texttt{saq-arbitrary-cardinality-feasibility-v2@dcaed57/@e98a3e4}

HOST-I-R1 correction target/direct-child review:
\texttt{saq-arbitrary-cardinality-feasibility-v2@e17f887/@e10bde7}

HOST-I-R1 correction-only repair protocol target/direct-child review:
\texttt{saq-arbitrary-cardinality-feasibility-v2@ddfef99/@b1a7429}

HOST-I source-static target/direct-child review:
\texttt{saq-arbitrary-cardinality-feasibility-v2@4602585/@e8e9c79}

Source-authority erratum target/review:
\texttt{saq-arbitrary-cardinality-feasibility-v2@6fe8544/@9fa9528}

Independent HOST-I-AUTH review record:
\texttt{saq-arbitrary-cardinality-feasibility-v2@71e6bec}

Preserved host-rebind erratum target/review:
\texttt{saq-arbitrary-cardinality-feasibility-v2@5a47fed/@b89dabe}

Preserved PREP-authorization target/review:
\texttt{saq-arbitrary-cardinality-feasibility-v2@e7f940e/@16a8201}

Preserved CACHE-I target/review:
\texttt{saq-arbitrary-cardinality-feasibility-v2@c33a2bf/@5db3025}

Preserved CACHE-P target/review:
\texttt{saq-arbitrary-cardinality-feasibility-v2@56210f8/@249d5b8}

Preserved source-repair/review snapshots:
\texttt{saq-arbitrary-cardinality-feasibility-v2@e48df452/@c401dae}

Preserved terminal PAR result/review:
\texttt{saq-arbitrary-cardinality-feasibility-v2@30dfada/@fd5367e}

Preserved terminal A4-1S snapshot:
\texttt{saq-arbitrary-cardinality-analysis@f1b464b}

\begin{itemize}
\tightlist
\item
  \textbf{Stage}: A4-0 synthetic instrument

  \textbf{Committed status}: \texttt{PASS\_INSTRUMENT\_ONLY}
\item
  \textbf{What it means}: formulation, tiny witness, and reference
  checks passed
\item
  \textbf{Stage}: A4-1P base-only protocol

  \textbf{Committed status}:
  \texttt{FROZEN\_NOT\_AUTHORIZED\ /\ NOT\_RUN}
\item
  \textbf{What it means}: scientific inputs and decision rule were
  registered; base read remained forbidden
\item
  \textbf{Stage}: A4-1S synthetic implementation

  \textbf{Committed status}: \texttt{NO\_GO\_EXACT\_SOLVER\_COST}
\item
  \textbf{What it means}: parity passed, but the reviewed frozen
  construction/evidence pipeline exceeded its cost ceiling
\item
  \textbf{Stage}: A4-V2 parent protocol

  \textbf{Committed status}:
  \texttt{PROTOCOL\_READY\_NOT\_AUTHORIZED\_FOR\_EXECUTION}
\item
  \textbf{What it means}: a new FOM and evidence boundary were
  independently reviewed at \texttt{f86a51d}
\item
  \textbf{Stage}: A4-V2-P-ERRATUM

  \textbf{Committed status}:
  \texttt{PROTOCOL\_ERRATUM\_INDEPENDENT\_REVIEW\_PASS}
\item
  \textbf{What it means}: an additive finite PAR-report closure fixes
  terminal-P receipt recursion without changing the parent scientific
  contract
\item
  \textbf{Stage}: A4-V2-I

  \textbf{Committed status}:
  \texttt{SOURCE\_IMPLEMENTED\_STATIC\_REVIEW\_PASS}
\item
  \textbf{What it means}: corrected exact source \texttt{482c401},
  binding, and 35-file manifest passed three independent static-only
  reviews; failed parent \texttt{3a4f7c5} remains non-evidence
\item
  \textbf{Stage}: A4-V2-PAR

  \textbf{Committed status}:
  \texttt{ARTIFACT\_INVALID\ /\ REVIEWED\_TERMINAL}
\item
  \textbf{What it means}: the one authorized invocation rejected the
  \texttt{/bin/python} symlink before a B/P worker; no build, parity
  artifact, valid PAR authority, or scientific decision exists
\item
  \textbf{Stage}: A4-V2-I-R1

  \textbf{Committed status}:
  \texttt{SOURCE\_REPAIR\_STATIC\_REVIEW\_PASS}
\item
  \textbf{What it means}: exact source \texttt{e48df452} binds leader
  identity to stable \texttt{/proc/self/exe} bytes and same-inode
  \texttt{sys.executable}; three-track review passed, but no code was
  imported, built, or executed
\item
  \textbf{Stage}: A4-V2-CACHE-P target

  \textbf{Committed status}:
  \texttt{EXACT\_TARGET\_INDEPENDENT\_REVIEW\_PASS}
\item
  \textbf{What it means}: exact target \texttt{56210f8} and review
  \texttt{249d5b8} freeze
  \texttt{GO\_SANITIZED\_REMOTE\_ISOLATION\_PROTOCOL\_REQUIRED\ /\ NO\_GO\_DIRECT\_PAR\_R1};
  governance only, with quarantine unread and 35-source tree unchanged
\item
  \textbf{Stage}: A4-V2-CACHE-I

  \textbf{Committed status}:
  \texttt{GENERIC\_CACHE\_POLICY\_SOURCE\_STATIC\_REVIEW\_PASS}
\item
  \textbf{What it means}: target \texttt{c33a2bf} and review
  \texttt{5db3025} bind 37 filesystem sources, 38 executable units, and
  source-tree SHA-256 \texttt{95680075...}; static artifact governance
  only, with no PREP or parity evidence
\item
  \textbf{Stage}: A4-V2-CACHE-PREP-AUTH

  \textbf{Committed status}:
  \texttt{PREP\_AUTHORIZATION\_EXACT\_TARGET\_REVIEW\_PASS}
\item
  \textbf{What it means}: exact three-path target \texttt{e7f940e} and
  direct-child review \texttt{16a8201} bound a dormant PREP contract;
  the later host mismatch made its invocation authority nontransferable
  and its unique probe superseded
\item
  \textbf{Stage}: A4-V2-PREP-HOST-P

  \textbf{Committed status}:
  \texttt{HOST\_IDENTITY\_REBIND\_PROTOCOL\_REVIEW\_PASS}
\item
  \textbf{What it means}: target/review \texttt{5a47fed/@b89dabe} bind
  the current \texttt{.el9\_8.2} leader and complete future rebind
  closure; host identity is not rebound and no Python, PREP, probe,
  clone, token, build, data, or scientific action ran
\item
  \textbf{Stage}: A4-V2-PREP-HOST-I-AUTH

  \textbf{Committed status}:
  \texttt{AUTHORIZATION\_EXACT\_TARGET\_REVIEW\_PASS}
\item
  \textbf{What it means}: exact target/review \texttt{212a67b/@71e6bec}
  freeze only the contract and projection for a possible later source
  target and grant no HOST-I authority; no
  five-source/seven-derived-object rebind, Python, build, PREP, data, or
  scientific action occurred
\item
  \textbf{Stage}: attempted A4-V2-PREP-HOST-I

  \textbf{Committed status}: \texttt{STOPPED\_BEFORE\_TARGET}
\item
  \textbf{What it means}: exact checking exposed an unsatisfiable
  seven-versus-eight protocol-component closure before any source edit,
  Python, or PREP
\item
  \textbf{Stage}: A4-V2-PREP-HOST-I-ERRATUM

  \textbf{Committed status}:
  \texttt{HOST\_I\_SOURCE\_AUTHORITY\_ERRATUM\_REVIEW\_PASS}
\item
  \textbf{What it means}: target/review \texttt{6fe8544/@9fa9528} admit
  the eighth component and freeze its exact accounting/projection only;
  no implementation-source/derived-authority rebind or execution
  occurred
\item
  \textbf{Stage}: A4-V2-PREP-HOST-I

  \textbf{Committed status}:
  \texttt{SOURCE\_STATIC\_TARGET\_REVIEW\_FAIL\_AUTHORITY\_IDENTITY\_MISMATCH}
\item
  \textbf{What it means}: target \texttt{4602585} changed exactly the
  registered 14 paths without execution; direct-child review
  \texttt{e8e9c79} found one HIGH because the binding and crosswalk cite
  a false governing-review SHA-256; that target remains immutable failed
  evidence and is not retroactively passed

  by R1
\item
  \textbf{Stage}: A4-V2-PREP-HOST-I-R1-P

  \textbf{Committed status}:
  \texttt{HOST\_I\_R1\_CORRECTION\_ONLY\_REPAIR\_PROTOCOL\_REVIEW\_PASS}
\item
  \textbf{What it means}: target/review \texttt{ddfef99/@b1a7429} freeze
  only a future exact five-path correction target, corrected identities,
  cascade boundary, and direct-child review projection; zero
  LOW-or-higher findings, no repair authority, and no execution
\item
  \textbf{Stage}: A4-V2-PREP-HOST-I-R1

  \textbf{Committed status}:
  \texttt{HOST\_I\_R1\_CORRECTION\_ONLY\_REPAIR\_REVIEW\_PASS}
\item
  \textbf{What it means}: target/review \texttt{e17f887/@e10bde7}
  changed only the exact five paths, reproduced all frozen identities
  and 13/13 registered host objects with zero LOW-or-higher findings,
  and establishes \texttt{HOST\_IDENTITY\_REBOUND} only for that
  registered bundle; whole-host identity remains unestablished and no
  execution ran
\item
  \textbf{Stage}: A4-V2-SOURCE-HISTORY-P

  \textbf{Committed status}:
  \texttt{SOURCE\_HISTORY\_EPOCH\_CORRECTION\_PROTOCOL\_REVIEW\_PASS}
\item
  \textbf{What it means}: target/review \texttt{dcaed57/@e98a3e4}
  independently reproduced the exact five legacy blobs, other 32
  preserved sources, and strict four-legacy/six-active role partition;
  it rejects either-epoch/adaptive matching and authorizes no source
  edit or execution
\item
  \textbf{Stage}: first A4-V2-SOURCE-HISTORY-I

  \textbf{Committed status}:
  \texttt{STOPPED\_PRE\_TARGET\_PATH\_CLOSURE\_UNSATISFIABLE}
\item
  \textbf{What it means}: exact projection proved the required ten-path
  target impossible because the runtime schema and maximal witness had
  no authorized semantic, output, observation, or ledger delta; no
  source target or execution occurred
\item
  \textbf{Stage}: A4-V2-SOURCE-HISTORY-I-R1-P

  \textbf{Committed status}:
  \texttt{SOURCE\_HISTORY\_I\_R1\_PATH\_CLOSURE\_ERRATUM\_REVIEW\_PASS}
\item
  \textbf{What it means}: target/review \texttt{150e5b3/@7749491} froze
  the necessary-and-sufficient future eight-path/five-derived-object
  cascade, preserved the two runtime artifacts exactly, and retained the
  original contract as the sole ninth review component; zero
  LOW-or-higher findings and governance only
\item
  \textbf{Stage}: A4-V2-SOURCE-HISTORY-I-R1

  \textbf{Committed status}:
  \texttt{SOURCE\_HISTORY\_EPOCH\_CORRECTION\_SOURCE\_STATIC\_REVIEW\_PASS}
\item
  \textbf{What it means}: target/review \texttt{e9b5c83/@a1198c4}
  changed the exact eight paths and only one verifier source
  (\texttt{+38} net lines), independently reproduced five legacy
  overrides, the strict four-legacy/six-active role partition,
  nine-component boundary, and canonical 37-source tree, and preserved
  the runtime schema/maximal

  witness byte-for-byte; zero LOW-or-higher findings, but no
  syntax/import, verifier, build, PREP, data, or scientific execution
\item
  \textbf{Stage}: A4-V2-CACHE-PREP-R1-AUTH

  \textbf{Committed status}:
  \texttt{PREP\_AUTHORIZATION\_EXACT\_TARGET\_REVIEW\_PASS}
\item
  \textbf{What it means}: target/review \texttt{cd1757e/@eb1b893}
  independently closed the reviewed source head, current
  source/authority identities, corrected source history, and
  registered-host ceiling with zero LOW-or-higher findings;
  documentation governance only, with the ordinary target-head
  observation distinct from the reserved immediate review-head probe
\item
  \textbf{Stage}: A4-V2-CACHE-PREP-ISO-R1

  \textbf{Committed status}:
  \texttt{CRASH\_OR\_UNKNOWN\ /\ START\_ONLY\ /\ NO\_SCIENTIFIC\_DECISION}
\item
  \textbf{What it means}: the exact immediate probe passed and the
  unique invocation was consumed, but only durable START and empty
  captures exist; no TERMINAL, isolation root, clone, receipt, or
  PREP-body entry
\item
  \textbf{Stage}: A4-V2-CACHE-PREP-ISO-R1-CRASH-P

  \textbf{Committed status}:
  \texttt{START\_ONLY\_CRASH\_RECORD\_PROTOCOL\_REVIEW\_PASS}
\item
  \textbf{What it means}: target/review \texttt{1982865/@8bec546} freeze
  only the future reviewed terminal-record topology, with zero
  LOW-or-higher findings and no retry or repair authority
\item
  \textbf{Stage}: A4-V2-CACHE-PREP-ISO-R1-CRASH-I / CACHE-BIND

  \textbf{Committed status}: \texttt{NOT\_AUTHORIZED\ /\ NOT\_RUN}
\item
  \textbf{What it means}: no terminal record or binding was formed;
  portfolio closure does not reclassify START-only runtime state
\item
  \textbf{Stage}: A4-V2-PAR-R1

  \textbf{Committed status}: \texttt{NOT\_AUTHORIZED\ /\ NOT\_RUN}
\item
  \textbf{What it means}: no build or corrected parity event occurred
\item
  \textbf{Stage}: A4-V2-SRUN / data

  \textbf{Committed status}: \texttt{NOT\_AUTHORIZED\ /\ NOT\_RUN}
\item
  \textbf{What it means}: no synthetic scientific event or
  benchmark/base/query/index read occurred
\item
  \textbf{Stage}: A4-V2 project portfolio

  \textbf{Committed status}:
  \texttt{CLOSED\ /\ STOP\_NO\_FURTHER\_ARTIFACT\_RECOVERY}
\item
  \textbf{What it means}: audited closure \texttt{3577edd} stops
  investment in the recovery chain; it is not an artifact terminal,
  scientific result, or performance result
\item
  \textbf{Stage}: SAQ/index/search integration

  \textbf{Committed status}: \texttt{NOT\_AUTHORIZED}
\item
  \textbf{What it means}: no production representation or query-path
  change may be made
\end{itemize}

Research question:

\begin{Shaded}
\begin{Highlighting}[]
\NormalTok{Under one fixed B\_g{-}bit word for a two{-}factor group, can positive{-}integer}
\NormalTok{scalar cardinalities use the joint{-}state budget more effectively than}
\NormalTok{power{-}of{-}two cardinalities while retaining one random{-}access word and one}
\NormalTok{expanded{-}table lookup per group?}
\end{Highlighting}
\end{Shaded}

The competing factorized feasible sets are:

\begin{Shaded}
\begin{Highlighting}[]
\NormalTok{dyadic:     K\_j in \{1, 2, 4, 8, ...\}}
\NormalTok{arbitrary:  K\_j in positive integers}
\NormalTok{both:       product\_j K\_j \textless{}= S = 2\^{}B\_g}
\end{Highlighting}
\end{Shaded}

This is a fixed-address coding opportunity study. It does not modify CAQ
or SAQ, and it produced no natural-data, Recall, QPS, or systems claim.
\end{frame}

\begin{frame}[fragile,allowframebreaks]{60. Correct Fixed-Rate
Formulation And Witness}
\protect\phantomsection\label{60-correct-fixed-rate-formulation-and-witness}
First correct a tempting but invalid rate comparison:

\begin{Shaded}
\begin{Highlighting}[]
\NormalTok{(3,3) has 9 joint states and needs 4 fixed bits.}
\NormalTok{It cannot be compared at equal fixed rate with a 3{-}bit (4,2) product.}
\end{Highlighting}
\end{Shaded}

The frozen equal-rate witness instead uses:

\begin{itemize}
\tightlist
\item
  \textbf{Quantity}: word width

  \textbf{Value}: \texttt{B\_g=4} bits
\item
  \textbf{Quantity}: joint capacity

  \textbf{Value}: \texttt{S=16} addresses
\item
  \textbf{Quantity}: best arbitrary witness allocation

  \textbf{Value}: \texttt{(3,5)}, 15 used states
\item
  \textbf{Quantity}: selected dyadic allocation

  \textbf{Value}: \texttt{(4,4)}, 16 used states
\end{itemize}

For two factors with zero-based labels \texttt{z\_1} and \texttt{z\_2},
mixed-radix addressing is direct:

\begin{Shaded}
\begin{Highlighting}[]
\NormalTok{u = z\_1 + K\_1 z\_2,        0 \textless{}= u \textless{} K\_1 K\_2 \textless{}= S.}
\end{Highlighting}
\end{Shaded}

The address is still one fixed word and the query interface is still one
\texttt{S}-entry lookup per group. Unused addresses remain paid
capacity; entropy or Huffman expected length is not a fixed-payload
saving.

The only unconditional optimization statement is:

\begin{Shaded}
\begin{Highlighting}[]
\NormalTok{D*\_arbitrary \textless{}= D*\_dyadic}
\end{Highlighting}
\end{Shaded}

for the same exact factorized reconstruction-SSE objective, because the
arbitrary-cardinality feasible set contains the dyadic set. It is not a
Recall, QPS, or block-VQ dominance theorem.
\end{frame}

\begin{frame}[fragile,allowframebreaks]{61. Prior Art And Claim
Boundary}
\protect\phantomsection\label{61-prior-art-and-claim-boundary}
The following are known ingredients, not Attempt 4 novelty:

\begin{Shaded}
\begin{Highlighting}[]
\NormalTok{non{-}power{-}of{-}two scalar quantization}
\NormalTok{transform coding and bit allocation}
\NormalTok{adaptive or irregular product quantization}
\NormalTok{mixed{-}radix fixed addressing}
\NormalTok{finite scalar quantization and non{-}dyadic palette constructions}
\NormalTok{entropy{-}constrained and variable{-}length coding}
\end{Highlighting}
\end{Shaded}

Therefore a publishable claim would need evidence beyond the primitive:

\begin{enumerate}
\tightlist
\item
  a material gap on natural ANN residuals at matched fixed payload;
\item
  a same-capacity trained block-VQ control, because unrestricted block
  VQ weakly dominates a factorized codebook in reconstruction
  expressiveness;
\item
  complete training, metadata, expanded-table, packing, and scan-cost
  accounting; and
\item
  only after a new frozen systems protocol, a replicated
  Recall-throughput effect against competitive baselines.
\end{enumerate}

Current boundary:

\begin{Shaded}
\begin{Highlighting}[]
\NormalTok{PRESERVE A4{-}1S TERMINAL COST STOP}
\NormalTok{A4 V2 CORRECTED PROTOCOL AUTHORITY REVIEWED}
\NormalTok{A4{-}V2{-}I SOURCE\_IMPLEMENTED\_STATIC\_REVIEW\_PASS}
\NormalTok{A4{-}V2{-}PAR ARTIFACT\_INVALID / REVIEWED TERMINAL / NO VALID PAR AUTHORITY}
\NormalTok{A4{-}V2{-}I{-}R1 SOURCE\_REPAIR\_STATIC\_REVIEW\_PASS ONLY}
\NormalTok{A4{-}V2{-}CACHE{-}P TARGET EXACTLY REVIEWED / NO DIRECT PAR{-}R1}
\NormalTok{A4{-}V2{-}CACHE{-}I GENERIC\_CACHE\_POLICY\_SOURCE\_STATIC\_REVIEW\_PASS ONLY}
\end{Highlighting}
\end{Shaded}

\begin{Shaded}
\begin{Highlighting}[]
\NormalTok{A4{-}V2{-}CACHE{-}PREP{-}AUTH PREP\_AUTHORIZATION\_EXACT\_TARGET\_REVIEW\_PASS ONLY}
\NormalTok{A4{-}V2{-}PREP{-}HOST{-}P HOST\_IDENTITY\_REBIND\_PROTOCOL\_REVIEW\_PASS ONLY}
\NormalTok{A4{-}V2{-}PREP{-}HOST{-}I{-}AUTH AUTHORIZATION\_EXACT\_TARGET\_REVIEW\_PASS ONLY}
\NormalTok{A4{-}V2{-}PREP{-}HOST{-}I ATTEMPT STOPPED BEFORE TARGET ON 7{-}VS{-}8 CONTRADICTION}
\NormalTok{A4{-}V2{-}PREP{-}HOST{-}I{-}ERRATUM HOST\_I\_SOURCE\_AUTHORITY\_ERRATUM\_REVIEW\_PASS ONLY}
\NormalTok{A4{-}V2{-}PREP{-}HOST{-}I SOURCE\_STATIC\_TARGET\_REVIEW\_FAIL\_AUTHORITY\_IDENTITY\_MISMATCH}
\NormalTok{A4{-}V2{-}PREP{-}HOST{-}I{-}R1{-}P HOST\_I\_R1\_CORRECTION\_ONLY\_REPAIR\_PROTOCOL\_REVIEW\_PASS ONLY}
\end{Highlighting}
\end{Shaded}

\begin{Shaded}
\begin{Highlighting}[]
\NormalTok{A4{-}V2{-}PREP{-}HOST{-}I{-}R1 HOST\_I\_R1\_CORRECTION\_ONLY\_REPAIR\_REVIEW\_PASS}
\NormalTok{HOST\_IDENTITY\_REBOUND ESTABLISHED FOR REGISTERED 13{-}OBJECT BUNDLE ONLY}
\NormalTok{WHOLE{-}HOST IDENTITY NOT ESTABLISHED}
\NormalTok{A4{-}V2{-}SOURCE{-}HISTORY{-}P SOURCE\_HISTORY\_EPOCH\_CORRECTION\_PROTOCOL\_REVIEW\_PASS}
\NormalTok{A4{-}V2{-}SOURCE{-}HISTORY{-}I STOPPED PRE{-}TARGET; TEN{-}PATH CLOSURE UNSATISFIABLE}
\NormalTok{A4{-}V2{-}SOURCE{-}HISTORY{-}I{-}R1{-}P SOURCE\_HISTORY\_I\_R1\_PATH\_CLOSURE\_ERRATUM\_REVIEW\_PASS}
\NormalTok{A4{-}V2{-}SOURCE{-}HISTORY{-}I{-}R1 SOURCE\_HISTORY\_EPOCH\_CORRECTION\_SOURCE\_STATIC\_REVIEW\_PASS}
\end{Highlighting}
\end{Shaded}

\begin{Shaded}
\begin{Highlighting}[]
\NormalTok{KNOWN SOURCE{-}HISTORY RULE CORRECTED IN STATIC SOURCE / VERIFIER NOT RUN}
\NormalTok{A4{-}V2{-}CACHE{-}PREP{-}R1{-}AUTH PREP\_AUTHORIZATION\_EXACT\_TARGET\_REVIEW\_PASS ONLY}
\NormalTok{CACHE{-}VERIFIER AND PAR READINESS NOT ESTABLISHED}
\NormalTok{A4{-}V2{-}CACHE{-}PREP{-}ISO{-}R1 CRASH\_OR\_UNKNOWN / START\_ONLY / NO SCIENTIFIC DECISION}
\NormalTok{A4{-}V2{-}CACHE{-}PREP{-}ISO{-}R1{-}CRASH{-}P START\_ONLY\_CRASH\_RECORD\_PROTOCOL\_REVIEW\_PASS}
\NormalTok{A4{-}V2{-}CACHE{-}PREP{-}ISO{-}R1{-}CRASH{-}I, REPAIR, RETRY, AND CACHE{-}BIND NOT AUTHORIZED / NOT RUN}
\NormalTok{A4{-}V2{-}PAR{-}R1 NOT AUTHORIZED / NOT RUN}
\end{Highlighting}
\end{Shaded}

\begin{Shaded}
\begin{Highlighting}[]
\NormalTok{A4{-}V2{-}SRUN, DATA, AND SAQ CHANGES NOT AUTHORIZED / NOT RUN}
\NormalTok{A4{-}V2 PORTFOLIO CLOSED / STOP\_NO\_FURTHER\_ARTIFACT\_RECOVERY}
\NormalTok{NO ARTIFACT TERMINAL / NO SCIENTIFIC DECISION / PERFORMANCE\_NOT\_YET\_MEASURED}
\end{Highlighting}
\end{Shaded}

The fixed-word opportunity remains a mathematical possibility, but the
frozen exact A4-1S construction/evidence pipeline did not meet its
preregistered cost ceiling. V2 does not rescue or reinterpret that
result. It freezes a new synthetic comparative-instrument question and
cannot

claim a new quantization primitive, feasibility, or SAQ integration. The
erratum repairs protocol closure only. The later PAR attempt reached
module import but failed at the leader executable\textquotesingle s
no-follow identity read, before build or parity. It adds only

artifact-validation failure evidence, not evidence about the candidate.
I-R1 then repaired that dedicated identity acquisition and passed exact-
commit static review, but no corrected conductor was executed. It
therefore adds source-correctness evidence only and does not supersede
the

PAR terminal status. CACHE-P later reviewed CPython cache semantics and
exact committed source, retained the unread quarantine, and froze a
staged sanitized remote- isolation authority DAG. Its exact-target
review found no LOW+ issue, but its claim ceiling is

artifact governance only. It changes no implementation or environment,
creates no clone, and rejects direct PAR-R1.

CACHE-I then added only the frozen generic cache-policy source/schema
layer. Exact target \texttt{c33a2bf} and direct-child review
\texttt{5db3025} close 37 filesystem sources plus one separate inline
bootstrap, four schema/maximal-instance pairs, and a 38-unit static
import/process/mutation inventory. A separately authorized

Python 3.9 \texttt{compile()}-only check accepted the exact locked
source snapshots without importing or executing them. The maximum claim
is \texttt{GENERIC\_CACHE\_POLICY\_SOURCE\_STATIC\_REVIEW\_PASS}, not
PREP readiness, parity, feasibility, an SAQ limitation, systems
performance, novelty, or a method.

The subsequent exact three-path PREP authorization target
\texttt{e7f940e} and direct-child review \texttt{16a8201} reached
\texttt{PREP\_AUTHORIZATION\_EXACT\_TARGET\_REVIEW\_PASS} with 0 LOW+
findings. This froze only a dormant execution contract. Before
invocation, static checking found that a root RPM update had replaced
the

frozen \texttt{.el9\_8} Python leader with \texttt{.el9\_8.2}; PREP was
not invoked, START was not created, and no runtime terminal status
exists. Host-rebind erratum target/review \texttt{5a47fed/@b89dabe} then
reached only \texttt{HOST\_IDENTITY\_REBIND\_PROTOCOL\_REVIEW\_PASS}. It
does not rebind source or establish PREP

readiness. The old invocation authority is unspent but nontransferable;
its probe is unspent but superseded and must never run or be reused. The
later exact HOST-I-AUTH target/review \texttt{212a67b/@71e6bec} reached
only \texttt{AUTHORIZATION\_EXACT\_TARGET\_REVIEW\_PASS}. It froze a
possible future fourteen-path

source edge but changed none of the five sources or seven derived
authority objects. When HOST-I was later instructed, exact projection
checking found that the frozen seven-component authority closure could
not admit the required eighth host-rebind component, so

work stopped before any source target. Additive source-authority erratum
target/review \texttt{6fe8544/@9fa9528} reached only
\texttt{HOST\_I\_SOURCE\_AUTHORITY\_ERRATUM\_REVIEW\_PASS}: it admits
that eighth component and freezes exact future byte/accounting
semantics, but still changes no implementation source or derived
authority object. The subsequently authorized

HOST-I produced exact source-static target \texttt{4602585}. Its
direct-child review \texttt{e8e9c79} found one HIGH: the binding and
crosswalk contain a false digest for the governing erratum review.
Therefore the status is
\texttt{SOURCE\_STATIC\_TARGET\_REVIEW\_FAIL\_AUTHORITY\_IDENTITY\_MISMATCH},
and Python/PREP remain \texttt{NOT\_AUTHORIZED\_NOT\_RUN}.

Correction-only protocol target/review \texttt{ddfef99/@b1a7429} later
reached
\texttt{HOST\_I\_R1\_CORRECTION\_ONLY\_REPAIR\_PROTOCOL\_REVIEW\_PASS}.
Exact R1 target/review \texttt{e17f887/@e10bde7} then reached
\texttt{HOST\_I\_R1\_CORRECTION\_ONLY\_REPAIR\_REVIEW\_PASS} and
established rebound only for the registered 13-object bundle; whole-host
identity remains unestablished. Source-history protocol target/review
\texttt{dcaed57/@e98a3e4} subsequently reached
\texttt{SOURCE\_HISTORY\_EPOCH\_CORRECTION\_PROTOCOL\_REVIEW\_PASS},
freezing an

exact role-preassigned two-epoch rule. The first implementation
authorization then stopped before target because its ten-path closure
was unsatisfiable. Path-closure erratum target/review
\texttt{150e5b3/@7749491} subsequently reached
\texttt{SOURCE\_HISTORY\_I\_R1\_PATH\_CLOSURE\_ERRATUM\_REVIEW\_PASS},
freezing only an exact eight-path/five-derived-object correction
boundary. Exact R1 target/review \texttt{e9b5c83/@a1198c4} subsequently

reached
\texttt{SOURCE\_HISTORY\_EPOCH\_CORRECTION\_SOURCE\_STATIC\_REVIEW\_PASS}
with zero LOW-or-higher findings. It corrected only the known verifier
source rule and five derived authority objects while preserving the two
runtime artifacts; the verifier was not run, so all cache-verifier/PAR
readiness claims remain forbidden. Fresh

R1 PREP authorization target/review \texttt{cd1757e/@eb1b893}
subsequently passed exact documentation review. Its later unique probe
and invocation were consumed and stopped at START-only
\texttt{CRASH\_OR\_UNKNOWN} before the PREP body. Reviewed protocol
\texttt{1982865/@8bec546} closes only the future terminal-record
topology; CRASH-I, repair,

retry, and CACHE-BIND remain unauthorized.
\end{frame}

\begin{frame}[fragile,allowframebreaks]{62. A4-0 Synthetic Instrument
Result}
\protect\phantomsection\label{62-a4-0-synthetic-instrument-result}
\begin{itemize}
\tightlist
\item
  \textbf{Quantity}: fixed capacity

  \textbf{Observed}: 16
\item
  \textbf{Quantity}: arbitrary cardinalities / used states

  \textbf{Observed}: \texttt{(3,5)} / 15
\item
  \textbf{Quantity}: arbitrary distortion per vector

  \textbf{Observed}: 0
\item
  \textbf{Quantity}: dyadic cardinalities / used states

  \textbf{Observed}: \texttt{(4,4)} / 16
\item
  \textbf{Quantity}: dyadic distortion per vector

  \textbf{Observed}: 0.1
\item
  \textbf{Quantity}: mixed-radix LUT/direct L2 maximum difference

  \textbf{Observed}: 0
\item
  \textbf{Quantity}: unrestricted 16-codeword block-VQ distortion

  \textbf{Observed}: 0
\end{itemize}

The constructed feasible-set gap is recovered exactly:

\begin{Shaded}
\begin{Highlighting}[]
\NormalTok{D*\_dyadic {-} D*\_arbitrary = 0.1 per vector.}
\end{Highlighting}
\end{Shaded}

Seven deterministic tests pass: exact scalar DP, duplicate/weighted
support, both allocation modes against exhaustive enumeration, all 15
mixed-radix tuples, expanded-LUT/direct-distance parity, and the frozen
witness/Huffman values.

Interpretation ceiling:

\begin{Shaded}
\begin{Highlighting}[]
\NormalTok{PASS\_INSTRUMENT\_ONLY}
\end{Highlighting}
\end{Shaded}

The unrestricted block-VQ oracle also reaches zero. No dataset, query,
ground truth, PCA, IVF, or SAQ artifact was used. The result validates
the reference instrument and strict dyadic-feasible-set witness, not a
practical method.
\end{frame}

\begin{frame}[fragile,allowframebreaks]{63. Frozen A4-1 Base-Only
Falsification Gate}
\protect\phantomsection\label{63-frozen-a4-1-base-only-falsification-gate}
Status:
\texttt{FROZEN\_NOT\_AUTHORIZED\ /\ NOT\_RUN\ /\ FORECLOSED\_BY\_A4-1S}

Retained only for protocol provenance, A4-1 would have used
query-unaware GIST sample50k and CIFAR60K residual panels at fixed 4-bit
and 8-bit two-factor words. Its registered primary arms were:

\begin{Shaded}
\begin{Highlighting}[]
\NormalTok{dyadic\_word}
\NormalTok{arbitrary\_word}
\NormalTok{trained\_block\_vq          same joint capacity}
\end{Highlighting}
\end{Shaded}

A global capped dyadic pack is retained as an attribution control. It
has the same 32/64-byte database payload but a different
128-scalar-lookup interface, so it cannot silently replace the
matched-word comparison.

For items 1-\/-4 below, passage requires Holm-adjusted one-sided
evidence to reject the corresponding non-positive null in every
dataset-by-rate cell. Items 5-\/-6 are additional hard point-estimate
gates:

\begin{enumerate}
\tightlist
\item
  arbitrary words remove more than 5\% of held-out dyadic error;
\item
  a positive trained-block-VQ opportunity exists;
\item
  arbitrary words close more than half of the dyadic-to-block-VQ gap;
\item
  the base-pair absolute-distance-error contrast is positive;
\item
  at least 48 of 64 groups have positive dyadic-error removal; and
\item
  every leave-one-group-out pooled removal estimate remains at least
  5\%.
\end{enumerate}

The first four contrasts are frozen as a 16-hypothesis Holm family
across the four dataset-by-rate cells. Every primary arm pays identical
fixed payload and full expanded lookup capacity.

The registered maximum positive outcome would have been only:

\begin{Shaded}
\begin{Highlighting}[]
\NormalTok{GO\_TO\_SYSTEMS\_PREREG}
\end{Highlighting}
\end{Shaded}

That would have established a data-level word-local opportunity and
permitted writing a later systems protocol. It still would not have
established Recall, QPS, or a method claim. A4-1S did not pass, so
reading the registered base inputs remains

forbidden and this gate must not be opened.

A4 V2 does not reopen A4-1. Its source-only implementation and
independent static review completed, and the user separately authorized
frozen PAR. That one invocation stopped at \texttt{ARTIFACT\_INVALID}
before any B/P worker, build, or parity fixture, so it created

no valid PAR authority and no scientific decision. The later clean I-R1
source correction and independent static review are now complete.
CACHE-P subsequently reviewed the cache/staging question at exact target
\texttt{56210f8}, with its independent review recorded at

\texttt{249d5b8}. It requires sanitized remote isolation and rejects
direct PAR-R1, but changes no implementation or environment and leaves
the quarantine unread and uncleaned. CACHE-I subsequently reached
\texttt{GENERIC\_CACHE\_POLICY\_SOURCE\_STATIC\_REVIEW\_PASS} at
target/review \texttt{c33a2bf/@5db3025}; it adds only static governance
sources

and schemas. The later PREP authorization target/review
\texttt{e7f940e/@16a8201} reached only
\texttt{PREP\_AUTHORIZATION\_EXACT\_TARGET\_REVIEW\_PASS}. Static
checking then found the \texttt{.el9\_8\ -\textgreater{}\ .el9\_8.2}
leader replacement before PREP or START. Host-rebind erratum
target/review \texttt{5a47fed/@b89dabe} reached only
\texttt{HOST\_IDENTITY\_REBIND\_PROTOCOL\_REVIEW\_PASS}; the old probe
is superseded and must

never be reused. HOST-I-AUTH target/review \texttt{212a67b/@71e6bec}
then reached only \texttt{AUTHORIZATION\_EXACT\_TARGET\_REVIEW\_PASS}:
no source/authority rebind, Python, START, or PREP occurred. A later
HOST-I instruction stopped before any source target on the exact
seven-versus-eight component contradiction. Source-authority erratum
target/review \texttt{6fe8544/@9fa9528} reached

only \texttt{HOST\_I\_SOURCE\_AUTHORITY\_ERRATUM\_REVIEW\_PASS}; it
changes no implementation source or derived authority object and grants
no HOST-I authority. The later source-static target/review
\texttt{4602585/@e8e9c79} ended at
\texttt{SOURCE\_STATIC\_TARGET\_REVIEW\_FAIL\_AUTHORITY\_IDENTITY\_MISMATCH}:
one HIGH false governing-review digest in the binding and crosswalk,
with host

rebound not established. Correction-only protocol target/review
\texttt{ddfef99/@b1a7429} subsequently passed with zero LOW-or-higher
findings. Exact R1 target/review \texttt{e17f887/@e10bde7} then
corrected only that authority identity and passed independent review
across all 13 registered host objects. Bundle rebound is established;
whole-host

identity and execution readiness are not. Source-history protocol
target/review \texttt{dcaed57/@e98a3e4} passed and froze a strict future
two-epoch repair. The first authorized implementation stopped before
target on its unsatisfiable ten-path closure; path-closure erratum
target/review \texttt{150e5b3/@7749491} then passed only for

a future exact eight-path/five-derived-object correction. Exact R1
target/review \texttt{e9b5c83/@a1198c4} then implemented only that
static repair and passed direct-child review with zero LOW-or-higher
findings. The corrected verifier was not run, so cache-verifier/PAR
readiness remains unestablished. Fresh dormant PREP authorization

target/review \texttt{cd1757e/@eb1b893} then reached
\texttt{PREP\_AUTHORIZATION\_EXACT\_TARGET\_REVIEW\_PASS} with zero
LOW-or-higher findings. Its later unique immediate probe passed and
invocation was consumed, but only durable START and empty captures
exist; the frozen result is
\texttt{CRASH\_OR\_UNKNOWN\ /\ NO\_SCIENTIFIC\_DECISION}. Crash-record
protocol \texttt{1982865/@8bec546} passed independent

review. CRASH-I, repair, retry, CACHE-BIND, PAR-R1, SRUN, base, and
query reads remain unauthorized.
\end{frame}

\begin{frame}[fragile,allowframebreaks]{64. A4-1S Synthetic
Implementation And Cost Gate}
\protect\phantomsection\label{64-a4-1s-synthetic-implementation-and-cost-gate}
Status: \texttt{NO\_GO\_EXACT\_SOLVER\_COST} -\/- committed and
independently reviewed

A4-1S exists only to establish that the exact registered machinery is
correct, deterministic, representation-complete, and affordable before
any base read. Its committed parity phase executed:

\begin{Shaded}
\begin{Highlighting}[]
\NormalTok{1,044 exact scalar cases: 780 exhaustive + 8 bit{-}pattern + 256 PCG64}
\NormalTok{3,961 persisted scalar solutions and exact{-}rational replays}
\NormalTok{2,433,600 allocation decisions, all executed in canonical order and digested}
\NormalTok{explicit mixed{-}radix, B4/B8 packing, and LUT fixtures}
\NormalTok{64 deterministic trained{-}block{-}VQ cases plus five named microfixtures}
\NormalTok{six status{-}precedence fixtures and 22/22 deterministic tests}
\end{Highlighting}
\end{Shaded}

All registered parity checks passed. The six canonical evidence files
are at \texttt{335837e}; independent review at \texttt{988ace0} found
zero blocker or high-severity issues. This supports
\texttt{PASS\_PARITY} for the instrument only.

The frozen one-panel \texttt{8192\ x\ 128} cost command then reached its
registered early-stop terminal after publishing coordinate 60, leaving
61 complete scalar-coordinate shards. The reviewed integer accounting
is:

\begin{Shaded}
\begin{Highlighting}[]
\NormalTok{timed CPU                         34,805,155,525 us}
\NormalTok{frozen projection factor                     5 / 2}
\NormalTok{projected cost                   24.170246892361 CPU{-}hours}
\NormalTok{registered ceiling              24.000000000000 CPU{-}hours}
\NormalTok{excess                              612.8888125 seconds}
\end{Highlighting}
\end{Shaded}

Canonical shard serialization used 18,096.207493 CPU-seconds, versus
16,706.810033 seconds for exact scalar construction. The result
therefore rejects the whole frozen construction/evidence pipeline, not
the scalar solver alone and not arbitrary-cardinality quantization
quality.

Status precedence prevents a broken instrument from becoming scientific
evidence:

\begin{itemize}
\tightlist
\item
  \textbf{Condition}: artifact/schema or implementation parity defect

  \textbf{Outcome}: \texttt{ARTIFACT\_INVALID} or
  \texttt{IMPLEMENTATION\_INVALID}
\item
  \textbf{Condition}: trained block-control defect

  \textbf{Outcome}: \texttt{CONTROL\_INVALID}
\item
  \textbf{Condition}: selected scalar alphabet collapses at binary32

  \textbf{Outcome}: \texttt{NO\_GO\_REPRESENTATION}
\item
  \textbf{Condition}: frozen exact construction/evidence pipeline cost
  ceiling is exceeded

  \textbf{Outcome}: \texttt{NO\_GO\_EXACT\_SOLVER\_COST}
\item
  \textbf{Condition}: every correctness, representation, control, and
  cost gate passes

  \textbf{Outcome}: \texttt{PASS\_SYNTHETIC\_GATE\_ONLY}
\end{itemize}

\texttt{PASS\_SYNTHETIC\_GATE\_ONLY} would have permitted only asking
the user whether to authorize the already-registered base gate. It did
not occur.

Cost evidence is committed at \texttt{9ce1052}; the independent review
and terminal decision are committed at \texttt{f1b464b}. Allocation,
block-VQ, and encoding cost phases did not execute after the
scalar-prefix early stop, so the result makes no claim

about them. A4-1 was not run, and untracked worktree files are
deliberately excluded.
\end{frame}

\begin{frame}[fragile,allowframebreaks]{65. A4 V2: Terminal PAR Failure,
Source Repair, And Governance Chain}
\protect\phantomsection\label{65-a4-v2-terminal-par-failure-source-repair-and-governance-chain}
Primary review:
\texttt{saq-arbitrary-cardinality-feasibility-v2@c4ccea3}

Reviewed protocol snapshot:
\texttt{saq-arbitrary-cardinality-feasibility-v2@f86a51d}

Corrected protocol-authority snapshot:
\texttt{saq-arbitrary-cardinality-feasibility-v2@f13a383}

Independent erratum-review head:
\texttt{saq-arbitrary-cardinality-feasibility-v2@98e6999}

Corrected source snapshot:
\texttt{saq-arbitrary-cardinality-feasibility-v2@482c401}

Independent source-review head:
\texttt{saq-arbitrary-cardinality-feasibility-v2@4ce2e69}

Reviewed PAR authorization base:
\texttt{saq-arbitrary-cardinality-feasibility-v2@2e983a0}

Terminal PAR result snapshot:
\texttt{saq-arbitrary-cardinality-feasibility-v2@30dfada}

Independent terminal-review head:
\texttt{saq-arbitrary-cardinality-feasibility-v2@fd5367e}

I-R1 source-repair authorization:
\texttt{saq-arbitrary-cardinality-feasibility-v2@fe10bc3}

Exact source-repair snapshot:
\texttt{saq-arbitrary-cardinality-feasibility-v2@e48df452}

Independent source-repair review head:
\texttt{saq-arbitrary-cardinality-feasibility-v2@c401dae}

CACHE-P authorization:
\texttt{saq-arbitrary-cardinality-feasibility-v2@ffd0f41}

CACHE-P exact content target:
\texttt{saq-arbitrary-cardinality-feasibility-v2@56210f8}

CACHE-P independent review record:
\texttt{saq-arbitrary-cardinality-feasibility-v2@249d5b8}

CACHE-I authorization/review:
\texttt{saq-arbitrary-cardinality-feasibility-v2@4f38ca9/@187e363}

CACHE-I exact source/static target:
\texttt{saq-arbitrary-cardinality-feasibility-v2@c33a2bf}

CACHE-I independent review record:
\texttt{saq-arbitrary-cardinality-feasibility-v2@5db3025}

PREP authorization target:
\texttt{saq-arbitrary-cardinality-feasibility-v2@e7f940e}

PREP authorization independent review record:
\texttt{saq-arbitrary-cardinality-feasibility-v2@16a8201}

Host-rebind erratum target/review:
\texttt{saq-arbitrary-cardinality-feasibility-v2@5a47fed/@b89dabe}

HOST-I-AUTH target:
\texttt{saq-arbitrary-cardinality-feasibility-v2@212a67b}

HOST-I-AUTH independent review record:
\texttt{saq-arbitrary-cardinality-feasibility-v2@71e6bec}

Source-authority erratum target/review:
\texttt{saq-arbitrary-cardinality-feasibility-v2@6fe8544/@9fa9528}

HOST-I source-static target/review:
\texttt{saq-arbitrary-cardinality-feasibility-v2@4602585/@e8e9c79}

HOST-I-R1 correction-only protocol target/review:
\texttt{saq-arbitrary-cardinality-feasibility-v2@ddfef99/@b1a7429}

HOST-I-R1 correction target/review:
\texttt{saq-arbitrary-cardinality-feasibility-v2@e17f887/@e10bde7}

The review preserves the old A4-1S result and changes the next question,
not the result. A4-1S timed exact construction together with full
canonical research serialization; serialization alone used 18,096
CPU-seconds. That supports neither a solver lower bound nor

candidate-only deployment cost.

V2 therefore freezes one primary FOM for the complete four-arm synthetic
comparative instrument:

\begin{Shaded}
\begin{Highlighting}[]
\NormalTok{T\_instrument = C\_setup + C\_core + C\_bundle\_io}
\NormalTok{PASS iff T\_instrument \textless{}= 34,560,000,000 CPU microseconds}
\end{Highlighting}
\end{Shaded}

The cap is only an internal admission boundary for the one frozen
synthetic panel. The old \texttt{5/2} factor does not transfer: there is
no V2 projection to GIST, CIFAR, two datasets, a full vector, an SAQ
index build,

or deployment.

Initial source-only inspection found that the parent protocol asked the
terminal \texttt{P\_parity} receipt to include publication work whose
terminal values did not yet exist. The additive erratum closes only that
self-reference with a finite one-file \texttt{PAR\_report} after
terminal

P. It leaves the three parent blobs, scientific computation, FOM,
threshold, retry policy, and status precedence unchanged.

No work disappears from reporting:

\begin{Shaded}
\begin{Highlighting}[]
\NormalTok{B\_build + P\_parity + T\_instrument}
\NormalTok{        + E\_emit + V\_replay + E\_archive\_body = T\_study\_metered}
\NormalTok{PAR\_report is a finite one{-}file PAR authority closure after terminal P}
\NormalTok{F\_trailer is a finite three{-}file final{-}study closure after archive{-}body timing}
\NormalTok{memory, temporary bytes, bundle bytes, evidence/archive bytes: separately capped}
\end{Highlighting}
\end{Shaded}

Every scientific/parity/control operation and every file validation,
discovery, hash, tree-walk, and byte-ledger derivation required to form
the seal must finish before terminal P. The post-P closure may only
inject captured integers, perform frozen bounded arithmetic, encode one
fixed

15-key seal, and atomically publish it. Its exact conservative
schema-language maximum is 5,171 bytes including LF; that maximal
witness is syntactic, not a semantically admissible resource
observation.

The schema freezes 396 atomic construction units, exact-prefix evidence,
independent exact-rational scalar/allocation replay, deterministic
bit-level block/encoding/packing replay, raw prelaunch evidence, and
admissibility-first status precedence. The four-arm bundle is
A4-reference-equivalent only; it is not current SAQ\textquotesingle s
serialized representation or

query estimator.

The corrected protocol authority passed independent timing-closure,
schema, and Git-authority review. Initial implementation commit
\texttt{3a4f7c5} then failed static review and remains non-evidence.
Corrected exact source commit \texttt{482c401} passed three independent
static-only reviews. Its canonical manifest binds 35 exact

sources at source-tree SHA-256 \texttt{d5b8374f...}.

The user later authorized exactly one frozen PAR invocation at reviewed
and pushed base \texttt{2e983a0}. It returned exit \texttt{3} with
frozen status \texttt{ARTIFACT\_INVALID}: \texttt{sys.executable} was
\texttt{/bin/python}, whose terminal path object was a symlink rejected
by

the conductor\textquotesingle s no-follow leader-identity read. The
failure occurred before host, NumPy, or process-inventory observation
and before any B/P worker. No CMake/build child, PCG64 fixture, parity
case, receipt, ledger, build manifest, parity summary, artifact index,
seal, or final PAR

directory was produced. The empty staging directory and ignored bytecode
caches are non-evidence.

The reviewed terminal claim is exactly:

\begin{Shaded}
\begin{Highlighting}[]
\NormalTok{ARTIFACT\_INVALID}
\NormalTok{REVIEWED\_TERMINAL\_NO\_VALID\_PAR\_AUTHORITY}
\NormalTok{NO\_SCIENTIFIC\_DECISION}
\end{Highlighting}
\end{Shaded}

This is artifact-validation failure evidence, not \texttt{PASS\_PARITY},
feasibility, quantization quality, natural-data, ANN, or systems
evidence. Static source review did not establish executable-identity
compatibility on the pinned host. Tooling correctness remains separate
from the database-systems contribution.

The later I-R1 correction changes only CPython leader identity. It opens
a stable descriptor for literal \texttt{/proc/self/exe}, requires
intentionally followed normalized \texttt{sys.executable} to name the
same nonempty regular inode, performs a bounded full EOF-checked read
with before/after

stability, and uses that identity consistently in conductor, runner
admission/receipts, and a physically independent positional-read
verifier. Generic artifact readers retain \texttt{O\_NOFOLLOW}.

Exact source \texttt{e48df452} and canonical 35-source tree
\texttt{8d8b5d3f8990e5d360e0a617d157a8b934c14d0f37b69f399bcc62c71f98360a}
passed three independent static reviews with no finding at LOW or above.
No Python module, compiler, build, test, fixture, RNG, native command,
or data path was executed. The maximum result

is therefore only:

\begin{Shaded}
\begin{Highlighting}[]
\NormalTok{SOURCE\_REPAIR\_STATIC\_REVIEW\_PASS}
\end{Highlighting}
\end{Shaded}

CACHE-P resolves only the protocol question raised by the five
quarantined bytecode files. Its primary-source review records that
CPython \texttt{-B} suppresses cache writes but does not disable reads
of an already present valid cache. The exact target therefore

retains those residuals as unread non-evidence, requires non-destructive
sanitized remote isolation, and rejects a direct transition to PAR-R1.
It preserves exact source repair \texttt{e48df452}, manifest blob
\texttt{2aa64e50...}, and the 35-source tree
\texttt{8d8b5d3f8990e5d360e0a617d157a8b934c14d0f37b69f399bcc62c71f98360a}.

The committed CACHE-P claim is exactly:

\begin{Shaded}
\begin{Highlighting}[]
\NormalTok{EXACT\_TARGET\_INDEPENDENT\_REVIEW\_PASS}
\NormalTok{GO\_SANITIZED\_REMOTE\_ISOLATION\_PROTOCOL\_REQUIRED}
\NormalTok{NO\_GO\_DIRECT\_PAR\_R1}
\end{Highlighting}
\end{Shaded}

This is artifact-governance evidence only. CACHE-P imported no module,
changed no source or environment, created no clone, prepared no remote
tree, and ran no build, fixture, parity, synthetic, base, query, or SAQ
work.

CACHE-I implements only the generic cache-governance layer frozen by
CACHE-P. Its exact target has 37 filesystem sources (ten Python and 27
native) plus one separately bound inline bootstrap, with source-tree
SHA-256
\texttt{9568007588c78ddda9fb4c4e20e8773ee2fa1da10a7656f181fef06883de38a2}.
The static closure binds four

schema/maximal-instance pairs, cache-aware PAR shapes, the independent
cache verifier, and a complete 38-unit import/process/mutation
crosswalk. Independent review found 0 issues at LOW or above.

After the source bytes were locked, a separately authorized Python 3.9
\texttt{compile()}-only check accepted the six changed/new outer sources
and inline bootstrap. No code object was executed and no module was
imported. The exact claim is only:

\begin{Shaded}
\begin{Highlighting}[]
\NormalTok{GENERIC\_CACHE\_POLICY\_SOURCE\_STATIC\_REVIEW\_PASS}
\end{Highlighting}
\end{Shaded}

This is artifact-governance source/static-review evidence, not PREP
readiness, runtime correctness, parity, feasibility, an SAQ limitation,
performance, novelty, or a method.

The later host-rebind erratum \texttt{5a47fed/@b89dabe} froze the
current \texttt{.el9\_8.2} leader and a bounded future rebind closure
but did not change the source. HOST-I-AUTH target/review
\texttt{212a67b/@71e6bec} then passed with zero LOW+ findings and only
\texttt{AUTHORIZATION\_EXACT\_TARGET\_REVIEW\_PASS}. It did not edit

the five sources or seven derived authority objects, establish
\texttt{HOST\_IDENTITY\_REBOUND}, or run Python, build, PREP, or data.
It authorizes no execution. When HOST-I was instructed, exact checking
found an unsatisfiable seven-versus-eight protocol-component closure
before any source target. Additive

source-authority erratum target/review \texttt{6fe8544/@9fa9528} passed
three static tracks and reached only
\texttt{HOST\_I\_SOURCE\_AUTHORITY\_ERRATUM\_REVIEW\_PASS}. It admits
the eighth component but changes no implementation source or derived
authority object. The later source-static target \texttt{4602585} was
independently reviewed at \texttt{e8e9c79} and failed

with one HIGH authority-identity mismatch: two derived documents cite
the wrong governing-review SHA-256. That target remains failed evidence.
The correction-only protocol target/review \texttt{ddfef99/@b1a7429}
passed with zero LOW-or-higher findings; exact R1 target/review
\texttt{e17f887/@e10bde7} then passed with zero LOW-or-higher findings

after reproducing the frozen four artifact substitutions, closure, and
13/13 host identities. \texttt{HOST\_IDENTITY\_REBOUND} is established
only for the registered 13-object bundle; whole-host identity is not
established. Source-history protocol target/review
\texttt{dcaed57/@e98a3e4} passed. The first authorized implementation
stopped before target

on its unsatisfiable ten-path closure, and path-closure erratum
target/review \texttt{150e5b3/@7749491} passed only for a future
eight-path/five-derived-object repair. Exact R1 target/review
\texttt{e9b5c83/@a1198c4} then reached
\texttt{SOURCE\_HISTORY\_EPOCH\_CORRECTION\_SOURCE\_STATIC\_REVIEW\_PASS},
correcting only the known static verifier rule with zero LOW-or-higher
findings. Because the verifier

was not run, cache-verifier/PAR readiness remains unestablished. Fresh
R1 PREP authorization target/review \texttt{cd1757e/@eb1b893}
subsequently passed exact documentation review. Its later unique
immediate probe passed, and the unique invocation was consumed but left
only durable START and empty captures

before the PREP body. Frozen status is
\texttt{CRASH\_OR\_UNKNOWN\ /\ NO\_SCIENTIFIC\_DECISION}. Crash-record
protocol target/review \texttt{1982865/@8bec546} then passed with zero
LOW-or-higher findings; it does not form the record or authorize
retry/repair. Audited closure \texttt{3577edd} instead stops further
artifact-recovery investment at portfolio

level. There is no next authorized research or execution step; reopening
requires a new explicit user decision, new question, and fresh
cost-justified clean protocol.

Closed downstream stages remain unauthorized and non-transitive:

\begin{itemize}
\tightlist
\item
  \textbf{Stage}: \texttt{A4-V2-I}

  \textbf{Scope}: source implementation and independent static review
  only
\item
  \textbf{Current authorization}: \texttt{COMPLETED\_REVIEW\_PASS}
\item
  \textbf{Stage}: \texttt{A4-V2-PAR}

  \textbf{Scope}: frozen build and parity only
\item
  \textbf{Current authorization}:
  \texttt{REVIEWED\_TERMINAL\_ARTIFACT\_INVALID}; no valid PAR authority
\item
  \textbf{Stage}: \texttt{A4-V2-I-R1}

  \textbf{Scope}: executable-identity source repair and static review
  only
\item
  \textbf{Current authorization}:
  \texttt{SOURCE\_REPAIR\_STATIC\_REVIEW\_PASS}
\item
  \textbf{Stage}: \texttt{A4-V2-CACHE-P}

  \textbf{Scope}: primary-source review and frozen staged-isolation
  protocol
\item
  \textbf{Current authorization}:
  \texttt{EXACT\_TARGET\_INDEPENDENT\_REVIEW\_PASS}; governance only
\item
  \textbf{Stage}: \texttt{A4-V2-CACHE-I}

  \textbf{Scope}: generic source/schema/static closure only
\item
  \textbf{Current authorization}:
  \texttt{GENERIC\_CACHE\_POLICY\_SOURCE\_STATIC\_REVIEW\_PASS};
  governance only
\item
  \textbf{Stage}: historical PREP authorization/review

  \textbf{Scope}: freeze one sanitized remote preparation event
\item
  \textbf{Current authorization}:
  \texttt{PREP\_AUTHORIZATION\_EXACT\_TARGET\_REVIEW\_PASS}; later
  nontransferable documentation governance only
\item
  \textbf{Stage}: \texttt{A4-V2-PREP-HOST-P}

  \textbf{Scope}: additive current-host rebind protocol and exact review
\item
  \textbf{Current authorization}:
  \texttt{HOST\_IDENTITY\_REBIND\_PROTOCOL\_REVIEW\_PASS}; host identity
  is not rebound
\item
  \textbf{Stage}: \texttt{A4-V2-PREP-HOST-I-AUTH}

  \textbf{Scope}: freeze only the contract/projection for a possible
  later source/review edge
\item
  \textbf{Current authorization}:
  \texttt{AUTHORIZATION\_EXACT\_TARGET\_REVIEW\_PASS}; documentation
  governance only; grants no HOST-I authority
\item
  \textbf{Stage}: attempted \texttt{A4-V2-PREP-HOST-I}

  \textbf{Scope}: validate the frozen projection before source editing
\item
  \textbf{Current authorization}: \texttt{STOPPED\_BEFORE\_TARGET};
  exact seven-versus-eight contradiction
\item
  \textbf{Stage}: \texttt{A4-V2-PREP-HOST-I-ERRATUM}

  \textbf{Scope}: admit the eighth protocol component and freeze exact
  future accounting
\item
  \textbf{Current authorization}:
  \texttt{HOST\_I\_SOURCE\_AUTHORITY\_ERRATUM\_REVIEW\_PASS};
  documentation governance only
\item
  \textbf{Stage}: \texttt{A4-V2-PREP-HOST-I}

  \textbf{Scope}: implement and independently review the coherent
  five-source/seven-derived-authority rebind
\item
  \textbf{Current authorization}:
  \texttt{SOURCE\_STATIC\_TARGET\_REVIEW\_FAIL\_AUTHORITY\_IDENTITY\_MISMATCH};
  target/review \texttt{4602585/@e8e9c79}, one HIGH, retained as
  immutable failed evidence
\item
  \textbf{Stage}: \texttt{A4-V2-PREP-HOST-I-R1-P}

  \textbf{Scope}: freeze a correction-only exact five-path target and
  its direct-child review projection
\item
  \textbf{Current authorization}:
  \texttt{HOST\_I\_R1\_CORRECTION\_ONLY\_REPAIR\_PROTOCOL\_REVIEW\_PASS};
  target/review \texttt{ddfef99/@b1a7429}, documentation governance only
\item
  \textbf{Stage}: \texttt{A4-V2-PREP-HOST-I-R1}

  \textbf{Scope}: apply only the frozen authority-identity correction
  and review it
\item
  \textbf{Current authorization}:
  \texttt{HOST\_I\_R1\_CORRECTION\_ONLY\_REPAIR\_REVIEW\_PASS};
  target/review \texttt{e17f887/@e10bde7}; rebound established for
  registered 13-object bundle only, with no execution
\item
  \textbf{Stage}: \texttt{A4-V2-SOURCE-HISTORY-P}

  \textbf{Scope}: freeze an exact fail-closed two-epoch correction
  protocol
\item
  \textbf{Current authorization}:
  \texttt{SOURCE\_HISTORY\_EPOCH\_CORRECTION\_PROTOCOL\_REVIEW\_PASS};
  target/review \texttt{dcaed57/@e98a3e4}, zero LOW+, protocol
  governance only
\item
  \textbf{Stage}: first \texttt{A4-V2-SOURCE-HISTORY-I}

  \textbf{Scope}: validate and implement the frozen ten-path projection
\item
  \textbf{Current authorization}:
  \texttt{STOPPED\_PRE\_TARGET\_PATH\_CLOSURE\_UNSATISFIABLE}; runtime
  schema and maximal witness had no legitimate delta
\item
  \textbf{Stage}: \texttt{A4-V2-SOURCE-HISTORY-I-R1-P}

  \textbf{Scope}: correct only the future path-closure projection
\item
  \textbf{Current authorization}:
  \texttt{SOURCE\_HISTORY\_I\_R1\_PATH\_CLOSURE\_ERRATUM\_REVIEW\_PASS};
  target/review \texttt{150e5b3/@7749491}, zero LOW+, protocol
  governance only
\item
  \textbf{Stage}: \texttt{A4-V2-SOURCE-HISTORY-I-R1}

  \textbf{Scope}: implement and statically review the corrected exact
  eight-path/five-derived-object repair
\item
  \textbf{Current authorization}:
  \texttt{SOURCE\_HISTORY\_EPOCH\_CORRECTION\_SOURCE\_STATIC\_REVIEW\_PASS};
  target/review \texttt{e9b5c83/@a1198c4}, zero LOW+, no execution or
  readiness claim
\item
  \textbf{Stage}: \texttt{A4-V2-CACHE-PREP-R1-AUTH}

  \textbf{Scope}: freeze the current dormant one-invocation contract and
  exact review projection
\item
  \textbf{Current authorization}:
  \texttt{PREP\_AUTHORIZATION\_EXACT\_TARGET\_REVIEW\_PASS};
  target/review \texttt{cd1757e/@eb1b893}, zero LOW+, documentation
  governance only
\item
  \textbf{Stage}: \texttt{A4-V2-CACHE-PREP-ISO-R1}

  \textbf{Scope}: run one sanitized preparation after the reserved
  immediate review-head probe
\item
  \textbf{Current authorization}:
  \texttt{CRASH\_OR\_UNKNOWN\ /\ START\_ONLY\ /\ NO\_SCIENTIFIC\_DECISION};
  unique probe/invocation consumed before PREP body
\item
  \textbf{Stage}: \texttt{A4-V2-CACHE-PREP-ISO-R1-CRASH-P}

  \textbf{Scope}: freeze a truthful future terminal-record topology
\item
  \textbf{Current authorization}:
  \texttt{START\_ONLY\_CRASH\_RECORD\_PROTOCOL\_REVIEW\_PASS};
  target/review \texttt{1982865/@8bec546}, governance only
\item
  \textbf{Stage}: \texttt{A4-V2-CACHE-PREP-ISO-R1-CRASH-I}

  \textbf{Scope}: form and independently review the START-only terminal
  record
\item
  \textbf{Current authorization}: \texttt{NOT\_AUTHORIZED}
\item
  \textbf{Stage}: \texttt{A4-V2-CACHE-BIND}

  \textbf{Scope}: bind only the independently reviewed sanitized
  preparation
\item
  \textbf{Current authorization}: \texttt{NOT\_AUTHORIZED}
\item
  \textbf{Stage}: \texttt{A4-V2-PAR-R1}

  \textbf{Scope}: at most one corrected frozen build/parity event after
  the full isolation chain
\item
  \textbf{Current authorization}: \texttt{NOT\_AUTHORIZED}; direct
  transition rejected
\item
  \textbf{Stage}: \texttt{A4-V2-SRUN}

  \textbf{Scope}: one logical synthetic admission event
\item
  \textbf{Current authorization}: \texttt{NOT\_AUTHORIZED}
\item
  \textbf{Stage}: data / SAQ

  \textbf{Scope}: base, query, index, estimator, integration
\item
  \textbf{Current authorization}: \texttt{NOT\_AUTHORIZED}
\end{itemize}
\end{frame}

\begin{frame}[fragile,allowframebreaks]{66. Current Synthesis}
\protect\phantomsection\label{66-current-synthesis}
\begin{itemize}
\tightlist
\item
  \textbf{Attempt}: 1A full-D PCA replacement

  \textbf{Theoretical guarantee}: L2 isometry for any orthogonal
  transform
\item
  \textbf{Strongest positive evidence}: GIST residual-PCA accurate/full
  RMSE \texttt{-0.60\%/-0.33\%}

  \textbf{Decision / current boundary}: no stable ranking gain; fast
  harm; CIFAR replication failure
\item
  \textbf{Attempt}: 1B lossy \texttt{D\ -\textgreater{}\ d}

  \textbf{Theoretical guarantee}: exact head and norm terms in favorable
  oracle
\item
  \textbf{Strongest positive evidence}: tail norms reduce RMSE from
  0.0445 to 0.00248

  \textbf{Decision / current boundary}: omitted tail IP still worsens
  ranking versus native SAQ
\item
  \textbf{Attempt}: 2 exact scalar DP

  \textbf{Theoretical guarantee}: exact bin-boundary partition SSE using
  raw bin moments; final midpoint-nearest-centroid raw SSE is evaluated,
  not reoptimized; outer optimum is conditional on \texttt{E{[}j,b{]}}
\item
  \textbf{Strongest positive evidence}: audio B=4 raw MSE
  \texttt{-15.5\%}, R@100 \texttt{+0.004}

  \textbf{Decision / current boundary}: inner-DP novelty is foreclosed
  by prior art; no full-scale raw optimum or recall guarantee;
  cross-regime reversals
\item
  \textbf{Attempt}: 3 distance-quality re-evaluation

  \textbf{Theoretical guarantee}: paper-exact metric semantics; no
  method guarantee
\item
  \textbf{Strongest positive evidence}: GIST measured point: higher
  \texttt{1/Ratio}, \texttt{1.078x} QPS at the frozen target

  \textbf{Decision / current boundary}: one positive setting; DEEP
  controls remain negative; metric is prior work
\item
  \textbf{Attempt}: 4 arbitrary-cardinality fixed-rate words

  \textbf{Theoretical guarantee}: the arbitrary positive-integer
  feasible set contains the dyadic set for exact factorized SSE
\item
  \textbf{Strongest positive evidence}: frozen \texttt{(3,5)} witness
  and A4-1S parity; later positive records are
  protocol/source/artifact-governance evidence only

  \textbf{Decision / current boundary}: preserve A4-1S
  \texttt{NO\_GO\_EXACT\_SOLVER\_COST} and PAR
  \texttt{ARTIFACT\_INVALID\ /\ NO\_SCIENTIFIC\_DECISION};
  source-history R1 corrected only the static rule; the unique PREP
  invocation later stopped at START-only \texttt{CRASH\_OR\_UNKNOWN},
  and reviewed protocol \texttt{1982865/@8bec546} closes only the future
  record topology; CRASH-I, repair/retry, binding,
\item
  and PAR-R1 remain unauthorized
\end{itemize}

Cross-attempt lesson:

\begin{Shaded}
\begin{Highlighting}[]
\NormalTok{Optimizing a mathematically valid surrogate is not enough, and changing the}
\NormalTok{evaluation metric is not itself a mechanism.}

\NormalTok{The missing contribution must connect its optimized quantity to top{-}k or}
\NormalTok{downstream quality and system work, survive a second baseline or regime, and}
\NormalTok{include all construction, storage, and query overhead.}
\end{Highlighting}
\end{Shaded}

\begin{Shaded}
\begin{Highlighting}[]
\NormalTok{Attempt 4 was deliberately staged so that a representation theorem or a}
\NormalTok{synthetic witness could not silently become a method claim. The committed and}
\NormalTok{reviewed A4{-}1S cost stop prevented the next expansion of scope. V2 changes the}
\NormalTok{future measurement contract, its additive erratum closes one reporting}
\NormalTok{self{-}reference, and its corrected source passed static review. The later PAR}
\NormalTok{failure shows that static review did not guarantee executable{-}identity}
\NormalTok{compatibility; it is tooling/artifact evidence only and says nothing about}
\end{Highlighting}
\end{Shaded}

\begin{Shaded}
\begin{Highlighting}[]
\NormalTok{instrument feasibility. I{-}R1 closes the identified source defect under static}
\NormalTok{review, but without corrected execution it still says nothing about parity or}
\NormalTok{instrument feasibility. CACHE{-}P freezes how a future clean artifact could be}
\NormalTok{isolated, and CACHE{-}I statically closes only the corresponding generic source}
\NormalTok{and schema layer. Their reviews and the narrow syntax check are governance}
\NormalTok{evidence, not feasibility evidence. The later PREP authorization record also}
\NormalTok{passed exact{-}target review, but a later static host mismatch retired its}
\end{Highlighting}
\end{Shaded}

\begin{Shaded}
\begin{Highlighting}[]
\NormalTok{activation path before PREP. The reviewed additive host erratum freezes only}
\NormalTok{a future rebind contract; it does not rebind the host or grant PREP or}
\NormalTok{execution authority. The later HOST{-}I{-}AUTH review freezes only the contract}
\NormalTok{and projection for a possible source edge and grants no HOST{-}I authority; no}
\NormalTok{source/authority rebind, Python, or PREP ran. Its attempted HOST{-}I edge stopped}
\NormalTok{before edits on a seven{-}versus{-}eight component contradiction. The reviewed}
\NormalTok{source{-}authority erratum fixed that contract admission. The later source{-}}
\end{Highlighting}
\end{Shaded}

\begin{Shaded}
\begin{Highlighting}[]
\NormalTok{static target was formed, but its direct{-}child review failed on one HIGH}
\NormalTok{authority{-}identity mismatch. Its later correction{-}only target passed exact}
\NormalTok{review and established rebound only for the registered 13{-}object bundle.}
\NormalTok{Whole{-}host and execution readiness remain unestablished. A later reviewed}
\NormalTok{source{-}history protocol froze a bounded two{-}epoch repair; after its first}
\NormalTok{ten{-}path implementation projection failed, an erratum and exact R1 target}
\NormalTok{closed the static eight{-}path repair under direct{-}child review. The corrected}
\end{Highlighting}
\end{Shaded}

\begin{Shaded}
\begin{Highlighting}[]
\NormalTok{verifier was not run, so cache/PAR readiness still cannot be claimed. Fresh}
\NormalTok{R1 PREP authorization target/review \textasciigrave{}cd1757e/@eb1b893\textasciigrave{} then passed exact}
\NormalTok{documentation review. The subsequently authorized unique probe/invocation was}
\NormalTok{consumed and stopped at START{-}only \textasciigrave{}CRASH\_OR\_UNKNOWN\textasciigrave{}; crash{-}record protocol}
\NormalTok{\textasciigrave{}1982865/@8bec546\textasciigrave{} passed, while CRASH{-}I, repair, retry, and CACHE{-}BIND remain}
\NormalTok{unauthorized.}
\end{Highlighting}
\end{Shaded}
\end{frame}

\begin{frame}[fragile,allowframebreaks]{67. Proposed Meeting Discussion}
\protect\phantomsection\label{67-proposed-meeting-discussion}
Decision 1:

\begin{Shaded}
\begin{Highlighting}[]
\NormalTok{Do the Attempt 1 results justify treating PCA replacement as a closed practical}
\NormalTok{direction, while avoiding a universal PCA{-}optimality claim?}
\end{Highlighting}
\end{Shaded}

Decision 2:

\begin{Shaded}
\begin{Highlighting}[]
\NormalTok{Should exact{-}hist DP remain only a stronger offline baseline, given that}
\NormalTok{Wu 1991 / GLM 2017 foreclose the inner algorithmic novelty and White{-}Singal}
\NormalTok{2026 explicitly recognizes that baseline, while the local evidence does not}
\NormalTok{establish retrieval or system dominance?}
\end{Highlighting}
\end{Shaded}

Decision 3:

\begin{Shaded}
\begin{Highlighting}[]
\NormalTok{Should future ANN evidence always report both identifier Recall and geometric}
\NormalTok{distance quality, while treating the GIST flip only as measurement evidence?}
\end{Highlighting}
\end{Shaded}

Decision 4:

\begin{Shaded}
\begin{Highlighting}[]
\NormalTok{Do we agree to preserve the terminal A4{-}1S pipeline{-}cost failure, while}
\NormalTok{treating the reviewed A4 V2 PAR \textasciigrave{}ARTIFACT\_INVALID\textasciigrave{} result only as an artifact{-}}
\NormalTok{validation failure, with no inference about arbitrary{-}cardinality}
\NormalTok{quantization or instrument feasibility; treating I{-}R1 only as a reviewed}
\NormalTok{source repair rather than artifact readiness; and accepting CACHE{-}P only at}
\NormalTok{its artifact{-}governance ceiling, including sanitized remote isolation and the}
\NormalTok{no{-}go on direct PAR{-}R1; and treating CACHE{-}I\textquotesingle{}s static pass as governance}
\end{Highlighting}
\end{Shaded}

\begin{Shaded}
\begin{Highlighting}[]
\NormalTok{source evidence rather than PREP readiness; and treating the later PREP}
\NormalTok{authorization review as a dormant contract rather than a completed clone or}
\NormalTok{feasibility result; treating HOST{-}I{-}AUTH and the source{-}authority erratum as}
\NormalTok{governance only; and treating the later HOST{-}I source{-}static target as a}
\NormalTok{failed authority{-}DAG artifact whose otherwise passing static checks do not}
\NormalTok{establish host rebound by themselves; while treating the reviewed R1 repair}
\NormalTok{PASS as bounded identity governance for the registered 13{-}object bundle—not}
\end{Highlighting}
\end{Shaded}

\begin{Shaded}
\begin{Highlighting}[]
\NormalTok{whole{-}host identity, execution, PREP readiness, or cache{-}verifier/PAR}
\NormalTok{readiness?}
\end{Highlighting}
\end{Shaded}

The current authorization boundary is narrow:

\begin{Shaded}
\begin{Highlighting}[]
\NormalTok{A4{-}V2{-}PAR was invoked once and terminated before build/parity.}
\NormalTok{There is no active execution authority and no valid PAR authority.}

\NormalTok{The clean I{-}R1 source commit and independent static review are complete.}
\NormalTok{CACHE{-}P target/review \textasciigrave{}56210f8/@249d5b8\textasciigrave{} froze}
\NormalTok{\textasciigrave{}GO\_SANITIZED\_REMOTE\_ISOLATION\_PROTOCOL\_REQUIRED / NO\_GO\_DIRECT\_PAR\_R1\textasciigrave{}.}
\NormalTok{CACHE{-}I target/review \textasciigrave{}c33a2bf/@5db3025\textasciigrave{} subsequently reached only}
\end{Highlighting}
\end{Shaded}

\begin{Shaded}
\begin{Highlighting}[]
\NormalTok{\textasciigrave{}GENERIC\_CACHE\_POLICY\_SOURCE\_STATIC\_REVIEW\_PASS\textasciigrave{}; the quarantine remains}
\NormalTok{unread and uncleaned. PREP authorization target/review \textasciigrave{}e7f940e/@16a8201\textasciigrave{}}
\NormalTok{subsequently reached only \textasciigrave{}PREP\_AUTHORIZATION\_EXACT\_TARGET\_REVIEW\_PASS\textasciigrave{}.}
\NormalTok{Static checking then found the \textasciigrave{}.el9\_8 {-}\textgreater{} .el9\_8.2\textasciigrave{} leader replacement before}
\NormalTok{PREP or START. Host{-}rebind erratum target/review \textasciigrave{}5a47fed/@b89dabe\textasciigrave{} reached}
\NormalTok{only \textasciigrave{}HOST\_IDENTITY\_REBIND\_PROTOCOL\_REVIEW\_PASS\textasciigrave{}; that protocol stage did not}
\NormalTok{rebind host identity.}
\end{Highlighting}
\end{Shaded}

\begin{Shaded}
\begin{Highlighting}[]
\NormalTok{The old authority is nontransferable and the old probe is superseded and must}
\NormalTok{never run or be reused. HOST{-}I{-}AUTH target/review \textasciigrave{}212a67b/@71e6bec\textasciigrave{} reached}
\NormalTok{only \textasciigrave{}AUTHORIZATION\_EXACT\_TARGET\_REVIEW\_PASS\textasciigrave{}; it performed no five{-}source or}
\NormalTok{seven{-}derived{-}object rebind and ran no Python or PREP. The attempted HOST{-}I}
\NormalTok{edge then stopped before any source target on an exact seven{-}versus{-}eight}
\NormalTok{component contradiction. Source{-}authority erratum target/review}
\NormalTok{\textasciigrave{}6fe8544/@9fa9528\textasciigrave{} reached only}
\end{Highlighting}
\end{Shaded}

\begin{Shaded}
\begin{Highlighting}[]
\NormalTok{\textasciigrave{}HOST\_I\_SOURCE\_AUTHORITY\_ERRATUM\_REVIEW\_PASS\textasciigrave{}. The subsequently authorized}
\NormalTok{HOST{-}I target/review \textasciigrave{}4602585/@e8e9c79\textasciigrave{} ended at}
\NormalTok{\textasciigrave{}SOURCE\_STATIC\_TARGET\_REVIEW\_FAIL\_AUTHORITY\_IDENTITY\_MISMATCH\textasciigrave{}: one HIGH,}
\NormalTok{host rebound not established. Correction{-}only protocol target/review}
\NormalTok{\textasciigrave{}ddfef99/@b1a7429\textasciigrave{} then reached}
\NormalTok{\textasciigrave{}HOST\_I\_R1\_CORRECTION\_ONLY\_REPAIR\_PROTOCOL\_REVIEW\_PASS\textasciigrave{}. Exact R1 target/review}
\NormalTok{\textasciigrave{}e17f887/@e10bde7\textasciigrave{} subsequently reached}
\end{Highlighting}
\end{Shaded}

\begin{Shaded}
\begin{Highlighting}[]
\NormalTok{\textasciigrave{}HOST\_I\_R1\_CORRECTION\_ONLY\_REPAIR\_REVIEW\_PASS\textasciigrave{}, establishing rebound only for}
\NormalTok{the registered 13{-}object bundle. Whole{-}host identity remains unestablished and}
\NormalTok{no Python, build, PREP, or scientific execution ran. A future}
\NormalTok{source{-}history mismatch then received reviewed protocol target/review}
\NormalTok{\textasciigrave{}dcaed57/@e98a3e4\textasciigrave{}. Its first implementation authorization stopped before}
\NormalTok{target on the unsatisfiable ten{-}path closure; reviewed path{-}closure erratum}
\NormalTok{\textasciigrave{}150e5b3/@7749491\textasciigrave{} froze the exact eight{-}path/five{-}derived{-}object R1 repair.}
\end{Highlighting}
\end{Shaded}

\begin{Shaded}
\begin{Highlighting}[]
\NormalTok{Exact target/review \textasciigrave{}e9b5c83/@a1198c4\textasciigrave{} then reached}
\NormalTok{\textasciigrave{}SOURCE\_HISTORY\_EPOCH\_CORRECTION\_SOURCE\_STATIC\_REVIEW\_PASS\textasciigrave{} with zero LOW+.}
\NormalTok{The verifier was not run, so cache{-}verifier/PAR readiness is not established.}
\NormalTok{Fresh dormant PREP authorization target/review \textasciigrave{}cd1757e/@eb1b893\textasciigrave{} then}
\NormalTok{reached \textasciigrave{}PREP\_AUTHORIZATION\_EXACT\_TARGET\_REVIEW\_PASS\textasciigrave{} with zero LOW+. Its}
\NormalTok{later unique immediate probe and invocation were consumed and left only START}
\NormalTok{plus empty captures before the PREP body, so the result is}
\end{Highlighting}
\end{Shaded}

\begin{Shaded}
\begin{Highlighting}[]
\NormalTok{\textasciigrave{}CRASH\_OR\_UNKNOWN / NO\_SCIENTIFIC\_DECISION\textasciigrave{}. Crash{-}record protocol}
\NormalTok{\textasciigrave{}1982865/@8bec546\textasciigrave{} passed independent review. No execution node is currently}
\NormalTok{authorized; stop for a user checkpoint. CRASH{-}I, source repair/retry,}
\NormalTok{CACHE{-}BIND, explicit PAR{-}R1, and SRUN remain separate and unauthorized;}
\NormalTok{base/query inputs, SAQ changes, and the old A4{-}1 gate remain unauthorized.}
\end{Highlighting}
\end{Shaded}

Speaker notes:

\begin{itemize}
\tightlist
\item
  Attempt 4 completed synthetic correctness validation and a
  falsification of affordability under the frozen pipeline-cost gate;
  allocation, block-VQ, and encoding cost phases were not reached after
  the scalar-prefix early stop.
\item
  A4 V2\textquotesingle s scientific PAR work remains unexecuted: the
  conductor was invoked but stopped at executable-identity admission
  before build, RNG, or parity. Its erratum repairs reporting closure
  and its source pass checks artifact machinery only; I-R1 statically
  repairs the

  identified path but has not been executed. The old \texttt{5/2}
  projection is not a V2 cost model.
\item
  Novelty, mechanism, replication, and total overhead should be
  challenged before experimental scope is expanded.
\end{itemize}
\end{frame}

\begin{frame}[fragile,allowframebreaks]{68. Evidence And Code Map}
\protect\phantomsection\label{68-evidence-and-code-map}
Attempt 1, SAQ branch \texttt{saq-transform-analysis@3d94840}:

\begin{Shaded}
\begin{Highlighting}[]
\NormalTok{docs/saq\_transform\_phase1\_limitation\_evidence\_2026\_07\_10.md}
\NormalTok{docs/saq\_transform\_phase1b\_external\_replication\_evidence\_2026\_07\_10.md}
\NormalTok{docs/saq\_lossy\_projection\_lp0\_gate\_a\_evidence\_2026\_07\_11.md}
\NormalTok{script/prepare\_phase1\_transform\_views.py}
\NormalTok{src/phase1\_transform\_diagnostic.cpp}
\NormalTok{src/lp0\_projection\_gate\_a.cpp}
\NormalTok{script/evaluate\_phase1b\_gate.py}
\end{Highlighting}
\end{Shaded}

\begin{Shaded}
\begin{Highlighting}[]
\NormalTok{script/evaluate\_lp0\_gate\_a.py}
\end{Highlighting}
\end{Shaded}

Lossy projection evidence is preserved on
\texttt{saq-lossy-projection-analysis@051ec6a}.

Attempt 2, sibling \texttt{vectordb@f51b487}:

\begin{Shaded}
\begin{Highlighting}[]
\NormalTok{src/scalar\_quantizer.cpp}
\NormalTok{src/dimensionwise\_quantizer.cpp}
\NormalTok{src/bit\_allocator.cpp}
\NormalTok{src/eval\_dimensionwise\_quantizer.cpp}
\NormalTok{docs/dimensionwise\_quantizer\_smoke\_run\_2026\_06\_26.md}
\NormalTok{reports/scalar\_training\_exact\_hist\_audit\_2026\_06\_30/README.md}
\NormalTok{reports/scalar\_training\_exact\_hist\_audit\_2026\_06\_30/comparison.csv}
\end{Highlighting}
\end{Shaded}

\begin{Shaded}
\begin{Highlighting}[]
\NormalTok{reports/scalar\_training\_exact\_hist\_audit\_2026\_06\_30/logs/}
\NormalTok{  deep1M\_pca\_B4\_variable\_exact{-}hist\_rank{-}boundary.stdout}
\NormalTok{reports/pca\_rank\_boundary\_deep1M\_2026\_06\_27/logs/}
\NormalTok{  deep1M\_pca\_B4\_variable\_lloyd\_rank{-}boundary.stdout}
\end{Highlighting}
\end{Shaded}

Attempt 2 external primary sources:

\begin{itemize}
\tightlist
\item
  \href{https://arxiv.org/html/2606.00289v1}{White and Singal, Inner
  Product Aware Quantization, arXiv:2606.00289v1}
\item
  \href{https://github.com/nathanllww/Inner-Product-Aware-Quantization/tree/e92ed904b85ad12469a0a350c81071ff9ef6aa37}{Pinned
  official code at \texttt{e92ed90}}
\item
  \href{https://arxiv.org/abs/1701.07204}{Grønlund et al., Fast Exact
  k-Means, k-Medians and Bregman Divergence Clustering in 1D}
\item
  \href{https://www.sciencedirect.com/science/article/pii/0196677491900392}{Wu,
  Optimal Quantization by Matrix Searching, Journal of Algorithms, 1991}
\item
  \href{https://arxiv.org/abs/2509.12086}{Li et al., SAQ: Pushing the
  Limits of Vector Quantization through Code Adjustment and Dimension
  Segmentation, arXiv:2509.12086}, direct ANN prior for DP segmentation
  and bit allocation under a space quota
\end{itemize}

Attempt 3, branch \texttt{saq-ratio-metric-analysis@146dc16}:

\begin{Shaded}
\begin{Highlighting}[]
\NormalTok{docs/saq\_attempt3\_ratio\_metric\_related\_work\_and\_gate\_2026\_07\_13.md}
\NormalTok{docs/saq\_attempt3\_ratio\_metric\_protocol\_2026\_07\_13.md}
\NormalTok{docs/saq\_attempt3\_a3\_0\_a3\_1\_gist\_evidence\_2026\_07\_13.md}
\NormalTok{docs/saq\_attempt3\_a3\_2\_a3\_3\_decision\_2026\_07\_13.md}
\NormalTok{docs/saq\_attempt3\_a3\_1b\_artifacts\_2026\_07\_13/}
\NormalTok{docs/saq\_attempt3\_a3\_2\_artifacts\_2026\_07\_13/}
\NormalTok{script/evaluate\_inverse\_ratio.py}
\end{Highlighting}
\end{Shaded}

\begin{Shaded}
\begin{Highlighting}[]
\NormalTok{script/summarize\_ratio\_frontier.py}
\NormalTok{src/test\_qps.cpp}
\end{Highlighting}
\end{Shaded}

Attempt 4, branch \texttt{saq-arbitrary-cardinality-analysis@f1b464b}:

\begin{Shaded}
\begin{Highlighting}[]
\NormalTok{docs/saq\_attempt4\_arbitrary\_cardinality\_related\_work\_and\_gate\_2026\_07\_13.md}
\NormalTok{docs/saq\_attempt4\_arbitrary\_cardinality\_sources\_2026\_07\_13.json}
\NormalTok{docs/saq\_attempt4\_a4\_0\_offline\_protocol\_2026\_07\_13.md}
\NormalTok{docs/saq\_attempt4\_a4\_0\_synthetic\_evidence\_2026\_07\_13.md}
\NormalTok{docs/saq\_attempt4\_a4\_0\_artifacts\_2026\_07\_13/synthetic\_witness.json}
\NormalTok{docs/saq\_attempt4\_a4\_1\_base\_only\_feasibility\_preregistration\_2026\_07\_13.md}
\NormalTok{docs/saq\_attempt4\_a4\_1\_base\_only\_input\_spec\_2026\_07\_13.json}
\end{Highlighting}
\end{Shaded}

\begin{Shaded}
\begin{Highlighting}[]
\NormalTok{docs/saq\_attempt4\_a4\_1\_base\_only\_hypotheses\_2026\_07\_13.json}
\NormalTok{docs/saq\_attempt4\_a4\_1s\_synthetic\_implementation\_protocol\_2026\_07\_13.md}
\NormalTok{docs/saq\_attempt4\_a4\_1s\_artifacts\_2026\_07\_13/parity\_artifact\_index.json}
\NormalTok{docs/saq\_attempt4\_a4\_1s\_implementation\_parity\_review\_2026\_07\_13.md}
\NormalTok{docs/saq\_attempt4\_a4\_1s\_artifacts\_2026\_07\_13/synthetic\_cost\_projection\_manifest.json}
\NormalTok{docs/saq\_attempt4\_a4\_1s\_artifacts\_2026\_07\_13/synthetic\_cost\_projection\_summary.json}
\NormalTok{docs/saq\_attempt4\_a4\_1s\_artifacts\_2026\_07\_13/synthetic\_cost\_projection\_detail\_ledger.json}
\end{Highlighting}
\end{Shaded}

\begin{Shaded}
\begin{Highlighting}[]
\NormalTok{docs/saq\_attempt4\_a4\_1s\_artifacts\_2026\_07\_13/cost\_projection\_artifact\_index.json}
\NormalTok{docs/saq\_attempt4\_a4\_1s\_cost\_projection\_review\_2026\_07\_13.md}
\end{Highlighting}
\end{Shaded}

Attempt 4 V2 corrected protocol-authority snapshot
\texttt{saq-arbitrary-cardinality-feasibility-v2@f13a383}, corrected
source snapshot \texttt{@482c401}, terminal PAR result/review
\texttt{@30dfada/@fd5367e}, and exact I-R1 source-repair/review
\texttt{@e48df452/@c401dae}, followed by CACHE-P exact content
target/review record \texttt{@56210f8/@249d5b8} and CACHE-I exact
source/static target/review \texttt{@c33a2bf/@5db3025},

then the exact PREP authorization target/review
\texttt{@e7f940e/@16a8201}, and host-rebind erratum target/review
\texttt{@5a47fed/@b89dabe}, followed by HOST-I-AUTH target/review
\texttt{@212a67b/@71e6bec}, then source-authority erratum target/review
\texttt{@6fe8544/@9fa9528}, and finally HOST-I source-static
target/review \texttt{@4602585/@e8e9c79}, followed by HOST-I-R1
correction-only repair

protocol target/review \texttt{@ddfef99/@b1a7429} and correction
target/review \texttt{@e17f887/@e10bde7}, then source-history
epoch-correction protocol target/review \texttt{@dcaed57/@e98a3e4},
followed by source-history path-closure erratum target/review
\texttt{@150e5b3/@7749491}, and exact source-history R1 target/review
\texttt{@e9b5c83/@a1198c4}, followed by fresh R1 PREP authorization
target/review \texttt{@cd1757e/@eb1b893},

the consumed START-only invocation, and crash-record protocol
target/review \texttt{@1982865/@8bec546}:

\begin{Shaded}
\begin{Highlighting}[]
\NormalTok{docs/saq\_a4\_v2\_primary\_source\_metadata\_2026\_07\_14.json}
\NormalTok{docs/saq\_a4\_v2\_cost\_evidence\_primary\_source\_review\_2026\_07\_14.md}
\NormalTok{docs/saq\_a4\_v2\_cost\_evidence\_go\_no\_go\_memo\_2026\_07\_14.md}
\NormalTok{docs/saq\_a4\_v2\_synthetic\_construction\_preregistration\_2026\_07\_14.md}
\NormalTok{docs/saq\_a4\_v2\_synthetic\_construction\_contract\_2026\_07\_14.json}
\NormalTok{docs/saq\_a4\_v2\_artifact\_schema\_2026\_07\_14.json}
\NormalTok{docs/saq\_a4\_v2\_protocol\_independent\_review\_2026\_07\_14.md}
\end{Highlighting}
\end{Shaded}

\begin{Shaded}
\begin{Highlighting}[]
\NormalTok{docs/saq\_a4\_v2\_par\_report\_erratum\_authorization\_2026\_07\_14.md}
\NormalTok{docs/saq\_a4\_v2\_par\_report\_timing\_closure\_erratum\_2026\_07\_14.md}
\NormalTok{docs/saq\_a4\_v2\_par\_report\_timing\_closure\_erratum\_2026\_07\_14.json}
\NormalTok{docs/saq\_a4\_v2\_par\_report\_seal\_schema\_2026\_07\_14.json}
\NormalTok{docs/saq\_a4\_v2\_par\_report\_seal\_maximal\_instance\_2026\_07\_14.json}
\NormalTok{docs/saq\_a4\_v2\_protocol\_authority\_manifest\_2026\_07\_14.json}
\NormalTok{docs/saq\_a4\_v2\_par\_report\_timing\_closure\_erratum\_independent\_review\_2026\_07\_14.md}
\end{Highlighting}
\end{Shaded}

\begin{Shaded}
\begin{Highlighting}[]
\NormalTok{docs/saq\_a4\_v2\_implementation\_binding\_2026\_07\_14.md}
\NormalTok{docs/saq\_a4\_v2\_implementation\_manifest\_2026\_07\_14.json}
\NormalTok{docs/saq\_a4\_v2\_source\_provenance\_crosswalk\_2026\_07\_14.md}
\NormalTok{docs/saq\_a4\_v2\_implementation\_independent\_review\_2026\_07\_15.md}
\NormalTok{docs/saq\_a4\_v2\_par\_authorization\_2026\_07\_15.md}
\NormalTok{docs/saq\_a4\_v2\_par\_prebuild\_artifact\_identity\_failure\_2026\_07\_15.md}
\NormalTok{docs/saq\_a4\_v2\_par\_prebuild\_artifact\_identity\_failure\_independent\_review\_2026\_07\_15.md}
\end{Highlighting}
\end{Shaded}

\begin{Shaded}
\begin{Highlighting}[]
\NormalTok{docs/saq\_a4\_v2\_executable\_identity\_source\_repair\_authorization\_2026\_07\_15.md}
\NormalTok{docs/saq\_a4\_v2\_executable\_identity\_source\_repair\_independent\_review\_2026\_07\_15.md}
\NormalTok{docs/saq\_a4\_v2\_cache\_staging\_policy\_authorization\_2026\_07\_15.md}
\NormalTok{docs/saq\_a4\_v2\_cache\_staging\_primary\_sources\_2026\_07\_15.json}
\NormalTok{docs/saq\_a4\_v2\_cache\_staging\_primary\_source\_review\_2026\_07\_15.md}
\NormalTok{docs/saq\_a4\_v2\_cache\_staging\_disposition\_protocol\_2026\_07\_15.md}
\NormalTok{docs/saq\_a4\_v2\_cache\_staging\_disposition\_contract\_2026\_07\_15.json}
\end{Highlighting}
\end{Shaded}

\begin{Shaded}
\begin{Highlighting}[]
\NormalTok{docs/saq\_a4\_v2\_cache\_staging\_disposition\_protocol\_independent\_review\_2026\_07\_15.md}
\NormalTok{docs/saq\_a4\_v2\_cache\_implementation\_authorization\_2026\_07\_15.md}
\NormalTok{docs/saq\_a4\_v2\_cache\_implementation\_authorization\_independent\_review\_2026\_07\_15.md}
\NormalTok{docs/saq\_a4\_v2\_cache\_protocol\_authority\_manifest\_2026\_07\_15.json}
\NormalTok{docs/saq\_a4\_v2\_cache\_static\_closure\_2026\_07\_15.json}
\NormalTok{docs/saq\_a4\_v2\_cache\_runtime\_schema\_2026\_07\_15.json}
\NormalTok{docs/saq\_a4\_v2\_cache\_prep\_binding\_schema\_2026\_07\_15.json}
\end{Highlighting}
\end{Shaded}

\begin{Shaded}
\begin{Highlighting}[]
\NormalTok{docs/saq\_a4\_v2\_isolated\_clone\_prep\_receipt\_schema\_2026\_07\_15.json}
\NormalTok{docs/saq\_a4\_v2\_isolated\_clone\_prep\_token\_schema\_2026\_07\_15.json}
\NormalTok{docs/saq\_a4\_v2\_cache\_implementation\_independent\_review\_2026\_07\_15.md}
\NormalTok{docs/saq\_a4\_v2\_isolated\_clone\_prep\_authorization\_2026\_07\_15.md}
\NormalTok{docs/saq\_a4\_v2\_isolated\_clone\_prep\_authorization\_independent\_review\_2026\_07\_15.md}
\NormalTok{docs/saq\_a4\_v2\_prep\_host\_identity\_rebind\_erratum\_protocol\_2026\_07\_16.md}
\NormalTok{docs/saq\_a4\_v2\_prep\_host\_identity\_rebind\_erratum\_contract\_2026\_07\_16.json}
\end{Highlighting}
\end{Shaded}

\begin{Shaded}
\begin{Highlighting}[]
\NormalTok{docs/saq\_a4\_v2\_prep\_host\_identity\_rebind\_erratum\_independent\_review\_2026\_07\_16.md}
\NormalTok{docs/saq\_a4\_v2\_prep\_host\_identity\_rebind\_implementation\_authorization\_2026\_07\_16.md}
\NormalTok{docs/saq\_a4\_v2\_prep\_host\_identity\_rebind\_implementation\_authorization\_independent\_review\_2026\_07\_16.md}
\NormalTok{docs/saq\_a4\_v2\_prep\_host\_identity\_rebind\_implementation\_erratum\_protocol\_2026\_07\_16.md}
\NormalTok{docs/saq\_a4\_v2\_prep\_host\_identity\_rebind\_implementation\_erratum\_contract\_2026\_07\_16.json}
\NormalTok{docs/saq\_a4\_v2\_prep\_host\_identity\_rebind\_implementation\_erratum\_independent\_review\_2026\_07\_16.md}
\NormalTok{docs/saq\_a4\_v2\_prep\_host\_identity\_rebind\_implementation\_independent\_review\_2026\_07\_16.md}
\end{Highlighting}
\end{Shaded}

\begin{Shaded}
\begin{Highlighting}[]
\NormalTok{docs/saq\_a4\_v2\_prep\_host\_identity\_rebind\_r1\_correction\_only\_repair\_protocol\_2026\_07\_17.md}
\NormalTok{docs/saq\_a4\_v2\_prep\_host\_identity\_rebind\_r1\_correction\_only\_repair\_contract\_2026\_07\_17.json}
\NormalTok{docs/saq\_a4\_v2\_prep\_host\_identity\_rebind\_r1\_correction\_only\_repair\_protocol\_independent\_review\_2026\_07\_17.md}
\NormalTok{docs/saq\_a4\_v2\_prep\_host\_identity\_rebind\_r1\_correction\_only\_repair\_independent\_review\_2026\_07\_17.md}
\NormalTok{docs/saq\_a4\_v2\_source\_history\_epoch\_correction\_protocol\_2026\_07\_17.md}
\NormalTok{docs/saq\_a4\_v2\_source\_history\_epoch\_correction\_contract\_2026\_07\_17.json}
\NormalTok{docs/saq\_a4\_v2\_source\_history\_epoch\_correction\_protocol\_independent\_review\_2026\_07\_17.md}
\end{Highlighting}
\end{Shaded}

\begin{Shaded}
\begin{Highlighting}[]
\NormalTok{docs/saq\_a4\_v2\_source\_history\_i\_r1\_path\_closure\_erratum\_protocol\_2026\_07\_18.md}
\NormalTok{docs/saq\_a4\_v2\_source\_history\_i\_r1\_path\_closure\_erratum\_contract\_2026\_07\_18.json}
\NormalTok{docs/saq\_a4\_v2\_source\_history\_i\_r1\_path\_closure\_erratum\_independent\_review\_2026\_07\_18.md}
\NormalTok{script/a4\_v2\_cache\_policy\_verifier.py}
\NormalTok{docs/saq\_a4\_v2\_source\_history\_epoch\_correction\_implementation\_independent\_review\_2026\_07\_17.md}
\NormalTok{docs/saq\_a4\_v2\_prep\_start\_only\_crash\_record\_protocol\_2026\_07\_18.md}
\NormalTok{docs/saq\_a4\_v2\_prep\_start\_only\_crash\_record\_contract\_2026\_07\_18.json}
\end{Highlighting}
\end{Shaded}

\begin{Shaded}
\begin{Highlighting}[]
\NormalTok{docs/saq\_a4\_v2\_prep\_start\_only\_crash\_record\_protocol\_independent\_review\_2026\_07\_18.md}
\end{Highlighting}
\end{Shaded}

A4-0 and A4-1S parity are outcome evidence only at their stated
instrument ceilings. A4-1S cost is a reviewed pipeline-cost no-go. A4-1
remains an unexecuted frozen protocol; no natural-data, ANN, or systems
result is established. A4 V2 has an

independently reviewed corrected protocol and source plus one reviewed
pre-build \texttt{ARTIFACT\_INVALID} PAR invocation. It produced no
valid PAR authority or scientific decision. I-R1 later reached static
source- repair review pass only; it did not run. CACHE-P later reached
\texttt{EXACT\_TARGET\_INDEPENDENT\_REVIEW\_PASS},

requiring sanitized remote isolation and rejecting direct PAR-R1 at an
artifact-governance ceiling. It did not read or clean the quarantine.
CACHE-I then reached
\texttt{GENERIC\_CACHE\_POLICY\_SOURCE\_STATIC\_REVIEW\_PASS} for its
generic source/schema layer, with a separate syntax-only pass but no
import or

code execution. The later PREP authorization record reached
\texttt{PREP\_AUTHORIZATION\_EXACT\_TARGET\_REVIEW\_PASS}. A later
static host mismatch prevented PREP and START; the old probe is
superseded and may never be reused. The host-rebind erratum reached only
\texttt{HOST\_IDENTITY\_REBIND\_PROTOCOL\_REVIEW\_PASS}, not host

rebound or PREP readiness. HOST-I-AUTH then reached only
\texttt{AUTHORIZATION\_EXACT\_TARGET\_REVIEW\_PASS}; it changed no
source/authority object and ran no Python or PREP. The attempted HOST-I
edge stopped before a target on the seven-versus-eight contradiction;
its additive source-authority erratum reached

only \texttt{HOST\_I\_SOURCE\_AUTHORITY\_ERRATUM\_REVIEW\_PASS} and
changed no implementation source or derived authority object. The later
HOST-I source-static target/review \texttt{4602585/@e8e9c79} failed on
one HIGH false governing-review digest, so host rebound remains
unestablished and correction was not performed. Correction-only protocol
target/review \texttt{ddfef99/@b1a7429} then

passed with zero LOW-or-higher findings. Exact R1 target/review
\texttt{e17f887/@e10bde7} subsequently passed and established rebound
only for the registered 13-object bundle; whole-host identity remains
unestablished and no execution ran. A separate
\texttt{SOURCE\_HISTORY\_MISMATCH} source repair was still required
before cache-verifier

or PAR readiness; its protocol target/review \texttt{dcaed57/@e98a3e4}
passed, the first implementation authorization stopped before target on
an unsatisfiable ten-path projection, and path-closure erratum
target/review \texttt{150e5b3/@7749491} passed only for a corrected
future eight-path repair. Exact R1 target/review
\texttt{e9b5c83/@a1198c4} then

corrected only that static rule and reached
\texttt{SOURCE\_HISTORY\_EPOCH\_CORRECTION\_SOURCE\_STATIC\_REVIEW\_PASS}
with zero LOW+. The verifier was not run, so readiness remains
unestablished. Fresh dormant PREP authorization target/review
\texttt{cd1757e/@eb1b893} then passed exact documentation review with
zero LOW+. Its later unique immediate

probe passed, but the consumed invocation left only durable START plus
empty captures before the PREP body. Frozen status is
\texttt{CRASH\_OR\_UNKNOWN\ /\ NO\_SCIENTIFIC\_DECISION}. Crash-record
protocol target/review \texttt{1982865/@8bec546} then passed with zero
LOW+, freezing only the future record topology. CRASH-I,

repair/retry, CACHE-BIND, PAR-R1, A4-V2-SRUN, and all data access remain
unauthorized until their own reviewed and explicit authorities exist.

Current deck:

\begin{Shaded}
\begin{Highlighting}[]
\NormalTok{docs/saq\_next\_meeting\_attempts\_1\_4\_slides\_2026\_07\_13.md}
\end{Highlighting}
\end{Shaded}
\end{frame}

\begin{frame}[fragile,allowframebreaks]{69. One-Slide Takeaway}
\protect\phantomsection\label{69-one-slide-takeaway}
\begin{Shaded}
\begin{Highlighting}[]
\NormalTok{Attempt 1:}
\NormalTok{PCA replacement found a small GIST estimator signal that failed CIFAR}
\NormalTok{replication. The favorable lossy D{-}\textgreater{}d oracle still ranked worse than native}
\NormalTok{SAQ because a norm{-}only tail cannot recover tail inner products.}

\NormalTok{Attempt 2:}
\NormalTok{The paper and pinned public code do not present an integrated shared{-}codebook,}
\end{Highlighting}
\end{Shaded}

\begin{Shaded}
\begin{Highlighting}[]
\NormalTok{heterogeneous{-}rate, outer{-}allocation pipeline; that absence does not establish}
\NormalTok{novelty. Exact 1D scalar DP is old prior art. Separately, the pinned code}
\NormalTok{exposes a full{-}raw{-}data exact solver primitive and an adjacent shared{-}block}
\NormalTok{helper, neither with a public integrated experiment path. Local raw{-}SSE gains}
\NormalTok{are real, but Recall can improve or regress. Keep this as offline baseline and}
\NormalTok{objective{-}mismatch evidence, not a method.}
\end{Highlighting}
\end{Shaded}

\begin{Shaded}
\begin{Highlighting}[]
\NormalTok{Attempt 3:}
\NormalTok{Paper{-}exact 1/Ratio changes one GIST matched{-}quality conclusion: a measured}
\NormalTok{fac{-}error point gives higher geometric quality and 1.078x QPS. DEEP B=4/B=5}
\NormalTok{remain negative, 39.6\% of paired GIST queries worsen, and the metric is prior}
\NormalTok{work. The result is measurement evidence, not a method.}

\NormalTok{Attempt 4:}
\end{Highlighting}
\end{Shaded}

\begin{Shaded}
\begin{Highlighting}[]
\NormalTok{Fixed{-}rate arbitrary cardinalities strictly beat the dyadic scalar{-}product}
\NormalTok{restriction on the frozen (3,5) witness, but unrestricted block VQ matches}
\NormalTok{them and the primitives are known. A4{-}0 validates only the instrument. A4{-}1S}
\NormalTok{exact implementation parity passed, but the frozen construction/evidence}
\NormalTok{pipeline projected to 24.170246892361 CPU{-}hours and failed its 24{-}hour gate.}
\NormalTok{That terminal result remains unchanged. A4 V2 now freezes a direct internal}
\NormalTok{comparative{-}instrument FOM and a complete evidence ledger. Its reviewed}
\end{Highlighting}
\end{Shaded}

\begin{Shaded}
\begin{Highlighting}[]
\NormalTok{erratum adds only a finite one{-}file terminal{-}P reporting closure; the old 5/2}
\NormalTok{projection does not transfer. Corrected V2 source passed independent static}
\NormalTok{review. The later authorized PAR conductor ran only far enough to reject the}
\NormalTok{\textasciigrave{}/bin/python\textasciigrave{} symlink at its no{-}follow identity boundary; no build, RNG,}
\NormalTok{fixture, parity case, or final artifact was produced. I{-}R1 subsequently bound}
\NormalTok{leader identity to stable \textasciigrave{}/proc/self/exe\textasciigrave{} bytes plus same{-}inode}
\NormalTok{\textasciigrave{}sys.executable\textasciigrave{} validation and passed exact{-}commit static review. It was not}
\end{Highlighting}
\end{Shaded}

\begin{Shaded}
\begin{Highlighting}[]
\NormalTok{executed and therefore establishes neither parity nor artifact readiness.}
\NormalTok{CACHE{-}P subsequently reviewed CPython cache semantics and the frozen source}
\NormalTok{without reading the quarantine. Exact target \textasciigrave{}56210f8\textasciigrave{}, independently reviewed}
\NormalTok{at \textasciigrave{}249d5b8\textasciigrave{}, freezes sanitized remote isolation and rejects direct PAR{-}R1.}
\NormalTok{That is artifact governance only: it creates no clone, preparation, parity, or}
\NormalTok{scientific evidence. CACHE{-}I target \textasciigrave{}c33a2bf\textasciigrave{}, independently reviewed at}
\NormalTok{\textasciigrave{}5db3025\textasciigrave{}, then closed 37 filesystem sources and one separate inline bootstrap}
\end{Highlighting}
\end{Shaded}

\begin{Shaded}
\begin{Highlighting}[]
\NormalTok{under the generic cache policy. Its separately authorized syntax{-}only check}
\NormalTok{executed no code object. The result is only}
\NormalTok{\textasciigrave{}GENERIC\_CACHE\_POLICY\_SOURCE\_STATIC\_REVIEW\_PASS\textasciigrave{}, not PREP readiness or}
\NormalTok{feasibility evidence. PREP authorization target \textasciigrave{}e7f940e\textasciigrave{}, independently}
\NormalTok{reviewed at \textasciigrave{}16a8201\textasciigrave{}, subsequently reached only}
\NormalTok{\textasciigrave{}PREP\_AUTHORIZATION\_EXACT\_TARGET\_REVIEW\_PASS\textasciigrave{}. Static prelaunch checking then}
\NormalTok{found the \textasciigrave{}.el9\_8 {-}\textgreater{} .el9\_8.2\textasciigrave{} leader replacement before PREP or START. The}
\end{Highlighting}
\end{Shaded}

\begin{Shaded}
\begin{Highlighting}[]
\NormalTok{host{-}rebind erratum target \textasciigrave{}5a47fed\textasciigrave{}, independently reviewed at \textasciigrave{}b89dabe\textasciigrave{},}
\NormalTok{reached only \textasciigrave{}HOST\_IDENTITY\_REBIND\_PROTOCOL\_REVIEW\_PASS\textasciigrave{}. It created no clone,}
\NormalTok{token, or receipt, did not rebind the host, and superseded the old probe;}
\NormalTok{HOST{-}I{-}AUTH target \textasciigrave{}212a67b\textasciigrave{}, independently reviewed at \textasciigrave{}71e6bec\textasciigrave{}, then}
\NormalTok{reached only \textasciigrave{}AUTHORIZATION\_EXACT\_TARGET\_REVIEW\_PASS\textasciigrave{}. That is documentation}
\NormalTok{governance only: no five{-}source/seven{-}derived{-}object rebind, Python, build, or}
\NormalTok{PREP occurred. The attempted HOST{-}I source edge stopped before a target on an}
\end{Highlighting}
\end{Shaded}

\begin{Shaded}
\begin{Highlighting}[]
\NormalTok{exact seven{-}versus{-}eight component contradiction. Source{-}authority erratum}
\NormalTok{target \textasciigrave{}6fe8544\textasciigrave{}, independently reviewed at \textasciigrave{}9fa9528\textasciigrave{}, reached only}
\NormalTok{\textasciigrave{}HOST\_I\_SOURCE\_AUTHORITY\_ERRATUM\_REVIEW\_PASS\textasciigrave{}; it changes no implementation}
\NormalTok{source or derived authority object. The later source{-}static target \textasciigrave{}4602585\textasciigrave{},}
\NormalTok{independently reviewed at \textasciigrave{}e8e9c79\textasciigrave{}, failed with one HIGH because two}
\NormalTok{authority documents cite the wrong governing{-}review SHA{-}256. Host rebound is}
\NormalTok{not established. Correction{-}only protocol target \textasciigrave{}ddfef99\textasciigrave{}, independently}
\end{Highlighting}
\end{Shaded}

\begin{Shaded}
\begin{Highlighting}[]
\NormalTok{reviewed at \textasciigrave{}b1a7429\textasciigrave{}, froze the exact five{-}path R1 edge. Correction target}
\NormalTok{\textasciigrave{}e17f887\textasciigrave{}, independently reviewed at \textasciigrave{}e10bde7\textasciigrave{}, then passed with zero LOW{-}or{-}}
\NormalTok{higher findings across the correction, closure, and 13/13 registered host}
\NormalTok{objects. This establishes only registered{-}bundle rebound, not whole{-}host or}
\NormalTok{execution readiness. Source{-}history protocol target \textasciigrave{}dcaed57\textasciigrave{}, independently}
\NormalTok{reviewed at \textasciigrave{}e98a3e4\textasciigrave{}, subsequently passed and froze a strict two{-}epoch future}
\NormalTok{repair. The first implementation authorization stopped before target on an}
\end{Highlighting}
\end{Shaded}

\begin{Shaded}
\begin{Highlighting}[]
\NormalTok{unsatisfiable ten{-}path closure. Path{-}closure erratum target \textasciigrave{}150e5b3\textasciigrave{},}
\NormalTok{independently reviewed at \textasciigrave{}7749491\textasciigrave{}, then passed only for a corrected future}
\NormalTok{eight{-}path/five{-}derived{-}object repair. Exact R1 target \textasciigrave{}e9b5c83\textasciigrave{}, independently}
\NormalTok{reviewed at \textasciigrave{}a1198c4\textasciigrave{}, then implemented only that correction and reached}
\NormalTok{\textasciigrave{}SOURCE\_HISTORY\_EPOCH\_CORRECTION\_SOURCE\_STATIC\_REVIEW\_PASS\textasciigrave{}. It ran no}
\NormalTok{verifier, build, PREP, or data, so cache{-}verifier/PAR readiness is still}
\NormalTok{unestablished. Fresh dormant PREP authorization target \textasciigrave{}cd1757e\textasciigrave{}, independently}
\end{Highlighting}
\end{Shaded}

\begin{Shaded}
\begin{Highlighting}[]
\NormalTok{reviewed at \textasciigrave{}eb1b893\textasciigrave{}, then reached}
\NormalTok{\textasciigrave{}PREP\_AUTHORIZATION\_EXACT\_TARGET\_REVIEW\_PASS\textasciigrave{} with zero LOW{-}or{-}higher}
\NormalTok{findings. Its later unique probe/invocation was consumed and stopped at}
\NormalTok{START{-}only \textasciigrave{}CRASH\_OR\_UNKNOWN\textasciigrave{} before the PREP body. Crash{-}record protocol}
\NormalTok{target/review \textasciigrave{}1982865/@8bec546\textasciigrave{} then passed, defining only a future terminal{-}}
\NormalTok{record topology and no scientific or retry/repair authority.}
\end{Highlighting}
\end{Shaded}

\begin{Shaded}
\begin{Highlighting}[]
\NormalTok{Project decision:}
\NormalTok{Keep Attempts 1{-}{-}3 as rigorous negative or partial evidence. Do not rescue}
\NormalTok{them with post{-}hoc sweeps. Preserve A4{-}1S as a preregistered synthetic cost}
\NormalTok{stop. Treat A4 V2 only as reviewed protocol plus a source{-}static review pass,}
\NormalTok{followed by a reviewed terminal \textasciigrave{}ARTIFACT\_INVALID\textasciigrave{} and a later reviewed static}
\NormalTok{source repair, not feasibility evidence. CACHE{-}P additionally provides a}
\NormalTok{reviewed isolation protocol, with}
\end{Highlighting}
\end{Shaded}

\begin{Shaded}
\begin{Highlighting}[]
\NormalTok{\textasciigrave{}GO\_SANITIZED\_REMOTE\_ISOLATION\_PROTOCOL\_REQUIRED / NO\_GO\_DIRECT\_PAR\_R1\textasciigrave{}, but}
\NormalTok{the later CACHE{-}I result remains static artifact governance only. Treat the}
\NormalTok{subsequent PREP authorization and host{-}rebind erratum reviews as documentation}
\NormalTok{governance, not PREP readiness or execution. Treat HOST{-}I{-}AUTH the same way:}
\NormalTok{it freezes only the contract/projection for a possible future source edge and}
\NormalTok{grants no HOST{-}I authority. Treat the source{-}authority erratum as a correction}
\NormalTok{of that projection only, not source implementation, host rebound, or}
\end{Highlighting}
\end{Shaded}

\begin{Shaded}
\begin{Highlighting}[]
\NormalTok{readiness. Treat the later HOST{-}I target as a failed authority{-}DAG target, not}
\NormalTok{host rebound. Its later correction{-}only protocol passed direct{-}child review,}
\NormalTok{and actual R1 subsequently passed direct{-}child review, establishing only the}
\NormalTok{registered 13{-}object bundle identity. Do not promote that result to whole{-}host,}
\NormalTok{PREP, cache/PAR, performance, or scientific readiness. Treat the reviewed}
\NormalTok{source{-}history protocol and path{-}closure erratum as correction governance;}
\NormalTok{the first implementation attempt stopped before target, while the later exact}
\end{Highlighting}
\end{Shaded}

\begin{Shaded}
\begin{Highlighting}[]
\NormalTok{R1 correction passed direct{-}child static review only. Do not promote that}
\NormalTok{source result to executed verifier, PREP, cache/PAR, performance, or}
\NormalTok{scientific readiness. Treat fresh R1 PREP authorization target/review}
\NormalTok{\textasciigrave{}cd1757e/@eb1b893\textasciigrave{} as the reviewed authority for the now{-}consumed unique}
\NormalTok{probe/invocation. That invocation produced only START{-}only}
\NormalTok{\textasciigrave{}CRASH\_OR\_UNKNOWN\textasciigrave{}, not PREP readiness or science. Treat reviewed protocol}
\NormalTok{\textasciigrave{}1982865/@8bec546\textasciigrave{} as record governance only. CRASH{-}I, source repair/retry,}
\end{Highlighting}
\end{Shaded}

\begin{Shaded}
\begin{Highlighting}[]
\NormalTok{CACHE{-}BIND, separately explicit PAR{-}R1, and SRUN remain a strictly ordered}
\NormalTok{unauthorized chain; base, query, and SAQ gates remain closed.}
\end{Highlighting}
\end{Shaded}
\end{frame}

\end{document}
